% Environmental Chemistry: Sulphur and Nitrogen — Chemistry Upper Sixth — Cours signé OnBuch+
% Official MINESEC syllabus. Compiler avec : tectonic CHE05-environmental-chemistry-sulphur-and-nitrogen.tex
\newcommand{\DOCMATIERE}{Chemistry}
\newcommand{\DOCNIVEAU}{Upper Sixth}
\newcommand{\DOCMODULE}{Upper Sixth — Chemistry}
\newcommand{\DOCLECON}{5}
\newcommand{\DOCTITRE}{Environmental Chemistry: Sulphur and Nitrogen}
% OnBuch+ — préambule « cours signés » ENRICHI (compilé avec tectonic / XeLaTeX)
% Système visuel « Pop » d'OnBuch+ : fond crème (jamais blanc), bordures encre
% épaisses, ombres « dures » décalées, titres Archivo Black en capitales.
% Version MATHÉMATIQUES : graphes (pgfplots/tikz), boîtes pédagogiques
% (définition, propriété, méthode, attention, le savais-tu, exercice résolu,
% exercice, TP, bilan), page de garde et signature OnBuch+ ancrée en bas.
%
% Macros attendues dans main.tex AVANT \input{preamble} :
%   \DOCMATIERE  \DOCNIVEAU  \DOCMODULE  \DOCLECON (numéro)  \DOCTITRE
\documentclass[11pt]{article}

\usepackage{fontspec}
\usepackage[english]{babel}
\usepackage[a4paper,margin=2cm,top=3.4cm,bottom=2.8cm,headheight=1.4cm,footskip=1.3cm]{geometry}
\usepackage{amsmath,amssymb,amsfonts}
\usepackage{mathtools} % \xmapsto, \coloneqq... souvent utilisés par les modèles
\usepackage{cancel} % simplifications barrées (bilans de forces)
% \abs{z} (module, valeur absolue) et \norm{u} (norme), utilisés par les modèles
\providecommand{\abs}{}\let\abs\relax\DeclarePairedDelimiter{\abs}{\lvert}{\rvert}
\providecommand{\norm}{}\let\norm\relax\DeclarePairedDelimiter{\norm}{\lVert}{\rVert}
\usepackage[table]{xcolor} % [table] : fournit aussi \rowcolors (alternance de lignes)
\usepackage{pagecolor}
\usepackage{tikz}
\usetikzlibrary{arrows.meta,positioning,calc,shapes.geometric,shapes.misc,shapes.arrows,decorations.pathmorphing,decorations.pathreplacing,decorations.markings,patterns,shadows,mindmap,trees,fit,backgrounds,matrix,intersections,angles,quotes}
\usepackage{pgfplots}
\usepackage[version=4]{mhchem} % équations nucléaires \ce{^{235}_{92}U}
\usepackage[european,siunitx]{circuitikz} % schémas électriques
\pgfplotsset{compat=1.18}
\usepgfplotslibrary{fillbetween,groupplots} % aires entre courbes (calcul intégral)
\usepackage{enumitem}
\usepackage{fancyhdr}
% [most] charge toutes les bibliothèques tcolorbox (listings, minted, raster,
% documentation...) et gonfle inutilement la mémoire TeX sur un document long
% (~300 boîtes) : on ne charge que ce qui est réellement utilisé.
\usepackage[skins,breakable]{tcolorbox}
\usepackage{graphicx}
\usepackage{colortbl}
\usepackage{booktabs}
\usepackage{tabularx}
\usepackage{diagbox} % case d'en-tête diagonale des tableaux à double entrée
% Colonnes extensibles centrée / gauche / droite (C, L, R), souvent utilisées par les modèles
\newcolumntype{C}{>{\centering\arraybackslash}X}
\newcolumntype{L}{>{\raggedright\arraybackslash}X}
\newcolumntype{R}{>{\raggedleft\arraybackslash}X}
\usepackage{array}
\usepackage{multicol}
\usepackage{siunitx}
% parse-numbers=false : accepte aussi \SI{2,5 \times 10^{-3}}{...}
\sisetup{locale=UK,output-decimal-marker={.},group-minimum-digits=4,parse-numbers=false}
\DeclareSIUnit\ohm{\ensuremath{\symup{\Omega}}} % U+2126 absent de la police
\DeclareSIUnit\division{div} % réglages d'oscilloscope (ms/div, V/div)
\DeclareSIUnit\tr{tr} % tours (vitesse de rotation, ex. \SI{1500}{\tr\per\minute})
\usepackage{qrcode}
% \degree : commande fréquemment utilisée par les modèles pour un angle ou une
% température en dehors de \SI{}{} (non fournie par les paquets chargés).
\providecommand{\degree}{^{\circ}}
% \qed (fin de démonstration), utilisé par les modèles sans amsthm
\providecommand{\qed}{\hfill$\square$}
% \barre : double barre des valeurs interdites dans un tableau de variations (tkz-tab)
\providecommand{\barre}{\ensuremath{\|}}
% environnement proof (amsthm non chargé) utilisé par les modèles
\ifdefined\proof\else\newenvironment{proof}[1][Proof]{\par\noindent\textit{#1.}\ }{\qed\par}\fi
\usepackage{lastpage}
\usepackage{hyperref}
\hypersetup{hidelinks}

% ============================================================
% Palette « Pop » OnBuch+ — mêmes hex que la classe P (plus_ui.dart)
% ============================================================
\definecolor{poporange}{HTML}{F59321}
\definecolor{popdark}{HTML}{E07A0C}
\definecolor{popcream}{HTML}{FFF6E8}
\definecolor{popink}{HTML}{1C1714}
\definecolor{poppurple}{HTML}{7A5AE0}
\definecolor{popgreen}{HTML}{2E9E6A}
\definecolor{poppink}{HTML}{E0457B}
\definecolor{popblue}{HTML}{2F7DE1}
\definecolor{popgold}{HTML}{C98A12}
\definecolor{popmuted}{HTML}{8A7F76}
\definecolor{popline}{HTML}{EADFD2}
\definecolor{popgray}{HTML}{8A7F76}
\colorlet{popgrey}{popgray}
% Alias tolérants (le modèle nomme parfois les couleurs différemment)
\colorlet{popwhite}{white}
\colorlet{popblack}{popink}
\colorlet{popred}{poppink}
\colorlet{popyellow}{popgold}
% Teintes claires pour fonds de tableaux / zones
\colorlet{popblueL}{popblue!10!white}
\colorlet{poporangeL}{poporange!14!white}
\colorlet{popgreenL}{popgreen!12!white}
\colorlet{poppurpleL}{poppurple!12!white}
\colorlet{poppinkL}{poppink!12!white}
\colorlet{popgoldL}{popgold!14!white}

\pagecolor{popcream}

% ============================================================
% Typographie
% ============================================================
\defaultfontfeatures{Ligatures=TeX}
\setmainfont{PlusJakartaSans-Medium.ttf}[Path=../../../fonts/,BoldFont=PlusJakartaSans-Bold.ttf]
% Maths en sans-serif (Fira Math) avec chiffres et lettres droites de Plus
% Jakarta, pour que les formules se fondent dans le texte.
\usepackage[math-style=ISO,bold-style=ISO]{unicode-math}
\setmathfont{FiraMath-Regular.otf}[Path=../../../fonts/]
\setmathfont{PlusJakartaSans-Medium.ttf}[Path=../../../fonts/,range={up/num,up/{Latin,latin},\mathup}]
% (icomma retiré : décimales avec point en anglais)
% Caractères mathématiques Unicode absents des polices de texte (Plus Jakarta,
% Archivo Black) : rendus par leur équivalent mathématique, en texte comme en
% formule — sinon ils s'affichent en carré vide (ex. « Arithmétique dans ℤ »).
\usepackage{newunicodechar}
\newunicodechar{ℝ}{\ensuremath{\mathbb{R}}}
\newunicodechar{ℤ}{\ensuremath{\mathbb{Z}}}
\newunicodechar{ℂ}{\ensuremath{\mathbb{C}}}
\newunicodechar{ℕ}{\ensuremath{\mathbb{N}}}
\newunicodechar{ℚ}{\ensuremath{\mathbb{Q}}}
\newunicodechar{𝔻}{\ensuremath{\mathbb{D}}}
\newunicodechar{α}{\ensuremath{\alpha}}
\newunicodechar{β}{\ensuremath{\beta}}
\newunicodechar{γ}{\ensuremath{\gamma}}
\newunicodechar{δ}{\ensuremath{\delta}}
\newunicodechar{ε}{\ensuremath{\varepsilon}}
\newunicodechar{θ}{\ensuremath{\theta}}
\newunicodechar{λ}{\ensuremath{\lambda}}
\newunicodechar{μ}{\ensuremath{\mu}}
\newunicodechar{σ}{\ensuremath{\sigma}}
\newunicodechar{τ}{\ensuremath{\tau}}
\newunicodechar{φ}{\ensuremath{\varphi}}
\newunicodechar{ψ}{\ensuremath{\psi}}
\newunicodechar{ω}{\ensuremath{\omega}}
\newunicodechar{Σ}{\ensuremath{\Sigma}}
% majuscules produites par \MakeUppercase dans les titres (α-aminés -> Α)
\newunicodechar{Α}{\ensuremath{\alpha}}
\newunicodechar{Β}{\ensuremath{\beta}}
\newunicodechar{Γ}{\ensuremath{\Gamma}}
\newunicodechar{Δ}{\ensuremath{\Delta}}
\newunicodechar{Ε}{\ensuremath{\varepsilon}}
\newunicodechar{Θ}{\ensuremath{\Theta}}
\newunicodechar{Λ}{\ensuremath{\Lambda}}
\newunicodechar{Μ}{\ensuremath{\mu}}
\newunicodechar{Τ}{\ensuremath{\tau}}
\newunicodechar{Φ}{\ensuremath{\Phi}}
\newunicodechar{Ψ}{\ensuremath{\Psi}}
\newunicodechar{Ω}{\ensuremath{\Omega}}
\newunicodechar{↦}{\ensuremath{\mapsto}}
\newunicodechar{∈}{\ensuremath{\in}}
\newunicodechar{∉}{\ensuremath{\notin}}
\newunicodechar{∧}{\ensuremath{\wedge}}
\newunicodechar{⇔}{\ensuremath{\Leftrightarrow}}
\newunicodechar{⟺}{\ensuremath{\Longleftrightarrow}}
\newunicodechar{⇒}{\ensuremath{\Rightarrow}}
\newunicodechar{⟹}{\ensuremath{\Longrightarrow}}
\newunicodechar{∘}{\ensuremath{\circ}}
\newunicodechar{∪}{\ensuremath{\cup}}
\newunicodechar{∩}{\ensuremath{\cap}}
\newunicodechar{≡}{\ensuremath{\equiv}}
\newunicodechar{⊂}{\ensuremath{\subset}}
\newunicodechar{⊥}{\ensuremath{\perp}}
\newunicodechar{∃}{\ensuremath{\exists}}
\newunicodechar{∀}{\ensuremath{\forall}}
\newunicodechar{∛}{\ensuremath{\sqrt[3]{\,}}}
\newunicodechar{∜}{\ensuremath{\sqrt[4]{\,}}}
\newunicodechar{⃗}{\ensuremath{{}^{\to}}}
\newunicodechar{ⁿ}{\ifmmode^{n}\else\textsuperscript{\ensuremath{n}}\fi}
\newunicodechar{ˣ}{\ifmmode^{x}\else\textsuperscript{\ensuremath{x}}\fi}
\newunicodechar{ᵇ}{\ifmmode^{b}\else\textsuperscript{\ensuremath{b}}\fi}
\newunicodechar{ʳ}{\ifmmode^{r}\else\textsuperscript{\ensuremath{r}}\fi}
\newunicodechar{ᵘ}{\ifmmode^{u}\else\textsuperscript{\ensuremath{u}}\fi}
\newunicodechar{ʸ}{\ifmmode^{y}\else\textsuperscript{\ensuremath{y}}\fi}
\newunicodechar{ᵃ}{\ifmmode^{a}\else\textsuperscript{\ensuremath{a}}\fi}
\newunicodechar{ᵉ}{\ifmmode^{e}\else\textsuperscript{\ensuremath{e}}\fi}
\newunicodechar{ᵏ}{\ifmmode^{k}\else\textsuperscript{\ensuremath{k}}\fi}
\newunicodechar{ᵖ}{\ifmmode^{p}\else\textsuperscript{\ensuremath{p}}\fi}
\newunicodechar{ᴺ}{\ifmmode^{N}\else\textsuperscript{\ensuremath{N}}\fi}
\newunicodechar{ᶜ}{\ifmmode^{c}\else\textsuperscript{\ensuremath{c}}\fi}
\newunicodechar{ʰ}{\ifmmode^{h}\else\textsuperscript{\ensuremath{h}}\fi}
\newunicodechar{ᵠ}{\ifmmode^{q}\else\textsuperscript{\ensuremath{q}}\fi}
\newunicodechar{ₙ}{\ifmmode_{n}\else\textsubscript{\ensuremath{n}}\fi}
\newunicodechar{ₐ}{\ifmmode_{a}\else\textsubscript{\ensuremath{a}}\fi}
\newunicodechar{ᵢ}{\ifmmode_{i}\else\textsubscript{\ensuremath{i}}\fi}
\newunicodechar{ᵣ}{\ifmmode_{r}\else\textsubscript{\ensuremath{r}}\fi}
\newunicodechar{ₖ}{\ifmmode_{k}\else\textsubscript{\ensuremath{k}}\fi}
\newunicodechar{ᵦ}{\ifmmode_{\beta}\else\textsubscript{\ensuremath{\beta}}\fi}
\newunicodechar{ₓ}{\ifmmode_{x}\else\textsubscript{\ensuremath{x}}\fi}
\newfontfamily\popdisplay{ArchivoBlack-Regular.ttf}[Path=../../../fonts/]
\newfontfamily\poplogo{SpaceGrotesk-Bold.ttf}[Path=../../../fonts/]
\newfontfamily\popsemi{PlusJakartaSans-SemiBold.ttf}[Path=../../../fonts/]
\newfontfamily\popxbold{PlusJakartaSans-ExtraBold.ttf}[Path=../../../fonts/]
\newfontfamily\popxboldls{PlusJakartaSans-ExtraBold.ttf}[Path=../../../fonts/,LetterSpace=3.0]

\setlist[enumerate]{leftmargin=1.7em,itemsep=5pt,topsep=4pt}
\setlist[itemize]{leftmargin=1.7em,itemsep=4pt,topsep=4pt,label={\color{poporange}\textbullet}}
\setlength{\parindent}{0pt}
\setlength{\parskip}{5pt}
\renewcommand{\arraystretch}{1.25}

% pgfplots : style Pop par défaut
\pgfplotsset{
  every axis/.append style={axis line style={popink,line width=1pt},tick style={popink},
    grid style={popline,line width=0.5pt},label style={font=\small},tick label style={font=\footnotesize}},
  popcurve/.style={poporange,line width=1.8pt,no markers,smooth},
}
% Même style disponible dans un \draw TikZ simple (hors pgfplots)
\tikzset{popcurve/.style={poporange,line width=1.8pt}}
\tikzset{popgrid/.style={popline,very thin}}
\colorlet{popcurve}{poporange} % le nom du style est parfois utilisé comme couleur

% ============================================================
% Pill / squiggle
% ============================================================
\newcommand{\poppill}[3]{%
  \tikz[baseline=-0.55ex]\node[draw=popink,line width=1.2pt,rounded corners=2.6pt,%
    fill=#2,inner xsep=6.5pt,inner ysep=3pt]{%
    \popxboldls\fontsize{7}{7}\selectfont\color{#3}\MakeUppercase{#1}};%
}
\newcommand{\popsquiggle}[1]{%
  \tikz[baseline=-0.4ex]{
    \draw[#1,line width=2.6pt,line cap=round]
      plot [smooth] coordinates {(0,0.09) (0.34,0.01) (0.68,0.21) (1.02,0.01) (1.36,0.10)};
  }%
}
% Mot-clé surligné (marqueur orange sous le texte)
\newcommand{\cle}[1]{{\popxbold\color{popdark}#1}}

% ============================================================
% En-tête / pied
% ============================================================
\pagestyle{fancy}
\fancyhf{}
\renewcommand{\headrulewidth}{0pt}
\renewcommand{\footrulewidth}{0pt}
\lhead{\raisebox{-0.15ex}{\poplogo\fontsize{13}{13}\selectfont\color{popink}OnBuch\color{poporange}+}}
\rhead{\poppill{\DOCMATIERE\ \textbullet\ \DOCNIVEAU\ \textbullet\ Chapter \DOCLECON}{white}{popink}}
\lfoot{\popxbold\fontsize{7.6}{7.6}\selectfont\color{popmuted}\MakeUppercase{Signed course OnBuch+}\ \textbullet\ {\popsemi\DOCTITRE}}
\rfoot{\tikz[baseline=-0.6ex]\node[rounded corners=8.5pt,fill=popink,minimum height=17pt,minimum width=17pt,inner xsep=5pt,inner sep=0]{\popxbold\fontsize{8}{8}\selectfont\color{popcream}\thepage\,/\,\pageref*{LastPage}};}

% ============================================================
% Titres
% ============================================================
\newcommand{\chaptitre}[2]{%
  \begin{tcolorbox}[enhanced,colback=poporange,colframe=popink,boxrule=2.6pt,arc=11pt,
      drop shadow=popink,left=15pt,right=15pt,top=12pt,bottom=11pt,before skip=3pt,after skip=18pt]
    \poppill{#2}{popink}{popcream}\par\vspace{9pt}
    {\raggedright\hyphenpenalty=10000\exhyphenpenalty=10000\popdisplay\fontsize{22}{24}\selectfont\color{popcream}\MakeUppercase{#1}\par}
  \end{tcolorbox}
}
% Section numérotée : pastille orange avec le numéro + titre Archivo
\newcounter{popsec}
\newcommand{\coursec}[1]{%
  \stepcounter{popsec}\setcounter{popsub}{0}%
  \par\vspace{16pt}\noindent
  \begin{minipage}{\linewidth}
  \tikz[baseline=-0.7ex]\node[draw=popink,line width=1.6pt,fill=poporange,rounded corners=4pt,minimum size=22pt,inner sep=0]{\popdisplay\fontsize{12}{12}\selectfont\color{popcream}\thepopsec};%
  \hspace{8pt}{\popdisplay\fontsize{15}{17}\selectfont\color{popink}\MakeUppercase{#1}}\par
  \vspace{4pt}\noindent{\color{poporange}\rule{36pt}{3.6pt}}
  \end{minipage}\par\nopagebreak\vspace{8pt}
}
\newcounter{popsub}
\newcommand{\courssub}[1]{%
  \stepcounter{popsub}%
  \par\vspace{9pt}\noindent{\popxbold\fontsize{11.5}{13}\selectfont\color{popink}{\color{poporange}\thepopsec.\thepopsub}\hspace{6pt}#1}\par\nopagebreak\vspace{4pt}
}

% ============================================================
% Boîtes pédagogiques — même recette : fond blanc, bordure épaisse colorée,
% ombre dure encre, badge plein. Titre optionnel [#1].
% ============================================================
\tcbset{popbase/.style={breakable,enhanced,colback=white,boxrule=2.3pt,arc=9pt,
  shadow={3pt}{-3pt}{0pt}{popink},left=14pt,right=14pt,top=11pt,bottom=11pt,before skip=11pt,after skip=15pt,
  fontupper=\popsemi\fontsize{10.6}{14.5}\selectfont\color{popink},
  fontlower=\popsemi\fontsize{10.6}{14.5}\selectfont\color{popink}}}
% Coche dessinée (le glyphe ✓ n'existe pas dans les polices du document)
\newcommand{\popcheck}[1]{\tikz[baseline=-0.5ex]\draw[#1,line width=1.6pt,line cap=round,line join=round](-0.13,0.0)--(-0.04,-0.09)--(0.13,0.1);}
\providecommand{\coche}{\popcheck{popgreen}}
\providecommand{\resultat}[1]{\par\smallskip\noindent\fcolorbox{popgreen}{popgreenL}{\parbox{\dimexpr\linewidth-2\fboxsep-2\fboxrule\relax}{\textbf{Result.} #1}}\par\smallskip}
\newcommand{\popboxhead}[3]{\poppill{#1}{#2}{white}\ifx\relax#3\relax\else\hspace{6pt}{\popxbold\fontsize{10.6}{12}\selectfont\color{popink}#3}\fi\par\vspace{7pt}}

\newtcolorbox{aretenir}[1][]{popbase,colframe=popgreen,before upper={\popboxhead{Key points}{popgreen}{#1}}}
\newtcolorbox{exemplebox}[1][]{popbase,colframe=poporange,before upper={\popboxhead{Exemple}{poporange}{#1}}}
\newtcolorbox{definition}[1][]{popbase,colframe=popblue,before upper={\popboxhead{Definition}{popblue}{#1}}}
\newtcolorbox{propriete}[1][]{popbase,colframe=poppurple,before upper={\popboxhead{Property}{poppurple}{#1}}}
\newtcolorbox{methode}[1][]{popbase,colframe=popgold,before upper={\popboxhead{Method}{popgold}{#1}}}
\newtcolorbox{attention}[1][]{popbase,colframe=poppink,before upper={\popboxhead{Warning}{poppink}{#1}}}
\newtcolorbox{experience}[1][]{popbase,colframe=popblue,colback=popblueL,before upper={\popboxhead{Experiment}{popblue}{#1}}}
% « Le savais-tu ? » : carte violette pleine (contexte, culture, Cameroun)
\newtcolorbox{savaistu}[1][]{popbase,colback=poppurple,colframe=popink,
  fontupper=\popsemi\fontsize{10.4}{14.2}\selectfont\color{white},
  before upper={\poppill{Did you know?}{white}{poppurple}\ifx\relax#1\relax\else\hspace{6pt}{\popxbold\color{white}#1}\fi\par\vspace{7pt}}}
% Exercice résolu : énoncé en haut, \tcblower puis la solution sur fond crème
\newcounter{popexo}\newcounter{popexr}
\newtcolorbox{exoresolu}[1][]{popbase,colframe=poporange,colbacklower=popcream,
  segmentation style={popink,line width=1.2pt,dashed},
  before upper={\stepcounter{popexr}\popboxhead{Worked example \thepopexr}{poporange}{#1}},
  before lower={\poppill{Solution}{popink}{popcream}\par\vspace{6pt}}}
% Exercice (non résolu en place) — la correction est dans la partie Corrigés
\newtcolorbox{exercice}[1][]{popbase,colframe=popink,
  before upper={\stepcounter{popexo}\popboxhead{Exercise \thepopexo}{popink}{#1}}}
% Corrigé
\newtcolorbox{corrige}[1][]{popbase,colframe=popgreen,colback=popgreenL,
  before upper={\popboxhead{Solution}{popgreen}{#1}}}
% Objectifs / prérequis / bilan
\newtcolorbox{objectifs}{popbase,colframe=popink,colback=poporangeL,
  before upper={\popboxhead{Chapter objectives}{popink}{}}}
\newtcolorbox{prerequis}{popbase,colframe=popink,colback=popblueL,
  before upper={\popboxhead{Prerequisites}{popblue}{}}}
\newtcolorbox{bilan}{popbase,colframe=popink,colback=white,boxrule=2.8pt,
  before upper={\popboxhead{Fiche bilan}{poporange}{L'essentiel à connaître pour l'examen}}}
% Figure encadrée avec légende
\newtcolorbox{popfigure}[1][]{enhanced,breakable=false,colback=white,colframe=popline,boxrule=1.4pt,arc=8pt,
  left=10pt,right=10pt,top=10pt,bottom=8pt,before skip=10pt,after skip=12pt,halign=center,
  lower separated=false,halign lower=center,
  fontlower=\popsemi\fontsize{9}{11.5}\selectfont\color{popmuted},
  #1}
\newcommand{\legende}[1]{\par\vspace{6pt}{\popsemi\fontsize{9}{11.5}\selectfont\color{popmuted}#1}}

% ============================================================
% Signature OnBuch+
% ============================================================
\newcommand{\poplogoinline}[1]{{\poplogo\fontsize{#1}{#1}\selectfont\color{popink}OnBuch\color{poporange}+}}
\newcommand{\signatureonbuch}{%
  \begin{tcolorbox}[enhanced,colback=popink,colframe=popink,boxrule=2.2pt,arc=10pt,
      drop shadow=poporange,left=14pt,right=14pt,top=10pt,bottom=10pt,before skip=0pt,after skip=0pt]
    \tikz[baseline=-0.7ex]\node[circle,fill=popgreen,draw=popcream,line width=1.3pt,minimum size=22pt,inner sep=0]{\popcheck{white}};%
    \hspace{9pt}\begin{minipage}[c]{0.62\linewidth}
      {\popxbold\fontsize{12}{13}\selectfont\color{popcream}Signed\ \textbullet\ {\poplogo OnBuch\color{poporange}+}}\par\vspace{2pt}
      {\raggedright\popsemi\fontsize{8}{10.5}\selectfont\color{popline}\MakeUppercase{OnBuch+ teaching team}\\\MakeUppercase{Follows the official MINESEC syllabus}\par}
    \end{minipage}\hfill
    {\poplogo\fontsize{22}{22}\selectfont\color{popcream}OnBuch\color{poporange}+}
  \end{tcolorbox}%
}

% ============================================================
% Page de garde
% #1 titre · #2 matière · #3 niveau · #4 module · #5 n° leçon · #6 durée ·
% #7 description / objectifs (texte ou itemize)
% ============================================================
\newcommand{\pagegarde}[7]{%
  \thispagestyle{empty}
  \newgeometry{a4paper,margin=1.7cm,top=1.6cm,bottom=1.4cm}%
  \begin{tikzpicture}[remember picture,overlay]
    \fill[poporange!18] ([xshift=-3.2cm,yshift=-3.6cm]current page.north east) circle (4.6cm);
    \fill[poppurple!14] ([xshift=2.4cm,yshift=7.6cm]current page.south west) circle (3.8cm);
  \end{tikzpicture}%
  {\poplogo\fontsize{30}{30}\selectfont\color{popink}OnBuch\color{poporange}+}\hfill
  \raisebox{0.6ex}{\poppill{Signed course \textbullet\ Premium}{popink}{popcream}}\par
  \vspace{1pt}\popsquiggle{poppurple}\par
  \vspace{8pt}
  \begin{tcolorbox}[enhanced,colback=poporange,colframe=popink,boxrule=2.8pt,arc=13pt,
      drop shadow=popink,left=18pt,right=18pt,top=12pt,bottom=13pt,after skip=10pt]
    \poppill{#2 \textbullet\ #3}{popink}{popcream}\hspace{5pt}\poppill{#4}{white}{popink}\par\vspace{8pt}
    {\popdisplay\fontsize{14}{14}\selectfont\color{popink}CHAPTER #5}\par\vspace{4pt}
    {\raggedright\hyphenpenalty=10000\exhyphenpenalty=10000\popdisplay\fontsize{25}{27}\selectfont\color{popcream}\MakeUppercase{#1}\par}
    \vspace{8pt}
    {\popxbold\fontsize{9.5}{9.5}\selectfont\color{popink}\MakeUppercase{Suggested time: #6}}
  \end{tcolorbox}
  \begin{tcolorbox}[enhanced,colback=white,colframe=popink,boxrule=2.2pt,arc=10pt,
      drop shadow=popink,left=16pt,right=16pt,top=10pt,bottom=10pt,after skip=8pt]
    \poppill{In this chapter}{poppurple}{white}\par\vspace{5pt}
    \popsemi\fontsize{9.3}{12.6}\selectfont\color{popink}#7
  \end{tcolorbox}
  \noindent\begin{minipage}[t]{0.487\linewidth}
    \begin{tcolorbox}[enhanced,colback=popgreen,colframe=popink,boxrule=2.2pt,arc=10pt,
        drop shadow=popink,left=11pt,right=11pt,top=10pt,bottom=10pt,height=3.1cm,valign=center]
      \tcbox[colback=white,colframe=popink,boxrule=1.3pt,arc=5pt,boxsep=0pt,nobeforeafter,
        left=3pt,right=3pt,top=3pt,bottom=3pt,valign=center]{\qrcode[height=1.9cm]{https://play.google.com/store/apps/details?id=com.luvvix.onbuch&hl=en}}%
      \hspace{8pt}\begin{minipage}[c]{0.5\linewidth}
        {\popdisplay\fontsize{11}{12.5}\selectfont\color{white}\MakeUppercase{Download OnBuch}}\par\vspace{4pt}
        {\raggedright\popsemi\fontsize{8.4}{11}\selectfont\color{white}Free \textbullet\ Google Play\\AI tutor, past papers, results\par}
      \end{minipage}
    \end{tcolorbox}
  \end{minipage}\hfill
  \begin{minipage}[t]{0.487\linewidth}
    \begin{tcolorbox}[enhanced,colback=poppurple,colframe=popink,boxrule=2.2pt,arc=10pt,
        drop shadow=popink,left=11pt,right=11pt,top=10pt,bottom=10pt,height=3.1cm,valign=center]
      \tikz[baseline=-0.7ex]\node[circle,fill=white,draw=popink,line width=1.3pt,minimum size=1.6cm,inner sep=0]{\popdisplay\fontsize{24}{24}\selectfont\color{poppurple}+};%
      \hspace{8pt}\begin{minipage}[c]{0.6\linewidth}
        {\popdisplay\fontsize{11}{12.5}\selectfont\color{white}\MakeUppercase{Go OnBuch+}}\par\vspace{4pt}
        {\raggedright\popsemi\fontsize{8.4}{11}\selectfont\color{white}All signed courses, unlimited AI tutor, PDF sheets — OnBuch+ tab in the app\par}
      \end{minipage}
    \end{tcolorbox}
  \end{minipage}
  \vspace{8pt}
  \popcontenu
  \vfill
  \signatureonbuch
  \restoregeometry
  \newpage
}

% Rangée « Ce cours contient » (page de garde)
\newcommand{\poptile}[3]{% #1 couleur #2 gros texte #3 légende
  \begin{tcolorbox}[enhanced,colback=#1,colframe=popink,boxrule=2pt,arc=9pt,shadow={2.5pt}{-2.5pt}{0pt}{popink},
      left=6pt,right=6pt,top=9pt,bottom=9pt,halign=center,valign=center,height=1.75cm,width=\linewidth,nobeforeafter]
    {\popdisplay\fontsize{12.5}{13}\selectfont\color{white}\MakeUppercase{#2}}\par\vspace{3pt}
    {\centering\popsemi\fontsize{7.8}{9.5}\selectfont\color{white}#3\par}
  \end{tcolorbox}}
\newcommand{\popcontenu}{%
  \par\noindent{\popxboldls\fontsize{7.5}{7.5}\selectfont\color{popmuted}THIS SIGNED COURSE CONTAINS}\par\vspace{7pt}
  \noindent\begin{minipage}[t]{0.235\linewidth}\poptile{popblue}{Course}{complete, illustrated, step by step}\end{minipage}\hfill
  \begin{minipage}[t]{0.235\linewidth}\poptile{poporange}{Methods}{and worked examples}\end{minipage}\hfill
  \begin{minipage}[t]{0.235\linewidth}\poptile{poppink}{Exercises}{exam style + solutions}\end{minipage}\hfill
  \begin{minipage}[t]{0.235\linewidth}\poptile{popgreen}{Summary sheet}{the essentials to remember}\end{minipage}\par}

% Sommaire
\newtcolorbox{sommaire}{enhanced,colback=white,colframe=popink,boxrule=2.2pt,arc=10pt,
  drop shadow=popink,left=18pt,right=18pt,top=13pt,bottom=13pt,before skip=6pt,after skip=20pt,
  before upper={\poppill{Contents}{poppurple}{white}\par\vspace{9pt}\popsemi\fontsize{10.6}{14.5}\selectfont\color{popink}}}

% Signature de fin de cours — poussée tout en bas de la dernière page
\newcommand{\findecours}{%
  \par\vfill\nopagebreak
  \begin{center}\popsquiggle{poporange}\end{center}\vspace{4pt}
  \signatureonbuch
}
\AtBeginDocument{\def\llbracket{\mathopen{[\mkern-3.5mu[}}\def\rrbracket{\mathclose{]\mkern-3.5mu]}}} % intervalles d'entiers (glyphe ⟦ absent de la police)
% Glyphes absents de Fira Math : redéfinis à partir de glyphes disponibles
\AtBeginDocument{%
  \def\setminus{\mathbin{\backslash}}\let\smallsetminus\setminus
  \def\vdots{\mathord{\vbox{\baselineskip3.2pt\lineskiplimit0pt\kern4pt\hbox{.}\hbox{.}\hbox{.}}}}%
  \def\ddots{\mathinner{\mkern1mu\raise6.5pt\hbox{.}\mkern2mu\raise3.6pt\hbox{.}\mkern2mu\raise0.7pt\hbox{.}\mkern1mu}}%
  \def\triangle{\mathord{\scalebox{0.9}{$\symup{\Delta}$}}}%
  \def\nabla{\mathord{\raisebox{\depth}{\scalebox{1}[-1]{$\symup{\Delta}$}}}}%
  \def\checkmark{\popcheck{popdark}}%
  \def\star{\mathbin{\ast}}\def\bigstar{\mathord{\ast}}%
  \def\diamond{\mathbin{\scalebox{0.6}{$\Diamond$}}}%
  \def\bigcup{\mathop{\mathchoice{\raisebox{-0.3em}{\scalebox{1.7}{$\cup$}}}{\scalebox{1.25}{$\cup$}}{\cup}{\cup}}\displaylimits}%
  \def\bigcap{\mathop{\mathchoice{\raisebox{-0.3em}{\scalebox{1.7}{$\cap$}}}{\scalebox{1.25}{$\cap$}}{\cap}{\cap}}\displaylimits}%
  \def\lesseqqgtr{\mathrel{\lessgtr}}\def\bigoplus{\mathop{\oplus}\displaylimits}%
}
\providecommand{\ssi}{if and only if} % abréviation fréquente
\AtBeginDocument{\providecommand{\wideparen}{\overparen}} % arc de cercle

%% FIGFIX-BEGIN v5 — correctifs globaux des figures (docs/onbuch-generation/tools/figfix)
\usepackage{adjustbox}\usepackage{etoolbox}\tcbuselibrary{raster}
% toute figure TikZ/pgfplots plus large que la ligne est réduite à la largeur disponible
\BeforeBeginEnvironment{tikzpicture}{\begin{adjustbox}{max width=\linewidth}}
\AfterEndEnvironment{tikzpicture}{\end{adjustbox}}
% légendes pgfplots lisibles par-dessus les courbes
\pgfplotsset{every axis/.append style={legend style={fill=white,fill opacity=0.92,text opacity=1,draw=black!15,inner sep=3pt}}}
% étiquettes d'arêtes (syntaxe edge["..."]) avec fond blanc
\tikzset{every edge quotes/.append style={fill=white,inner sep=1.2pt}}
%% FIGFIX-END
\begin{document}

\pagegarde{Environmental Chemistry: Sulphur and Nitrogen}{Chemistry}{Upper Sixth}{Upper Sixth — Chemistry}{5}{6 periods}{This chapter explores the chemistry of sulphur and nitrogen, two elements with wide ranges of oxidation states that underpin the manufacture of sulphuric acid, ammonia, nitric acid and fertilizers. It also examines how oxides of sulphur and nitrogen pollute the atmosphere, cause acid rain and eutrophication, and how these effects can be controlled.
\vspace{4pt}
\begin{multicols}{2}\small
\begin{itemize}
\item By the end of this chapter I can assign oxidation numbers to sulphur (-2, +2, +4, +6) and nitrogen (-3 to +5) in their common compounds and use them to classify redox behaviour.
\item By the end of this chapter I can describe the Contact process for manufacturing sulphuric acid, with conditions, equations and the role of the V2O5 catalyst.
\item By the end of this chapter I can state and explain the physical and chemical properties of concentrated and dilute sulphuric acid (acid, oxidant, dehydrating agent).
\item By the end of this chapter I can describe the Haber process for ammonia synthesis and justify the compromise conditions using Le Chatelier's principle.
\item By the end of this chapter I can describe the Ostwald process for nitric acid manufacture from ammonia.
\item By the end of this chapter I can write balanced equations for the reactions of ammonia (basicity, reducing agent, complex formation) and of nitric acid (acid and oxidizing agent).
\end{itemize}\end{multicols}}

\begin{sommaire}
\begin{multicols}{2}
\begin{enumerate}
\item Sulphur and its Variable Oxidation States
\item Sulphuric Acid: Manufacture, Properties and Uses
\item Oxides of Sulphur and Atmospheric Pollution
\item Nitrogen, Ammonia and Nitric Acid
\item Fertilizers and their Environmental Effects
\item Integration activity
\item Methods and worked examples
\item Exercises
\item Solutions to the exercises
\item Summary sheet
\end{enumerate}
\end{multicols}
\end{sommaire}

\begin{objectifs}
\begin{itemize}
\item By the end of this chapter I can assign oxidation numbers to sulphur (-2, +2, +4, +6) and nitrogen (-3 to +5) in their common compounds and use them to classify redox behaviour.
\item By the end of this chapter I can describe the Contact process for manufacturing sulphuric acid, with conditions, equations and the role of the V2O5 catalyst.
\item By the end of this chapter I can state and explain the physical and chemical properties of concentrated and dilute sulphuric acid (acid, oxidant, dehydrating agent).
\item By the end of this chapter I can describe the Haber process for ammonia synthesis and justify the compromise conditions using Le Chatelier's principle.
\item By the end of this chapter I can describe the Ostwald process for nitric acid manufacture from ammonia.
\item By the end of this chapter I can write balanced equations for the reactions of ammonia (basicity, reducing agent, complex formation) and of nitric acid (acid and oxidizing agent).
\item By the end of this chapter I can explain how SO2 and NOx form, how they produce acid rain and photochemical smog, and describe control measures.
\item By the end of this chapter I can explain the composition and manufacture of common NPK fertilizers and analyse their environmental effects, especially eutrophication.
\item By the end of this chapter I can carry out mole and percentage-yield calculations on industrial processes (Contact, Haber) in GCE-style problems.
\end{itemize}
\end{objectifs}

\begin{prerequis}
\begin{itemize}
\item Oxidation numbers and balancing redox equations (Lower Sixth: redox reactions)
\item Chemical equilibrium and Le Chatelier's principle (Lower Sixth)
\item Rates of reaction and catalysis (Lower Sixth)
\item Mole calculations, molar gas volume and percentage yield (Fifth Form / Lower Sixth)
\item Acid-base behaviour: acids, bases, salts and pH (Lower Sixth)
\end{itemize}
\end{prerequis}

\newpage

\coursec{Sulphur and its Variable Oxidation States}

In Douala's industrial zone and around the SONARA refinery in Limbe, residents sometimes notice a sharp, choking smell after gas flaring --- a sign of sulphur dioxide escaping into the air. Meanwhile, farmers in the Western Highlands spread ammonium- and nitrate-based fertilizers to boost maize and potato yields, yet streams downhill sometimes turn green with algae. Both stories involve two elements, sulphur and nitrogen, whose many \cle{oxidation states} make them industrially precious but environmentally dangerous. How are sulphuric acid and ammonia manufactured on a massive scale, and why do their by-products acidify rain and suffocate lakes? This chapter answers these questions by following sulphur and nitrogen from the factory to the field to the atmosphere. We begin with sulphur itself: where it comes from, what forms it takes, how it reacts, and how to keep track of its oxidation numbers.

\courssub{Occurrence and Extraction of Sulphur}

Sulphur is one of the most abundant elements in the Earth's crust. It occurs in two main forms:

\begin{itemize}
\item \textbf{Native (elemental) sulphur}, found in underground deposits, often associated with salt domes and volcanic regions;
\item \textbf{Combined sulphur}, in minerals such as iron pyrites \ce{FeS2}, galena \ce{PbS}, zinc blende \ce{ZnS} and gypsum \ce{CaSO4.2H2O}, and as hydrogen sulphide \ce{H2S} in crude oil and natural gas (``sour gas'').
\end{itemize}

Today, most of the world's sulphur no longer comes from mined deposits but is recovered from \cle{crude oil and natural gas refining}. Natural gas containing \ce{H2S} is treated so that the \ce{H2S} is separated and then converted to sulphur by the \cle{Claus process}: part of the \ce{H2S} is burnt to \ce{SO2}, which then reacts with the remaining \ce{H2S} over a catalyst:

\[
\ce{2H2S(g) + SO2(g) -> 3S(s) + 2H2O(g)}
\]

Where deep underground deposits of elemental sulphur exist, they are mined by the elegant \cle{Frasch process}, which exploits sulphur's low melting point (about \SI{115}{\ensuremath{{}^{\circ}\mathrm{C}}}).

\begin{experience}[The Frasch process]
Three concentric steel pipes are sunk into a borehole reaching the sulphur deposit.
\begin{enumerate}
\item \textbf{Superheated water} at about \SI{160}{\ensuremath{{}^{\circ}\mathrm{C}}} (kept liquid under pressure) is pumped down the \emph{outer} pipe; it melts the sulphur.
\item \textbf{Compressed air} is forced down the \emph{innermost} pipe.
\item The air mixes with the molten sulphur, making it light and frothy, so that the mixture of molten sulphur, water and air is pushed up the \emph{middle} pipe to the surface.
\end{enumerate}
At the surface the sulphur solidifies in large vats; it is about 99.5\% pure.
\end{experience}

\begin{popfigure}
\begin{tikzpicture}[scale=0.95]
% ground
\fill[popblueL] (-4.5,0) rectangle (4.5,0.3);
\draw[popline] (-4.5,0.3) -- (4.5,0.3);
\node[popmuted, anchor=west] at (-4.4,0.55) {\small surface};
% sulphur deposit
\fill[popgold!40] (-4.5,-4.2) rectangle (4.5,-2.6);
\draw[popline] (-4.5,-2.6) -- (4.5,-2.6);
\draw[popline] (-4.5,-4.2) -- (4.5,-4.2);
\node[popdark] at (2.9,-3.4) {\small sulphur deposit};
% outer pipe
\draw[very thick, popink] (-0.9,1.6) -- (-0.9,-3.8);
\draw[very thick, popink] (0.9,1.6) -- (0.9,-3.8);
% middle pipe
\draw[thick, popblue] (-0.5,1.6) -- (-0.5,-3.4);
\draw[thick, popblue] (0.5,1.6) -- (0.5,-3.4);
% inner pipe
\draw[thick, poporange] (-0.15,1.6) -- (-0.15,-3.0);
\draw[thick, poporange] (0.15,1.6) -- (0.15,-3.0);
% arrows
\draw[-stealth, very thick, popink] (-0.7,2.6) -- (-0.7,1.7);
\node[popink, anchor=west, align=left] at (-3.6,2.4) {\small superheated water \SI{160}{\ensuremath{{}^{\circ}\mathrm{C}}}};
\draw[popink] (-3.6,2.2) -- (-0.85,2.2);
\draw[-stealth, very thick, poporange] (0,2.9) -- (0,1.7);
\node[poporange, anchor=west] at (0.3,2.9) {\small compressed air};
\draw[-stealth, very thick, popblue] (0.7,1.7) -- (0.7,2.6);
\node[popblue, anchor=west] at (1.0,2.4) {\small molten sulphur + air + water};
% molten pool
\fill[popgold!70] (-1.4,-3.9) rectangle (1.4,-3.1);
\node[popdark] at (0,-3.5) {\small molten \ce{S}};
\end{tikzpicture}
\legende{The Frasch process: three concentric pipes deliver superheated water and compressed air to melt and lift underground sulphur.}
\end{popfigure}

\courssub{Allotropes of Sulphur}

If you dissolve sulphur in carbon disulphide and let the solvent evaporate slowly, you obtain beautiful yellow crystals. If instead you melt sulphur and let it crystallise just below its melting point, you get needle-shaped crystals of a \emph{different} structure. Same element, same physical state (solid), but different crystalline forms: this is allotropy.

\begin{definition}[Allotropy]
\cle{Allotropy} is the existence of an element in more than one form in the \emph{same physical state}. The different forms are called \cle{allotropes}; they differ in the arrangement or bonding of their atoms and therefore in some physical properties.
\end{definition}

Sulphur has two crystalline allotropes, both built from puckered \ce{S8} rings packed differently in the crystal:

\begin{itemize}
\item \textbf{Rhombic ($\alpha$) sulphur}: stable \emph{below} \SI{96}{\ensuremath{{}^{\circ}\mathrm{C}}}; octahedral-shaped crystals; density about \SI{2.06}{g.cm^{-3}}.
\item \textbf{Monoclinic ($\beta$) sulphur}: stable \emph{above} \SI{96}{\ensuremath{{}^{\circ}\mathrm{C}}} up to the melting point; needle-shaped (prismatic) crystals; density about \SI{1.96}{g.cm^{-3}}.
\end{itemize}

The temperature \SI{96}{\ensuremath{{}^{\circ}\mathrm{C}}} is the \cle{transition temperature}: below it rhombic sulphur is stable, above it monoclinic sulphur is stable. The change is slow, so each form can persist for a while in the other's range.

A third, non-crystalline form is \cle{plastic sulphur}: when boiling sulphur is poured into cold water (quenched), the \ce{S8} rings break into long chains that have no time to crystallise, giving a brown, rubbery, \emph{amorphous} solid which slowly reverts to rhombic sulphur on standing.

\begin{popfigure}
\begin{tikzpicture}[scale=1.0]
\node[draw=popblue, fill=popblueL, rounded corners, thick, align=center, minimum width=3.2cm] (rh) at (0,0) {\small rhombic ($\alpha$)\\ \small stable $< \SI{96}{\ensuremath{{}^{\circ}\mathrm{C}}}$};
\node[draw=poppurple, fill=poppurple!15, rounded corners, thick, align=center, minimum width=3.2cm] (mo) at (6.5,0) {\small monoclinic ($\beta$)\\ \small stable $> \SI{96}{\ensuremath{{}^{\circ}\mathrm{C}}}$};
\draw[-stealth, very thick, poporange] (rh.30) -- (mo.150) node[midway, above, popdark]{\small heat above \SI{96}{\ensuremath{{}^{\circ}\mathrm{C}}}};
\draw[-stealth, very thick, popgreen] (mo.210) -- (rh.330) node[midway, below, popdark]{\small cool below \SI{96}{\ensuremath{{}^{\circ}\mathrm{C}}}};
\node[draw=popgold, fill=popgold!25, rounded corners, thick, align=center, minimum width=3.2cm] (pl) at (3.25,-3.2) {\small plastic sulphur\\ \small (amorphous)};
\draw[-stealth, very thick, poppink] (mo.250) -- (pl.20) node[midway, right, popdark, align=left]{\small boil, then quench\\ \small in cold water};
\draw[-stealth, very thick, popmuted] (pl.160) -- (rh.290) node[midway, left, popdark, align=right]{\small slow, on\\ \small standing};
\end{tikzpicture}
\legende{Interconversion of the allotropes of sulphur. Rhombic and monoclinic sulphur interconvert at the transition temperature \SI{96}{\ensuremath{{}^{\circ}\mathrm{C}}}; quenching boiling sulphur gives amorphous plastic sulphur.}
\end{popfigure}

\courssub{Physical Properties of Sulphur}

\begin{aretenir}[Physical properties of sulphur]
\begin{itemize}
\item Yellow, brittle, non-metallic solid with no taste and no smell (the ``smell of sulphur'' is actually due to compounds such as \ce{SO2} or \ce{H2S}).
\item \textbf{Insoluble in water}, but \textbf{soluble in carbon disulphide} \ce{CS2} (and in organic solvents such as benzene).
\item Low melting point, about \SI{115}{\ensuremath{{}^{\circ}\mathrm{C}}}; poor conductor of heat and electricity.
\end{itemize}
\end{aretenir}

\courssub{Chemical Properties of Sulphur}

\textbf{Burning in air or oxygen.} Sulphur burns with a blue flame to give the choking gas sulphur dioxide:

\[
\ce{S(s) + O2(g) -> SO2(g)}
\]

\textbf{Reaction with metals.} Sulphur combines directly with many metals on heating to give \cle{sulphides}. A classic school demonstration is the vigorous, glowing reaction of iron filings with sulphur:

\[
\ce{Fe(s) + S(s) -> FeS(s)}
\]

Note that the product, iron(II) sulphide, is a \emph{compound}: it can no longer be separated by a magnet, unlike the original mixture.

\textbf{Reaction with hydrogen.} On passing hydrogen over boiling sulphur, hydrogen sulphide forms:

\[
\ce{H2(g) + S(l) -> H2S(g)}
\]

\courssub{Assigning Oxidation Numbers of Sulphur}

The key to understanding sulphur chemistry is its \cle{variable oxidation state}: sulphur appears at $-2$, $0$, $+2$, $+4$ and $+6$. We track these using oxidation numbers.

\begin{methode}[Assigning oxidation numbers]
\begin{enumerate}
\item A free (uncombined) element has oxidation number $0$: \ce{S}, \ce{O2}, \ce{Fe}.
\item Oxygen is usually $-2$ (except in peroxides, $-1$, and with fluorine).
\item Hydrogen is usually $+1$ (except in metal hydrides, $-1$).
\item The sum of oxidation numbers equals the overall charge of the species ($0$ for a neutral molecule).
\end{enumerate}
\end{methode}

\begin{exemplebox}[Oxidation number of S in \ce{H2SO4}]
Let $x$ be the oxidation number of sulphur. Then
\[
2(+1) + x + 4(-2) = 0 \quad\Longrightarrow\quad x = +6.
\]
Similarly, in \ce{S2O3^{2-}} (thiosulphate): $2x + 3(-2) = -2$, so $x = +2$ (an \emph{average} value).
\end{exemplebox}

\begin{center}
\begin{tabularx}{\linewidth}{|l|X|c|}
\hline
\rowcolor{popblueL} \textbf{Species} & \textbf{Name} & \textbf{Oxidation number of S} \\
\hline
\ce{H2S} & hydrogen sulphide & $-2$ \\
\hline
\ce{S} & elemental sulphur & $0$ \\
\hline
\ce{S2O3^{2-}} & thiosulphate ion & $+2$ (average) \\
\hline
\ce{SO2}, \ce{H2SO3} & sulphur dioxide, sulphurous acid & $+4$ \\
\hline
\ce{SO3}, \ce{H2SO4} & sulphur trioxide, sulphuric acid & $+6$ \\
\hline
\end{tabularx}
\end{center}

\begin{popfigure}
\begin{tikzpicture}[scale=1.0]
% ladder
\foreach \y/\ox/\cmp in {0/-2/\ce{H2S}, 1.4/0/\ce{S}, 2.8/+2/\ce{S2O3^{2-}}, 4.2/+4/\ce{SO2}, 5.6/+6/\ce{SO3}}{
  \draw[very thick, popline] (0,\y) -- (1.2,\y);
  \node[popdark, anchor=east] at (-0.15,\y) {\small $\ox$};
  \node[popink, anchor=west] at (1.4,\y) {\small \cmp};
}
\draw[very thick, popdark] (0.6,-0.4) -- (0.6,6.2);
% oxidation arrow
\draw[-stealth, ultra thick, poporange] (-1.6,0.2) -- (-1.6,5.8);
\node[poporange, rotate=90, align=center] at (-2.1,3) {\small oxidation (loss of \ce{e-})};
% reduction arrow
\draw[-stealth, ultra thick, popgreen] (3.6,5.8) -- (3.6,0.2);
\node[popgreen, rotate=90, align=center] at (4.15,3) {\small reduction (gain of \ce{e-})};
% extra labels
\node[popmuted, anchor=west] at (1.4,4.9) {\small \ce{H2SO3}, \ce{SO3^{2-}}};
\node[popmuted, anchor=west] at (1.4,6.3) {\small \ce{H2SO4}, \ce{SO4^{2-}}};
\end{tikzpicture}
\legende{The oxidation-state ladder of sulphur, from $-2$ to $+6$, with example species at each rung.}
\end{popfigure}

\courssub{Hydrogen Sulphide \ce{H2S}: Sulphur at $-2$}

\textbf{Preparation.} In the laboratory, \ce{H2S} is prepared by the action of dilute hydrochloric acid on iron(II) sulphide, usually in a Kipp's apparatus:

\[
\ce{FeS(s) + 2HCl(aq) -> FeCl2(aq) + H2S(g)}
\]

\begin{attention}[\ce{H2S} is highly toxic]
Hydrogen sulphide is \textbf{extremely poisonous} --- as toxic as hydrogen cyanide --- despite its familiar rotten-egg smell. Worse, it deadens the sense of smell at high concentration, so the disappearance of the smell is \emph{not} a sign of safety. \textbf{Never prepare or use \ce{H2S} outside a fume cupboard.}
\end{attention}

\textbf{Properties.}
\begin{itemize}
\item Colourless gas with a rotten-egg smell; denser than air; fairly soluble in water.
\item Its aqueous solution is a \textbf{weak dibasic acid}, ionising in two steps:
\[
\ce{H2S(aq) <=> H+(aq) + HS-(aq)} \qquad \ce{HS-(aq) <=> H+(aq) + S^{2-}(aq)}
\]
It forms two series of salts: hydrogensulphides (e.g.\ \ce{NaHS}) and sulphides (e.g.\ \ce{Na2S}).
\item It is a \textbf{strong reducing agent}, because sulphur at $-2$ can only be \emph{oxidised}. For example, it decolourises acidified potassium manganate(VII) and reduces iron(III) to iron(II):
\[
\ce{2FeCl3(aq) + H2S(g) -> 2FeCl2(aq) + 2HCl(aq) + S(s)}
\]
\end{itemize}

\textbf{Test for \ce{H2S}.} The gas turns filter paper moistened with \cle{lead(II) ethanoate} solution black, owing to the formation of lead(II) sulphide:

\[
\ce{Pb(CH3COO)2(aq) + H2S(g) -> PbS(s) + 2CH3COOH(aq)}
\]

\courssub{Sulphur Dioxide \ce{SO2} and Sulphurous Acid \ce{H2SO3}: Sulphur at $+4$}

\textbf{Preparation.} In the laboratory, \ce{SO2} is conveniently prepared by the action of a dilute acid on a sulphite:

\[
\ce{Na2SO3(s) + 2HCl(aq) -> 2NaCl(aq) + H2O(l) + SO2(g)}
\]

Industrially it is obtained by burning sulphur or roasting sulphide ores.

\textbf{Acidic oxide.} \ce{SO2} dissolves in water to form \cle{sulphurous acid}, a weak dibasic acid:

\[
\ce{SO2(g) + H2O(l) <=> H2SO3(aq)}
\]

It turns moist blue litmus paper red and reacts with alkalis to give sulphites \ce{SO3^{2-}} and hydrogensulphites \ce{HSO3-}.

\textbf{Reducing agent (the usual behaviour).} With sulphur at $+4$, \ce{SO2} is readily oxidised to $+6$, so it is a good \textbf{reducing agent}. It decolourises acidified potassium manganate(VII):

\[
\ce{2MnO4-(aq) + 5SO2(g) + 2H2O(l) -> 2Mn^{2+}(aq) + 5SO4^{2-}(aq) + 4H+(aq)}
\]

(purple $\rightarrow$ colourless), and turns acidified dichromate(VI) from orange to green.

\textbf{Oxidizing agent (with stronger reducers).} With a stronger reducing agent such as \ce{H2S}, \ce{SO2} behaves as an \emph{oxidizing agent}, being reduced to sulphur ($+4 \rightarrow 0$):

\[
\ce{SO2(g) + 2H2S(g) -> 3S(s) + 2H2O(l)}
\]

\begin{propriete}[Dual redox character of \ce{SO2}]
Sulphur dioxide is \textbf{both a reducing agent} (its usual role: S goes from $+4$ to $+6$, e.g.\ with \ce{KMnO4}, \ce{K2Cr2O7}, \ce{Cl2}) \textbf{and an oxidizing agent} towards stronger reducing agents such as \ce{H2S} (S goes from $+4$ to $0$). The oxidation number $+4$ is \emph{intermediate}, so it can move either way.
\end{propriete}

\textbf{Bleaching action.} \ce{SO2} bleaches moist vegetable

\coursec{Sulphuric Acid: Manufacture, Properties and Uses}

In the previous section we explored the rich redox chemistry of sulphur, noting how it comfortably adopts oxidation states from \cle{-2} (in \ce{H2S}) through \cle{+4} (in \ce{SO2}) to \cle{+6} (in \ce{SO3} and \ce{H2SO4}). The \cle{+6} state represents the fully oxidised, thermodynamically stable form of sulphur in the presence of oxygen. It is from this pinnacle that the single most important industrial chemical — sulphuric acid — is manufactured. A nation’s annual consumption of \ce{H2SO4} is still regarded by economists as one of the most reliable barometers of its industrial development: it is the lifeblood of fertilizer production, petroleum refining, metal processing, and the manufacture of countless everyday materials. In this section we trace the journey from elemental sulphur to the concentrated acid, justify every condition in the Contact process with Le Chatelier’s principle and kinetic theory, and then examine the remarkable trio of chemical personalities — acid, oxidiser, dehydrating agent — that make concentrated \ce{H2SO4} unique in the laboratory and in industry.

\courssub{The Contact Process: Industrial Manufacture of Sulphuric Acid}

The modern Contact process is a masterpiece of chemical engineering that converts sulphur (or sulphide ores) into sulphuric acid in three main stages, each optimised for yield, rate, and economic viability.

\begin{experience}[Outline of the Contact Process]
The overall transformation can be summarised as:
\begin{enumerate}
    \item \textbf{Production of \ce{SO2}:} Burning elemental sulphur or roasting sulphide ores (e.g. \ce{ZnS}, \ce{FeS2}, \ce{Cu2S}) in excess dry air.
    \item \textbf{Catalytic oxidation to \ce{SO3}:} \ce{2SO2(g) + O2(g) <=> 2SO3(g)} \quad $\Delta H = \SI{-197}{kJ\,mol^{-1}}$ (exothermic).
    \item \textbf{Absorption and dilution:} \ce{SO3} is absorbed in concentrated \ce{H2SO4} to form \cle{oleum} (\ce{H2S2O7}), which is then diluted with water to give the desired acid concentration.
\end{enumerate}
\end{experience}

\courssub{Step 1: Production and Purification of Sulphur Dioxide}

\begin{propriete}[Sources of \ce{SO2}]
Elemental sulphur (recovered from oil and gas desulphurisation, or mined by the Frasch process) is burnt in a furnace with a controlled supply of dry air:
\[
\ce{S(s) + O2(g) -> SO2(g)} \qquad \Delta H < 0
\]
Alternatively, metal sulphide ores are roasted:
\[
\ce{2ZnS(s) + 3O2(g) -> 2ZnO(s) + 2SO2(g)}
\]
\[
\ce{4FeS2(s) + 11O2(g) -> 2Fe2O3(s) + 8SO2(g)}
\]
The heat released is often recovered to generate steam for electricity or process heating.
\end{propriete}

\begin{attention}[Purification is Critical — The Catalyst Poison]
The gas leaving the burner/roaster contains dust, arsenic oxides (\ce{As2O3}), selenium, fluorine compounds, and moisture. Before entering the converter, it \textbf{must} be rigorously purified:
\begin{itemize}
    \item \textbf{Electrostatic precipitators / cyclones} remove dust.
    \item \textbf{Washing towers} with water or acid spray remove soluble impurities and cool the gas.
    \item \textbf{Drying towers} with concentrated \ce{H2SO4} remove moisture (water vapour would form acid mist later and corrode equipment).
    \item \textbf{Arsenic removal:} Arsenic impurities poison the \ce{V2O5} catalyst irreversibly. The gas is passed through a bed of \ce{Fe2O3} or \ce{Fe(OH)3} which absorbs \ce{As2O3} as ferric arsenate.
\end{itemize}
\textbf{GCE Trap:} A common exam question asks ``Why is the gas dried?'' Answer: To prevent corrosion and, crucially, to avoid formation of a stable \ce{H2SO4} mist in the absorber that would escape with the tail gases.
\end{attention}

\courssub{Step 2: Catalytic Oxidation of \ce{SO2} to \ce{SO3} — The Heart of the Process}

The reversible, exothermic reaction is:
\[
\ensuremath{\mathrm{2SO_{2}(g) + O_{2}(g) \rightleftharpoons [V_{2}O_{5}][450^\circ C,\ 1-2\ atm] 2SO_{3}(g)}} \qquad \Delta H = \SI{-197}{kJ\,mol^{-1}}
\]

\begin{popfigure}
\begin{tikzpicture}[scale=0.9, every node/.style={font=\small}]
% Axes
\draw[->] (0,0) -- (10,0) node[right] {Reaction coordinate};
\draw[->] (0,0) -- (0,5) node[above] {Energy};
% Reactants
\draw[thick, popdark] (0.5,3.5) -- (2,3.5) node[midway, above] {$\ce{2SO2 + O2}$};
% Products
\draw[thick, popdark] (8,1.5) -- (9.5,1.5) node[midway, above] {$\ce{2SO3}$};
% Uncatalysed path (high barrier)
\draw[thick, poporange, dashed] (2,3.5) .. controls (4,5.5) and (6,5.5) .. (8,1.5);
\node[poporange, text width=3cm, align=center] at (5,5.8) {$E_a$ (uncatalysed)};
% Catalysed path (lower barrier, two steps via V2O5)
\draw[thick, popblue] (2,3.5) .. controls (3.5,4.2) and (4.5,4.2) .. (5,3.5);
\draw[thick, popblue] (5,3.5) .. controls (6,4.0) and (7,4.0) .. (8,1.5);
\node[popblue, text width=3cm, align=center] at (5,4.5) {$E_a$ (catalysed)};
% Delta H arrow
\draw[<->, popgreen] (1,3.5) -- (1,1.5) node[midway, left] {$\Delta H = -197\ \mathrm{kJ\,mol^{-1}}$};
% Labels
\node[align=center] at (5,0.5) {Energy profile for \ce{2SO2 + O2 <=> 2SO3}};
\node[poporange, text width=3cm, align=center] at (2,5.2) {Without catalyst:\\high $E_a$, slow};
\node[popblue, text width=3cm, align=center] at (8,4.5) {With \ce{V2O5}:\\lower $E_a$, fast};
\end{tikzpicture}
\legende{Energy profile showing how the \ce{V2O5} catalyst provides an alternative reaction pathway with lower activation energy ($E_a$), speeding up both forward and reverse reactions equally without shifting the equilibrium position.}
\end{popfigure}

\begin{methode}[Choosing Industrial Conditions — The Le Chatelier / Kinetics / Economics Balance]
The Contact process is the textbook example of optimising an industrial equilibrium. For \ce{2SO2(g) + O2(g) <=> 2SO3(g)} ($\Delta H < 0$, $\Delta n_{\text{gas}} = -1$):

\begin{enumerate}
    \item \textbf{Temperature:} Low temperature favours the forward reaction (exothermic $\rightarrow$ higher equilibrium yield). \textbf{BUT} too low a temperature makes the rate unacceptably slow even with a catalyst. \textbf{Compromise:} \SIrange{400}{450}{\ensuremath{{}^{\circ}\mathrm{C}}}. At \SI{450}{\ensuremath{{}^{\circ}\mathrm{C}}} with \ce{V2O5}, the equilibrium conversion is $\approx 97\%$ and the rate is economically viable.
    \item \textbf{Pressure:} High pressure favours the forward reaction (fewer gas moles on product side). \textbf{BUT} the equilibrium yield at \SI{1}{atm} is already $>95\%$ at \SI{450}{\ensuremath{{}^{\circ}\mathrm{C}}}. Compressing large gas volumes is expensive (energy, thicker vessels). \textbf{Decision:} Operate at \SIrange{1}{2}{atm} (slightly above atmospheric to push gas through the plant). High pressure is \emph{not} used.
    \item \textbf{Catalyst:} Vanadium(V) oxide (\ce{V2O5}) on a silica support. It lowers $E_a$ dramatically, allowing the \SI{450}{\ensuremath{{}^{\circ}\mathrm{C}}} compromise. It is cheap, resistant to poisoning (if gas is purified), and has a long life. \textbf{Note:} The catalyst does \emph{not} change the equilibrium constant or the yield at equilibrium; it only helps the system reach equilibrium faster.
    \item \textbf{Concentration (Excess Oxygen):} Using excess air (typically $2\times$ stoichiometric \ce{O2}) shifts equilibrium to the right (Le Chatelier) and ensures near-complete conversion of the valuable \ce{SO2}. The unreacted gases are recycled.
    \item \textbf{Heat Management:} The reaction is exothermic. The converter contains multiple catalyst beds with \textbf{heat exchangers} between them. The hot exit gas preheats the incoming feed, maintaining the optimal temperature and recovering energy.
\end{enumerate}
\end{methode}

\begin{definition}[Oleum — Disulphuric(VI) Acid]
\textbf{Oleum} (fuming sulphuric acid) is a solution of sulphur trioxide in concentrated sulphuric acid, with the formula \ce{H2S2O7}. It is formed by the reaction:
\[
\ce{SO3(g) + H2SO4(l) -> H2S2O7(l)}
\]
Oleum is a stable, transportable intermediate. When water is added, it forms sulphuric acid:
\[
\ce{H2S2O7(l) + H2O(l) -> 2H2SO4(l)}
\]
The percentage labelling of oleum (e.g. \emph{20\% oleum}) refers to the mass of free \ce{SO3} present per 100 g of oleum, or equivalently the mass of water needed to convert 100 g of oleum to 100\% \ce{H2SO4}.
\end{definition}

\courssub{Step 3: Absorption of \ce{SO3} — Why Not Water Directly?}

\begin{attention}[The Mist Trap — Never Absorb \ce{SO3} Directly in Water]
If \ce{SO3} is bubbled into water, the reaction $\ce{SO3 + H2O -> H2SO4}$ is \textbf{violently exothermic}. The heat vaporises water and acid, creating a fine, stable \textbf{acid mist} (aerosol) that does not condense easily and escapes the absorber, causing pollution and product loss.
\textbf{Solution:} \ce{SO3} is absorbed in \textbf{98--99\% concentrated \ce{H2SO4}} in a packed absorption tower. The vapour pressure of water over concentrated \ce{H2SO4} is extremely low, so no mist forms. The product is oleum (\ce{H2S2O7}), which is then \textbf{carefully diluted} with water to the desired concentration (usually 98\% for storage/transport).
\end{attention}

\begin{popfigure}
\begin{tikzpicture}[scale=0.85, every node/.style={font=\small, align=center}]
% Nodes
\node[draw, rounded corners, fill=popblueL, minimum width=2.5cm, minimum height=1cm, text width=2.2cm] (S) at (0,0) {Sulphur /\\Sulphide ores};
\node[draw, rounded corners, fill=poporange, minimum width=2.5cm, minimum height=1cm, text width=2.2cm] (B) at (4,0) {Burner /\\Roaster};
\node[draw, rounded corners, fill=popgreenL, minimum width=2.5cm, minimum height=1cm, text width=2.2cm] (P) at (8,0) {Purification\\(Dust, As, H2O)};
\node[draw, rounded corners, fill=poppink, minimum width=3cm, minimum height=1.5cm, text width=2.8cm] (C) at (12,0) {Converter\\(\ce{V2O5}, 450°C,\\Heat exchangers)};
\node[draw, rounded corners, fill=popgold, minimum width=2.5cm, minimum height=1cm, text width=2.2cm] (A) at (16,0) {Absorption\\Tower};
\node[draw, rounded corners, fill=popblueL, minimum width=2.5cm, minimum height=1cm, text width=2.2cm] (D) at (20,0) {Dilution\\Tank};
\node[draw, rounded corners, fill=popgreen, minimum width=2.5cm, minimum height=1cm, text width=2.2cm] (Prod) at (20,-2.5) {Conc. \ce{H2SO4}\\(98\%)};
% Arrows
\draw[->, thick] (S) -- (B) node[midway, above] {+ Air};
\draw[->, thick] (B) -- (P) node[midway, above] {\ce{SO2} gas};
\draw[->, thick] (P) -- (C) node[midway, above] {Clean, dry \ce{SO2}/\ce{O2}};
\draw[->, thick] (C) -- (A) node[midway, above] {\ce{SO2}/\ce{SO3}/\ce{O2}};
\draw[->, thick] (A) -- (D) node[midway, above] {Oleum};
\draw[->, thick] (D) -- (Prod) node[midway, right] {+ Water};
% Recycle
\draw[->, thick, popmuted] (C) -- ++(0,-1.5) -| (P) node[pos=0.25, below, text width=2cm, align=center] {Unreacted gases\\recycled};
% Side streams
\draw[->, thick, popmuted] (B) -- ++(0,1.5) node[above] {Heat $\to$ Steam};
\draw[->, thick, popmuted] (C) -- ++(0,2) node[above] {Heat recovered};
\end{tikzpicture}
\legende{Flow diagram of the Contact process. Key features: purification unit protects the catalyst; converter uses multiple \ce{V2O5} beds with inter-stage heat exchangers; absorption in concentrated acid avoids mist formation; final dilution yields 98\% \ce{H2SO4}.}
\end{popfigure}

\courssub{Physical Properties of Concentrated Sulphuric Acid}

\begin{aretenir}[Key Physical Properties of Conc. \ce{H2SO4} (98\%)]
\begin{itemize}
    \item \textbf{Appearance:} Colourless, dense, oily liquid (viscosity $\approx 27\ \mathrm{mPa\cdot s}$ at \SI{20}{\ensuremath{{}^{\circ}\mathrm{C}}}).
    \item \textbf{Density:} \SI{1.84}{g\,cm^{-3}} — nearly twice as dense as water.
    \item \textbf{Boiling point:} \SI{337}{\ensuremath{{}^{\circ}\mathrm{C}}} (decomposes slightly to \ce{SO3} + \ce{H2O} vapour).
    \item \textbf{Hygroscopicity:} Extremely strong affinity for water — excellent drying agent for gases (except \ce{NH3} and \ce{H2S}, see below).
    \item \textbf{Dilution:} \textbf{Highly exothermic} ($\Delta H_{\text{dil}} \approx \SI{-80}{kJ\,mol^{-1}}$ for infinite dilution). The heat released can boil water and spatter concentrated acid.
\end{itemize}
\end{aretenir}

\begin{attention}[Safe Dilution Technique — ALWAYS Acid into Water]
\textbf{Rule:} ``Do as you \emph{oughta}, add acid to water.'' 
\begin{enumerate}
    \item Use a heat-resistant container (Pyrex beaker).
    \item Add the \textbf{water first}, then \textbf{slowly pour the concentrated acid} down the side of the beaker while stirring constantly.
    \item The dense acid sinks, and the heat is dissipated into the larger volume of water.
\end{enumerate}
\textbf{Never} pour water into concentrated acid — the water floats, boils violently at the interface, and ejects droplets of concentrated acid.
\end{attention}

\courssub{Chemical Properties: The Three Personalities of \ce{H2SO4}}

Sulphuric acid exhibits three distinct chemical behaviours depending on concentration and temperature. Recognising \emph{which role it is playing} is essential for predicting products.

\begin{propriete}[Personality 1: Dilute \ce{H2SO4} — A Typical Dibasic Acid]
When dilute, \ce{H2SO4} behaves as a strong diprotic acid, donating two protons stepwise:
\[
\ce{H2SO4(aq) -> H+(aq) + HSO4-(aq)} \quad (\text{strong, complete})
\]
\[
\ce{HSO4-(aq) <=> H+(aq) + SO4^2-(aq)} \quad (K_{a2} = 1.2 \times 10^{-2})
\]
It undergoes all typical acid reactions:
\begin{itemize}
    \item \textbf{With reactive metals (above H in electrochemical series):} \ensuremath{\mathrm{Mg + H_{2}SO_{4} \rightarrow  MgSO_{4} + H_{2}\uparrow}}
    \item \textbf{With bases (neutralisation):} \ce{2NaOH + H2SO4 -> Na2SO4 + 2H2O} (normal salt) \\
    \ce{NaOH + H2SO4 -> NaHSO4 + H2O} (acid salt / hydrogensulphate)
    \item \textbf{With carbonates/hydrogencarbonates:} \ensuremath{\mathrm{Na_{2}CO_{3} + H_{2}SO_{4} \rightarrow  Na_{2}SO_{4} + CO_{2}\uparrow + H_{2}O}}
    \item \textbf{With metal oxides:} \ce{CuO + H2SO4 -> CuSO4 + H2O}
\end{itemize}
\textbf{Note:} With \ce{BaCl2} or \ce{Pb(NO3)2}, dilute \ce{H2SO4} gives a \textbf{white precipitate} of \ce{BaSO4} or \ce{PbSO4} (insoluble sulphates) — a standard test for \ce{SO4^2-}.
\end{propriete}

\begin{propriete}[Personality 2: Hot Concentrated \ce{H2SO4} — An Oxidizing Agent]
When hot and concentrated, \ce{H2SO4} is reduced, with the sulphur oxidation state dropping from \cle{+6} to \cle{+4} (in \ce{SO2}) or lower. The acid provides both the oxidising power and the acidic medium.
\begin{center}
\begin{tabularx}{\linewidth}{|l|X|X|}\hline
\rowcolor{popblueL}
\textbf{Reducing Agent} & \textbf{Observation} & \textbf{Redox Half-Equations} \\ \hline
\ce{Cu(s)} (hot) & Blue solution (\ce{Cu^2+}), colourless \ce{SO2} gas & $\ce{Cu -> Cu^2+ + 2e-}$ \\
& & $\ce{SO4^2- + 4H+ + 2e- -> SO2 + 2H2O}$ \\
\textbf{Overall:} $\ensuremath{\mathrm{Cu + 2H_{2}SO_{4}(conc) \rightarrow [heat] CuSO_{4} + SO_{2}\uparrow + 2H_{2}O}}$ & & \\ \hline
\ce{C(s)} (charcoal) & \ce{SO2} gas, \ce{CO2} gas & $\ce{C + 2H2O -> CO2 + 4H+ + 4e-}$ \\
& & $\ce{2SO4^2- + 8H+ + 4e- -> 2SO2 + 4H2O}$ \\
\textbf{Overall:} $\ensuremath{\mathrm{C + 2H_{2}SO_{4}(conc) \rightarrow [heat] CO_{2}\uparrow + 2SO_{2}\uparrow + 2H_{2}O}}$ & & \\ \hline
\ce{S(s)} & \ce{SO2} gas (from both S and acid) & $\ce{S + 2H2O -> SO2 + 4H+ + 4e-}$ \\
\textbf{Overall:} $\ensuremath{\mathrm{S + 2H_{2}SO_{4}(conc) \rightarrow [heat] 3SO_{2}\uparrow + 2H_{2}O}}$ & & \\ \hline
\end{tabularx}
\end{center}
\textbf{Key point:} Cold concentrated \ce{H2SO4} does \emph{not} oxidise copper. Heat is required to overcome the kinetic barrier.
\end{propriete}

\begin{propriete}[Personality 3: Concentrated \ce{H2SO4} — A Dehydrating Agent]
Concentrated \ce{H2SO4} has an immense affinity for water, stripping the elements of water (H and O in 2:1 ratio) from compounds. This is \textbf{not} an acid-base reaction; it is a dehydration.
\begin{itemize}
    \item \textbf{Charring of carbohydrates (sugar, glucose, starch, paper, wood):}
    \[
    \ce{C12H22O11(s) ->[conc. H2SO4] 12C(s) + 11H2O(l/g)}
    \]
    The black porous carbon ``snake'' rises from the beaker as steam and \ce{SO2} (from slight oxidation) escape.
    \item \textbf{Drying gases:} Used in the laboratory to dry \ce{Cl2}, \ce{O2}, \ce{N2}, \ce{CO2}, \ce{SO2}, \ce{HCl}, etc.
\end{itemize}
\end{propriete}

\begin{attention}[Gases That CANNOT Be Dried with Conc. \ce{H2SO4}]
\begin{enumerate}
    \item \textbf{Ammonia (\ce{NH3}):} Acid-base reaction occurs: $\ce{NH3 + H2SO4 -> (NH4)2SO4}$ or $\ce{NH4HSO4}$. The gas is consumed.
    \item \textbf{Hydrogen sulphide (\ce{H2S}):} Redox reaction occurs: $\ensuremath{\mathrm{H_{2}S + H_{2}SO_{4} \rightarrow  S\downarrow + SO_{2}\uparrow + 2H_{2}O}}$. The gas is oxidised.
    \item \textbf{Alkenes (e.g. \ce{C2H4}):} Addition reaction / polymerisation occurs.
\end{enumerate}
\textbf{Alternative drying agents:} \ce{CaO} (for \ce{NH3}), \ce{P4O10} or silica gel (for \ce{H2S} and general use).
\end{attention}

\begin{popfigure}
\begin{tikzpicture}[scale=0.7, every node/.style={font=\small}]
% Beaker
\draw[thick] (2,0) -- (2,5) -- (6,5) -- (6,0);
\draw[thick] (1.5,0) -- (2,0) -- (2,-0.3) (6,0) -- (6.5,0) -- (6.5,-0.3);
% Sugar initially
\fill[white] (2.2,0.5) rectangle (5.8,2.5);
\draw[popmuted] (2.2,0.5) rectangle (5.8,2.5);
\node at (4,1.5) {Sucrose};
% Acid
\fill[popblueL!30] (2.2,0) rectangle (5.8,0.5);
\node at (4,0.25) {Conc. \ce{H2SO4}};
% Reaction: black snake rising
\draw[fill=popdark!80] (3,2.5) .. controls (2.5,3.5) and (3.5,4.5) .. (3,5.5) 
    .. controls (3.5,4.5) and (2.5,3.5) .. (3,2.5);
\draw[fill=popdark!80] (5,2.5) .. controls (5.5,3.5) and (4.5,4.5) .. (5,5.5) 
    .. controls (4.5,4.5) and (5.5,3.5) .. (5,2.5);
% Steam / gas bubbles
\foreach \y in {0.7,1.2,1.7,2.2} {
    \foreach \x in {2.5,3.5,4.5,5.5} {
        \draw[popblue, thick] (\x,\y) ellipse (0.15 and 0.08);
    }
}
% Labels
\node[align=center, text width=3cm] at (4,6.2) {Black porous carbon ``snake''\\rising from beaker};
\node[align=center, poporange, text width=4cm] at (8,3) {$\ce{C12H22O11 ->[conc. H2SO4] 12C + 11H2O}$\\(Dehydration, not oxidation)};
\node[align=center, popgreen, text width=3cm] at (8,1) {Steam (\ce{H2O})\\and traces of \ce{SO2}};
\end{tikzpicture}
\legende{The ``black snake'' demonstration: concentrated \ce{H2SO4} dehydrates sucrose, stripping H and O as water and leaving a rising column of elemental carbon. The reaction is highly exothermic, producing steam and a small amount of \ce{SO2} from partial oxidation.}
\end{popfigure}

\courssub{Uses of Sulphuric Acid}

\begin{aretenir}[Major Industrial Uses of \ce{H2SO4}]
\begin{enumerate}
    \item \textbf{Fertilizer manufacture (largest use, $\sim 60\%$):}
    \begin{itemize}
        \item \textbf{Superphosphate:} $\ce{Ca3(PO4)2 + 2H2SO4 -> Ca(H2PO4)2 + 2CaSO4}$ (single superphosphate, SSP).
        \item \textbf{Triple superphosphate:} $\ce{Ca3(PO4)2 + 3H3PO4 -> 3Ca(H2PO4)2}$ (uses phosphoric acid made from \ce{H2SO4}).
        \item \textbf{Ammonium sulphate:} $\ce{2NH3 + H2SO4 -> (NH4)2SO4}$.
    \end{itemize}
    \item \textbf{Lead-acid storage batteries (car batteries):} Electrolyte is $\approx 30\%$ \ce{H2SO4} (\SI{4.2}{mol\,dm^{-3}}). Reaction: $\ce{PbO2 + Pb + 2H2SO4 <=> 2PbSO4 + 2H2O}$.
    \item \textbf{Paints and pigments:} Manufacture of titanium dioxide (\ce{TiO2}, white pigment) via the sulphate process; production of barium sulphate (blanc fixe).
    \item \textbf{Detergents:} Sulphonation of alkylbenzenes to make linear alkylbenzene sulphonates (LAS).
    \item \textbf{Metal processing:} Pickling (removing oxide scale) of steel before galvanising or plating; leaching of copper and uranium ores.
    \item \textbf{Other:} Rayon (viscose) production, explosives (nitration mixture with \ce{HNO3}), petroleum refining (alkylation catalyst), laboratory reagent.
\end{enumerate}
\end{aretenir}

\courssub{GCE Exam Focus: Percentage Purity and Yield Calculations}

Contact process calculations are a staple of Paper 2. They typically involve:
\begin{itemize}
    \item Impure sulphur / sulphide ore $\rightarrow$ \ce{SO2} (account for purity).
    \item Equilibrium conversion in converter (often given as \% yield or calculated from $K_p$).
    \item Absorption efficiency.
    \item Final acid concentration / mass.
\end{itemize}

\begin{exoresolu}[Contact Process Yield Calculation]
\textbf{Statement:} \SI{100}{kg} of iron pyrites (\ce{FeS2}) containing \SI{85}{\%} pure \ce{FeS2} by mass is roasted in excess air. The \ce{SO2} produced is converted to \ce{SO3} in a Contact process converter with an equilibrium yield of \SI{96}{\%} (based on \ce{SO2} converted). The \ce{SO3} is absorbed in concentrated \ce{H2SO4} to form oleum, which is diluted to give \SI{98}{\%} \ce{H2SO4} by mass. Calculate the mass of \SI{98}{\%} acid produced. (Ar: Fe=56, S=32, O=16, H=1)

\tcblower
\textbf{Detailed Solution:}

\textbf{Step 1: Mass of pure \ce{FeS2}}
$M_r(\ce{FeS2}) = 56 + 64 = 120\ \mathrm{g\,mol^{-1}}$
Mass of pure \ce{FeS2} = $\SI{100}{kg} \times 0.85 = \SI{85}{kg} = \SI{85000}{g}$

\textbf{Step 2: Moles of \ce{FeS2} and \ce{SO2} produced}
Roasting: $\ce{4FeS2 + 11O2 -> 2Fe2O3 + 8SO2}$
Mole ratio $\ce{FeS2} : \ce{SO2} = 4 : 8 = 1 : 2$
$n(\ce{FeS2}) = \frac{85000}{120} = \SI{708.33}{mol}$
$n(\ce{SO2})\text{ (theoretical)} = 2 \times 708.33 = \SI{1416.67}{mol}$

\textbf{Step 3: Moles of \ce{SO3} formed (96\% conversion)}
$\ce{2SO2 + O2 <=> 2SO3}$ \quad Mole ratio $\ce{SO2} : \ce{SO3} = 1 : 1$
$n(\ce{SO3}) = 0.96 \times 1416.67 = \SI{1360.0}{mol}$

\textbf{Step 4: Moles of \ce{H2SO4} formed}
Absorption + dilution: $\ce{SO3 + H2O -> H2SO4}$ \quad (1:1 mole ratio)
$n(\ce{H2SO4}) = \SI{1360.0}{mol}$

\textbf{Step 5: Mass of pure \ce{H2SO4}}
$M_r(\ce{H2SO4}) = 2 + 32 + 64 = 98\ \mathrm{g\,mol^{-1}}$
Mass pure \ce{H2SO4} = $1360.0 \times 98 = \SI{133280}{g} = \SI{133.28}{kg}$

\textbf{Step 6: Mass of \SI{98}{\%} acid}
$\text{Mass of 98\% acid} = \frac{\text{mass pure acid}}{0.98} = \frac{133.28}{0.98} = \SI{136.0}{kg}$

\textbf{Answer:} \SI{136}{kg} of \SI{98}{\%} \ce{H2SO4} (to 3 s.f.).
\end{exoresolu}

\begin{exercice}[Practice: Sulphur Purity and Acid Yield]
\textbf{Statement:} A sample of crude sulphur of mass \SI{50.0}{kg} and purity \SI{92}{\%} is burnt in air. The \ce{SO2} is passed through a Contact process converter where the equilibrium conversion to \ce{SO3} is \SI{97.5}{\%}. The \ce{SO3} is absorbed in \SI{98}{\%} \ce{H2SO4} to form oleum, which is then diluted with water to produce \SI{95}{\%} \ce{H2SO4}. Calculate the mass of \SI{95}{\%} acid obtained. Assume no losses in absorption or dilution.
\end{exercice}

\begin{corrige}[Exercise 1]
\textbf{Solution:}
\begin{enumerate}
    \item Pure S mass = $50.0 \times 0.92 = \SI{46.0}{kg} = \SI{46000}{g}$.
    \item $n(\ce{S}) = 46000 / 32 = \SI{1437.5}{mol}$.
    \item $\ce{S + O2 -> SO2}$ : $n(\ce{SO2}) = \SI{1437.5}{mol}$.
    \item $\ce{2SO2 + O2 <=> 2SO3}$ (1:1 ratio): $n(\ce{SO3}) = 0.975 \times 1437.5 = \SI{1399.56}{mol}$.
    \item $\ce{SO3 + H2O -> H2SO4}$ : $n(\ce{H2SO4}) = \SI{1399.56}{mol}$.
    \item Mass pure \ce{H2SO4} = $1399.56 \times 98 = \SI{137157}{g} = \SI{137.16}{kg}$.
    \item Mass \SI{95}{\%} acid = $137.16 / 0.95 = \SI{144.4}{kg}$.
\end{enumerate}
\textbf{Answer:} \SI{144}{kg} (3 s.f.).
\end{corrige}

\courssub{Summary: Key Points to Retain}

\begin{aretenir}[Chapter Essentials — Sulphuric Acid]
\begin{itemize}
    \item \textbf{Contact Process:} $\ensuremath{\mathrm{S \rightarrow  SO_{2} \rightarrow [V_{2}O_{5}, 450^\circ C, 1-2 atm] SO_{3} \rightarrow [conc. H_{2}SO_{4}] H_{2}S_{2}O_{7} \rightarrow [H_{2}O] H_{2}SO_{4}}}$.
    \item \textbf{Conditions justified:} \SI{450}{\ensuremath{{}^{\circ}\mathrm{C}}} (yield/rate compromise), 1--2 atm (yield already high, compression costly), \ce{V2O5} catalyst (lowers $E_a$, no equilibrium shift), excess \ce{O2} (shifts equilibrium right).
    \item \textbf{Purification:} Essential to remove catalyst poisons (As) and prevent mist (drying).
    \item \textbf{Absorption:} \ce{SO3} into conc. \ce{H2SO4} (oleum) $\rightarrow$ dilution. \textbf{Never} direct into water (mist).
    \item \textbf{Three roles:} (1) Dilute = dibasic acid; (2) Hot conc. = oxidising agent (\ce{SO2} evolution); (3) Conc. = dehydrating agent (charring, gas drying).
    \item \textbf{Gases not dried by conc. \ce{H2SO4}:} \ce{NH3} (acid-base), \ce{H2S} (redox), alkenes (addition).
    \item \textbf{Safe dilution:} Acid into water, slowly, with stirring.
    \item \textbf{Uses:} Fertilizers (superphosphate, \ce{(NH4)2SO4}), batteries, pigments, detergents, metal pickling.
    \item \textbf{Calculations:} Track moles through each stage; apply \% purity at start, \% yield at converter, \% concentration at end.
\end{itemize}
\end{aretenir}

\begin{savaistu}[Sulphuric Acid in Cameroon]
Cameroon does not currently produce sulphuric acid on a large industrial scale; it is imported for use in the petroleum refinery (SONARA) for alkylation, in the aluminum smelter (ALUCAM) for anode production, and crucially for the manufacture of phosphate fertilizers (e.g. at the proposed Kribi fertilizer plant). The \emph{National Fertilizer Subsidy Programme} distributes NPK fertilizers whose phosphate component ultimately derives from sulphuric acid attack on phosphate rock. Thus, while the Contact process towers may not dot the Cameroonian landscape, the acid's fingerprints are on every bag of fertilizer that boosts maize yields in the North-West or cocoa in the South-West.
\end{savaistu}

\coursec{Oxides of Sulphur and Atmospheric Pollution}

In the previous section we saw that the Contact process deliberately converts \ce{SO2} into \ce{SO3} and then into \ce{H2SO4} inside a carefully controlled plant. Exactly the same chemistry occurs, uncontrolled, in the open air whenever sulphur-containing fuels are burned. The atmosphere becomes an accidental chemical reactor, and the product --- sulphuric acid --- falls back to Earth in the rain. This section studies that unwanted twin of the Contact process: the formation of sulphur oxides, their transformation into \cle{acid rain}, their effects on health and the environment, and the technologies used to limit the damage.

\courssub{Sources of Sulphur Dioxide}

\begin{definition}[Sulphur dioxide]
\cle{Sulphur dioxide}, \ce{SO2}, is a colourless, choking gas in which sulphur has the oxidation state $+4$. It is produced whenever elemental sulphur or sulphur compounds are burned in air:
\[ \ce{S(s) + O2(g) -> SO2(g)} \]
\end{definition}

The main sources of \ce{SO2} in the atmosphere are:

\begin{itemize}
\item \textbf{Burning of fossil fuels.} Coal and fuel oil contain sulphur (as organic sulphur compounds and as pyrite, \ce{FeS2}) at levels typically between $0.5\%$ and $5\%$ by mass. On combustion this sulphur is oxidised quantitatively to \ce{SO2}. Thermal power stations and industrial boilers are therefore major emitters.
\item \textbf{Roasting of sulphide ores.} In the extraction of metals such as copper and zinc, sulphide ores are roasted in air, releasing \ce{SO2}:
\[ \ce{2ZnS(s) + 3O2(g) -> 2ZnO(s) + 2SO2(g)} \]
In modern smelters this gas is captured and fed to a Contact plant --- a profitable way of combining metal extraction with sulphuric acid manufacture.
\item \textbf{Volcanic emissions.} Volcanoes inject very large quantities of \ce{SO2} directly into the atmosphere. This is a \emph{natural} source that no human policy can control; on a global scale it is comparable to industrial emissions.
\item \textbf{Gas flaring and refining.} Crude oil and natural gas contain sulphur compounds (notably \ce{H2S}). In the Cameroonian context, the refinery at Limbe (SONARA), gas flaring associated with offshore oil production, and thermal power stations burning fuel oil all contribute sulphur oxides to the local atmosphere.
\end{itemize}

\begin{savaistu}[Sulphur in Cameroon]
Natural gas from the Logbaba and Sanaga Sud fields contains hydrogen sulphide, \ce{H2S}, which must be removed before the gas is used --- both because \ce{H2S} is extremely toxic and because burning sulphurous gas would release \ce{SO2}. The removed sulphur can itself be recovered and sold to sulphuric acid plants.
\end{savaistu}

\courssub{From \ce{SO2} to \ce{SO3}: Atmospheric Oxidation}

\ce{SO2} does not remain unchanged in the air. Over hours to days it is oxidised to \ce{SO3} by two main routes:

\begin{enumerate}
\item \textbf{Photo-oxidation (gas phase).} Solar ultraviolet radiation initiates a radical chain reaction. The hydroxyl radical \ce{^.OH} (formed from ozone photolysis) attacks \ce{SO2}:
\[ \ce{SO2(g) + ^.OH(g) -> HOSO2^. (g)} \]
The \ce{HOSO2^.} radical reacts rapidly with \ce{O2} to give \ce{SO3} and \ce{HO2^.}, which regenerates \ce{^.OH}. The overall stoichiometry is:
\[ \ensuremath{\mathrm{2SO_{2}(g) + O_{2}(g) \rightarrow [h\nu] 2SO_{3}(g)}} \]
\item \textbf{Catalytic oxidation on particulate matter.} Soot, fly ash and metal-rich dust particles in polluted air provide surfaces on which \ce{SO2} and \ce{O2} react much faster. Water droplets containing dissolved metal ions (e.g.\ \ce{Fe^3+}, \ce{Mn^2+}) also catalyse the oxidation of dissolved \ce{SO2} to \ce{H2SO4}. This heterogeneous route is often dominant in humid, smoky air.
\end{enumerate}

Notice the parallel with the Contact process: there, a \ce{V2O5} catalyst at about $450\ \mathrm{^\circ C}$ performs the same oxidation efficiently; in the atmosphere, dust and sunlight do it slowly and for free.

\courssub{Acid Rain: Formation and pH}

\begin{propriete}[Natural acidity of rainwater]
Even in the absence of any pollution, rainwater is \emph{slightly acidic}, with $\mathrm{pH} \approx 5.6$, because it dissolves atmospheric carbon dioxide to form the weak acid carbonic acid:
\[ \ce{CO2(g) + H2O(l) <=> H2CO3(aq) <=> H+(aq) + HCO3-(aq)} \]
Acidity caused by \ce{SO2} and nitrogen oxides is \emph{additional} to this natural background.
\end{propriete}

Both sulphur oxides are acid anhydrides: they react with water to give acids.

\begin{align*}
\ce{SO2(g) + H2O(l) &-> H2SO3(aq)} && \text{sulphurous acid (weak, diprotic)}\\
\ce{SO3(g) + H2O(l) &-> H2SO4(aq)} && \text{sulphuric acid (strong, diprotic)}
\end{align*}

\ce{H2SO4} is a strong acid and dissociates extensively, so even small concentrations of \ce{SO3} lower the pH of rain dramatically. Since pH is logarithmic, a fall from $5.6$ to $4.0$ means the concentration of \ce{H+} has been multiplied by about $40$.

\begin{definition}[Acid rain]
\cle{Acid rain} is precipitation (rain, snow, fog or dew) with a pH below about $5.0$, caused by dissolved sulphur oxides (\ce{SO2}, \ce{SO3}) and nitrogen oxides (\ce{NO}, \ce{NO2}) which form sulphurous, sulphuric, nitrous and nitric acids in atmospheric water.
\end{definition}

\begin{popfigure}
\begin{center}
\begin{tikzpicture}[
  stage/.style={draw=popblue, fill=popblueL, rounded corners=2pt, align=center, text width=2.5cm, minimum height=1.1cm, font=\small},
  eff/.style={draw=poporange, fill=white, rounded corners=2pt, align=center, text width=2.2cm, minimum height=0.9cm, font=\small},
  arr/.style={-{Stealth[length=2.5mm]}, thick, popdark}]
\node[stage, align=center] (fuel) at (0,0){Sulphur in fuel\\(coal, fuel oil)};
\node[stage, align=center] (so2) at (4,0){\ce{SO2} emitted\\(chimneys, ores)};
\node[stage, align=center] (so3) at (8,0){Oxidation in air\\$\ce{SO2} \to \ce{SO3}$};
\node[stage, align=center] (rain) at (12,0){\ce{H2SO4} in rain\\(acid rain)};
\node[eff, align=center] (lake) at (9,-2.4){Acidified lakes\\fish die};
\node[eff, align=center] (soil) at (12,-2.4){Soil leaching\\\ce{Al^3+} toxicity};
\node[eff, align=center] (bldg) at (15,-2.4){Corroded\\buildings};
\draw[arr] (fuel) -- (so2);
\draw[arr] (so2) -- (so3);
\draw[arr] (so3) -- (rain);
\draw[arr, poporange] (rain) -- (lake);
\draw[arr, poporange] (rain) -- (soil);
\draw[arr, poporange] (rain) -- (bldg);
\end{tikzpicture}
\end{center}
\legende{From sulphur in fuel to environmental damage: the acid-rain pathway.}
\end{popfigure}

\begin{popfigure}
\begin{center}
\begin{tikzpicture}
\begin{axis}[
  ybar, bar width=0.9cm,
  width=11cm, height=6.5cm,
  ymin=0, ymax=8,
  ylabel={pH},
  symbolic x coords={Pure water, Normal rain, Acid rain, Lemon juice},
  xtick=data,
  x tick label style={font=\small, align=center},
  nodes near coords, nodes near coords style={font=\small},
  axis lines=left,
  ymajorgrids, grid style={popline!40}]
\addplot[fill=popblue] coordinates {(Pure water,7) (Normal rain,5.6) (Acid rain,4) (Lemon juice,2.3)};
\end{axis}
\end{tikzpicture}
\end{center}
\legende{Comparison of pH: pure water (7), normal rain ($\approx 5.6$), acid rain ($\approx 4$) and lemon juice ($\approx 2.3$).}
\end{popfigure}

\courssub{Environmental Effects of Acid Rain}

\textbf{Acidification of lakes and aquatic life.}\par
When acid rain drains into lakes with little natural buffering (e.g.\ lakes on granitic bedrock), the pH of the water falls. Below pH $5$, fish eggs fail to hatch and the gills of adult fish are damaged; below pH $4.5$, most fish species disappear and the lake becomes biologically impoverished --- clear, blue and almost lifeless.

\textbf{Soil leaching and aluminium toxicity.}\par
Acid rain dissolves and washes away (leaches) essential soil cations such as \ce{Ca^2+}, \ce{Mg^2+} and \ce{K+}, impoverishing the soil. Worse, the extra \ce{H+} attacks clay minerals and releases \ce{Al^3+} ions into soil water:
\[ \ce{Al(OH)3(s) + 3H+(aq) -> Al^3+(aq) + 3H2O(l)} \]
\ce{Al^3+} is toxic to plant roots and to fish: it damages root membranes, blocks nutrient uptake, and in gills it causes mucus clogging. Forest decline observed in heavily polluted regions is largely attributed to this combination of nutrient loss and aluminium toxicity.

\textbf{Corrosion of buildings and metals.}\par
Limestone and marble (both \ce{CaCO3}) react directly with the acid:
\[ \ce{CaCO3(s) + H2SO4(aq) -> CaSO4(s) + H2O(l) + CO2(g)} \]
The surface of statues and carved facades literally dissolves, and fine detail is lost. Iron and steel structures are attacked too, with the acid accelerating rusting --- a real economic cost for bridges, roofs and railway lines.

\textbf{Damage to crops and forests.}\par
Acid deposition damages the waxy cuticle of leaves, leaches nutrients from foliage, and weakens trees so that they succumb more easily to drought, frost, insects and disease. Crop yields fall on acidified soils unless lime is regularly added.

\courssub{Direct Health Effects of \ce{SO2} and Secondary Pollutants}

\ce{SO2} itself is a powerful respiratory irritant. It dissolves in the moisture lining the airways to form sulphurous acid, causing coughing, constriction of the bronchi and, after prolonged exposure, chronic \cle{bronchitis}. Asthmatics are particularly sensitive: even short exposures to elevated \ce{SO2} levels can trigger severe attacks.

In addition, \ce{SO2} is converted in the atmosphere into \cle{sulphate aerosols} --- tiny suspended particles of \ce{(NH4)2SO4} and \ce{H2SO4} droplets. These \emph{secondary pollutants}:
\begin{itemize}
\item scatter sunlight and reduce \textbf{visibility} (the haze over industrial regions);
\item penetrate deep into the lungs when inhaled, contributing to respiratory and cardiovascular disease;
\item can remain airborne for days and travel hundreds of kilometres, exporting pollution far from its source.
\end{itemize}

\begin{attention}[Greenhouse effect $\neq$ acid rain]
Do not confuse the \textbf{greenhouse effect} (caused mainly by \ce{CO2} and \ce{CH4}, leading to global warming) with \textbf{acid rain} (caused by \ce{SO2} and nitrogen oxides). They are different problems with different gases, different chemistries and different remedies. Examiners penalise answers that mix the two.
\end{attention}

\courssub{Control Measures}

\textbf{1. Use of low-sulphur fuels.} Switching from high-sulphur coal or fuel oil to natural gas or low-sulphur coal reduces emissions at the source.

\textbf{2. Hydrodesulphurization of fuels.} In refineries, sulphur compounds are removed from liquid fuels by reaction with hydrogen over a catalyst; the sulphur leaves as \ce{H2S}, which is then converted to elemental sulphur (Claus process) and sold.

\textbf{3. Flue-gas desulphurization (FGD).} The most important end-of-pipe technology: the exhaust gases are ``scrubbed'' with a slurry of limestone (\ce{CaCO3}) or lime (\ce{CaO}) before release.

\begin{methode}[Flue-gas desulphurization: the chemistry]
The scrubber neutralises the acidic \ce{SO2} with a basic calcium compound:
\begin{enumerate}
\item With limestone:
\[ \ce{CaCO3(s) + SO2(g) -> CaSO3(s) + CO2(g)} \]
\item With quicklime:
\[ \ce{CaO(s) + SO2(g) -> CaSO3(s)} \]
\item The calcium sulphite is then usually oxidised with air to gypsum, a saleable product used in plasterboard:
\[ \ce{2CaSO3(s) + O2(g) -> 2CaSO4(s)} \]
\end{enumerate}
Up to $90\%$ of the \ce{SO2} can be removed this way.
\end{methode}

\begin{popfigure}
\begin{center}
\begin{tikzpicture}[font=\small]
\draw[thick, popdark, fill=popblueL] (2,0) rectangle (5,6);
\draw[-{Stealth[length=2.5mm]}, very thick, poporange] (0,1) -- (2,1);
\node[align=center, poporange] at (1,1.7) {dirty flue gas\\(\ce{SO2})};
\draw[-{Stealth[length=2.5mm]}, very thick, popgreen] (5,5.4) -- (7,5.4);
\node[align=center, popgreen] at (6,6.1) {clean gas\\out};
\draw[-{Stealth[length=2.5mm]}, very thick, popblue] (3.5,7.6) -- (3.5,5.6);
\node[align=center, popblue] at (3.5,8.3) {limestone slurry\\spray \ce{CaCO3}};
\foreach \x/\y in {2.6/4.6, 3.4/4.2, 4.2/4.7, 2.9/3.4, 3.8/3.1, 4.4/3.8, 2.5/2.4, 3.6/2.2, 4.3/2.5}
  \fill[popblue] (\x,\y) circle (0.05);
\draw[-{Stealth[length=2.5mm]}, very thick, poppurple] (3.5,0) -- (3.5,-1.2);
\node[align=center, poppurple] at (3.5,-1.9) {\ce{CaSO3}/\ce{CaSO4}\\sludge out};
\node[align=center] at (3.5,5.1) {scrubbing\\chamber};
\end{tikzpicture}
\end{center}
\legende{Schematic of a limestone wet scrubber: flue gas enters, slurry spray absorbs \ce{SO2}, clean gas leaves.}
\end{popfigure}

\textbf{4. Tall chimneys.} Building very tall chimneys carries the plume high into the atmosphere so that ground-level concentrations near the plant are reduced.

\begin{attention}[Tall chimneys disperse, they do not remove]
A common mistake is to claim that tall chimneys ``reduce pollution''. They do \textbf{not} reduce the total \ce{SO2} emitted: the same mass of gas is released, merely spread over a wider area and carried further downwind --- sometimes across national borders. Dispersal is not removal; only scrubbing or cleaner fuels actually cut emissions.
\end{attention}

\begin{exoresolu}[Sulphur dioxide from burning fuel oil]
A thermal power station burns $1000\ \mathrm{kg}$ of fuel oil per hour containing $2.0\%$ sulphur by mass. Calculate (a) the mass of \ce{SO2} produced per hour, assuming all the sulphur is converted; (b) the mass of limestone, \ce{CaCO3}, needed per hour to scrub this \ce{SO2} completely. ($A_r$: C $=12$, O $=16$, S $=32$, Ca $=40$.)
\tcblower
\textbf{(a) Mass of \ce{SO2}.}

Mass of sulphur burned per hour:
\[ m(\ce{S}) = 1000 \times \frac{2.0}{100} = 20\ \mathrm{kg} \]
From $\ce{S + O2 -> SO2}$, one mole of S gives one mole of \ce{SO2}.
\[ n(\ce{S}) = \frac{20\,000\ \mathrm{g}}{32\ \mathrm{g\,mol^{-1}}} = 625\ \mathrm{mol} \]
\[ m(\ce{SO2}) = 625 \times (32 + 2\times16) = 625 \times 64 = 40\,000\ \mathrm{g} = 40\ \mathrm{kg} \]

\textbf{(b) Mass of limestone.}

From $\ce{CaCO3 + SO2 -> CaSO3 + CO2}$, one mole of \ce{CaCO3} removes one mole of \ce{SO2}.
\[ n(\ce{CaCO3}) = 625\ \mathrm{mol} \]
\[ m(\ce{CaCO3}) = 625 \times (40 + 12 + 3\times16) = 625 \times 100 = 62\,500\ \mathrm{g} = 62.5\ \mathrm{kg} \]
So the station emits $40\ \mathrm{kg}$ of \ce{SO2} per hour and would need $62.5\ \mathrm{kg}$ of pure limestone per hour to scrub it all.
\end{exoresolu}

\begin{savaistu}[A GCE extra: nitrogen oxides also make acid rain]
Sulphur is not the only culprit. Nitrogen oxides (\ce{NO} and \ce{NO2}, together called \ensuremath{\mathrm{NO_{x}}}), produced in vehicle engines and lightning, dissolve in rain to give nitric acid, \ce{HNO3}. In many industrialised countries \ensuremath{\mathrm{NO_{x}}} now contributes as much to acid rain as \ce{SO2}. We shall meet these oxides again when we study nitrogen compounds and fertilisers in the next sections.
\end{savaistu}

\begin{exercice}[Checking understanding]
\begin{enumerate}
\item Explain why ``pure'' unpolluted rainwater is not neutral.
\item Write balanced equations for the formation of sulphuric acid in the atmosphere starting from sulphur in coal.
\item State two environmental effects and one health effect of sulphur dioxide pollution.
\item A smelter roasts zinc blende according to $\ce{2ZnS + 3O2 -> 2ZnO + 2SO2}$. Calculate the volume of \ce{SO2} at s.t.p.\ released per tonne ($10^6\ \mathrm{g}$) of \ce{ZnS} roasted. (Molar volume at s.t.p.\ $= 22.4\ \mathrm{dm^3\,mol^{-1}}$; $A_r$: Zn $=65$, S $=32$.)
\end{enumerate}
\end{exercice}

\begin{corrige}[Exercise n°1]
\begin{enumerate}
\item Rainwater dissolves atmospheric \ce{CO2}, forming the weak acid \ce{H2CO3}; hence its pH is about $5.6$ rather than $7$.
\item $\ce{S + O2 -> SO2}$; $\ce{2SO2 + O2 -> 2SO3}$; $\ce{SO3 + H2O -> H2SO4}$.
\item Environmental: acidification of lakes killing fish; corrosion of limestone buildings (or soil leaching, forest damage). Health: respiratory irritation, bronchitis, aggravation of asthma.
\item $M(\ce{ZnS}) = 65 + 32 = 97\ \mathrm{g\,mol^{-1}}$; $n(\ce{ZnS}) = 10^6/97 \approx 10\,309\ \mathrm{mol}$. From the equation, $n(\ce{SO2}) = n(\ce{ZnS})$, so $V(\ce{SO2}) = 10\,309 \times 22.4 \approx 2.31\times10^{5}\ \mathrm{dm^3}$ at s.t.p.
\end{enumerate}
\end{corrige}

\begin{aretenir}[Key points]
\begin{itemize}
\item \ce{SO2} comes from burning sulphur-containing fossil fuels, roasting sulphide ores, volcanoes and gas flaring.
\item In air, \ce{SO2} is oxidised (photo-oxidation via \ce{^.OH} radicals, catalysis on particles) to \ce{SO3}; the oxides dissolve in rain to give \ce{H2SO3} and \ce{H2SO4}.
\item Normal rain has pH $\approx 5.6$ (dissolved \ce{CO2}); acid rain has pH $< 5$.
\item Effects: acidified lakes, soil leaching with \ce{Al^3+} toxicity, corrosion of \ce{CaCO3} buildings and metals, crop and forest damage, respiratory illness, sulphate aerosol haze.
\item Control: low-sulphur fuels, hydrodesulphurization, flue-gas scrubbing ($\ce{CaCO3 + SO2 -> CaSO3 + CO2}$). Tall chimneys only disperse pollution.
\item Never confuse acid rain (\ce{SO2}, \ensuremath{\mathrm{NO_{x}}}) with the greenhouse effect (\ce{CO2}, \ce{CH4}).
\end{itemize}
\end{aretenir}

\coursec{Nitrogen, Ammonia and Nitric Acid}

In the previous section we saw how sulphur, burned in fuels, escapes into the atmosphere as \ce{SO2} and returns to Earth as acid rain. But sulphur is not the only non-metal whose chemistry shapes our environment. Every breath you take is about $78\ \%$ nitrogen gas, \ce{N2} — yet this enormous reservoir is almost useless to plants and animals in its elemental form. Turning inert \ce{N2} into ammonia and nitric acid is one of the greatest achievements of industrial chemistry, and it feeds the world through fertilisers. In this section we study nitrogen itself, its two most important compounds — ammonia \ce{NH3} and nitric acid \ce{HNO3} — and the two industrial processes (Haber and Ostwald) that the GCE examiners love to test.

\courssub{Nitrogen in Nature and its Oxidation States}

\textbf{An inert but abundant gas.} Air contains about $78\ \%$ \ce{N2} by volume. Nitrogen is remarkably unreactive at ordinary temperatures: it does not burn, does not support combustion, and reacts with very few reagents directly. The reason is structural. The two nitrogen atoms are joined by a \cle{triple bond}, \ce{N#N}, with a very high bond dissociation energy (about \SI{945}{kJ.mol^{-1}}). Breaking this bond requires a large activation energy, so reactions of \ce{N2} are slow even when they are energetically favourable. This is why nitrogen fixation — natural or industrial — always demands drastic conditions (lightning, high temperature and pressure, or specialised enzymes).

\textbf{The nitrogen cycle in outline.} Although \ce{N2} is inert, nitrogen atoms circulate continuously between the atmosphere, the soil and living organisms:

\begin{itemize}
\item \textbf{Nitrogen fixation:} conversion of \ce{N2} into nitrogen compounds. Naturally by lightning (\ce{N2 + O2 -> 2NO} at very high temperature, then \ce{NO -> NO2 -> HNO3} washed into the soil as nitrates) and by \emph{Rhizobium} bacteria in the root nodules of legumes (beans, groundnuts — well known to Cameroonian farmers). Industrially by the Haber process.
\item \textbf{Nitrification:} soil bacteria (\emph{Nitrosomonas}, \emph{Nitrobacter}) oxidise ammonium compounds first to nitrites \ce{NO2-}, then to nitrates \ce{NO3-}, which plants absorb.
\item \textbf{Assimilation and decay:} plants convert nitrates into proteins; animals eat the plants; decomposition and excretion return ammonium compounds to the soil.
\item \textbf{Denitrification:} other bacteria reduce nitrates back to \ce{N2}, completing the cycle.
\end{itemize}

\begin{definition}[Nitrogen fixation]
\cle{Nitrogen fixation} is the conversion of atmospheric nitrogen gas, \ce{N2}, into nitrogen compounds (such as \ce{NH3}, \ce{NO} or nitrates) that plants can absorb and use. It occurs \textbf{naturally} (lightning discharges; \emph{Rhizobium} bacteria in legume root nodules) and \textbf{industrially} (the Haber process).
\end{definition}

\textbf{Oxidation states from $-3$ to $+5$.} Like sulphur, nitrogen displays a wide range of oxidation states — in fact an even wider one, every integer value from $-3$ to $+5$. You must be able to assign the oxidation state of N in any compound using the rules (O is $-2$, H is $+1$).

\begin{center}
\begin{tabularx}{\linewidth}{|c|X|X|}
\hline
\rowcolor{popblueL} \textbf{Oxidation state} & \textbf{Compound} & \textbf{Name} \\
\hline
$-3$ & \ce{NH3}, \ce{NH4+} & ammonia, ammonium ion \\
\hline
$-2$ & \ce{N2H4} & hydrazine \\
\hline
$-1$ & \ce{NH2OH} & hydroxylamine \\
\hline
$0$ & \ce{N2} & nitrogen gas \\
\hline
$+1$ & \ce{N2O} & dinitrogen monoxide \\
\hline
$+2$ & \ce{NO} & nitrogen monoxide \\
\hline
$+3$ & \ce{N2O3}, \ce{HNO2} & dinitrogen trioxide, nitrous acid \\
\hline
$+4$ & \ce{NO2}, \ce{N2O4} & nitrogen dioxide \\
\hline
$+5$ & \ce{HNO3}, \ce{N2O5}, \ce{NO3-} & nitric acid, dinitrogen pentoxide, nitrate \\
\hline
\end{tabularx}
\end{center}

\begin{popfigure}
\begin{center}
\begin{tikzpicture}[scale=1.0]
\draw[->, thick, popdark] (0,-0.5) -- (0,9.6) node[above, align=center, text width=2.6cm] {oxidation state of N};
\foreach \y/\n/\f/\c in {0/{-3}/{\ce{NH3}, \ce{NH4+}}/popgreen, 1/{-2}/{\ce{N2H4}}/popgreen, 2/{-1}/{\ce{NH2OH}}/popgreen, 3/{0}/{\ce{N2}}/popmuted, 4/{+1}/{\ce{N2O}}/popblue, 5/{+2}/{\ce{NO}}/popblue, 6/{+3}/{\ce{N2O3}, \ce{HNO2}}/popblue, 7/{+4}/{\ce{NO2}, \ce{N2O4}}/poporange, 8/{+5}/{\ce{HNO3}, \ce{N2O5}}/poporange}{
  \draw[thick, popdark] (-0.15,\y) -- (0.15,\y);
  \node[left] at (-0.2,\y) {\n};
  \node[right, \c] at (0.4,\y) {\f};
  \fill[\c] (0,\y) circle (2.2pt);
}
\node[right, align=left, text width=5.2cm, popmuted] at (3.6,4) {N is \textbf{reduced} downwards\\[2pt] N is \textbf{oxidised} upwards};
\draw[->, very thick, poporange] (3.2,4.6) -- (3.2,8.6);
\draw[->, very thick, popgreen] (3.2,3.4) -- (3.2,0.4);
\end{tikzpicture}
\end{center}
\legende{Oxidation-state ladder of nitrogen, from $-3$ to $+5$, with the important oxides and oxyacids.}
\end{popfigure}

\begin{attention}[Assigning oxidation states]
In \ce{N2O}, do not write N $= +2$ by treating it like \ce{NO}. Solve $2x + (-2) = 0$, giving $x = +1$ (an average value). Similarly in \ce{NH4+}: $x + 4(+1) = +1$, so $x = -3$. Always use the \emph{overall} charge of the species.
\end{attention}

\courssub{Ammonia: Laboratory Preparation}

Ammonia, \ce{NH3}, is a colourless gas with a characteristic pungent, choking smell (think of smelling salts or stale urine, where it forms by decomposition of urea). In the laboratory it is prepared by heating an \textbf{ammonium salt with a base} — classically ammonium chloride with calcium hydroxide (slaked lime):

\[ \ce{2NH4Cl(s) + Ca(OH)2(s) ->[heat] CaCl2(s) + 2NH3(g) + 2H2O(l)} \]

The mixture is heated gently in a hard-glass test tube tilted slightly downwards (so that the water produced does not run back onto the hot glass). Ammonia is \textbf{lighter than air} and \textbf{extremely soluble in water}, so it is collected by \cle{upward delivery} (downward displacement of air) — never over water.

\begin{methode}[Drying ammonia]
Ammonia is a basic gas, so it must be dried with a \textbf{basic drying agent}: quicklime, \ce{CaO}, placed in a drying tower.
\begin{enumerate}
\item \textbf{Do not use concentrated \ce{H2SO4}:} it is acidic and reacts: \ce{2NH3 + H2SO4 -> (NH4)2SO4}.
\item \textbf{Do not use anhydrous \ce{CaCl2}:} it combines with ammonia to form an addition compound, \ce{CaCl2.8NH3}.
\item \textbf{Do not use \ce{P2O5}} (also acidic).
\item \textbf{Use \ce{CaO}:} it is basic and absorbs water without reacting with \ce{NH3}.
\end{enumerate}
\end{methode}

\courssub{The Haber Process: Industrial Manufacture of Ammonia}

The Haber process is the industrial answer to nitrogen fixation — and arguably the most important chemical process ever invented, since the fertilisers it enables feed billions of people.

\textbf{Raw materials.} Nitrogen is obtained by fractional distillation of liquid air. Hydrogen is obtained from natural gas by steam reforming: \ce{CH4(g) + H2O(g) -> CO(g) + 3H2(g)}.

\textbf{The equilibrium.} The synthesis is a reversible, exothermic reaction:

\[ \ce{N2(g) + 3H2(g) <=> 2NH3(g)} \qquad \Delta H = -92\ \mathrm{kJ\,mol^{-1}} \]

\begin{popfigure}
\begin{center}
\begin{tikzpicture}[font=\small,
 box/.style={draw=popblue, thick, rounded corners, fill=popblueL, align=center, text width=3.0cm, minimum height=1.0cm}]
\node[box, align=center] (mix) at (0,0){\ce{N2} + \ce{H2}\\ (1 : 3), compressed\\ to $\sim 200$ atm};
\node[box, align=center] (conv) at (5.2,0){Converter\\ Fe catalyst\\ $\sim 450\ ^\circ$C};
\node[box, align=center] (cool) at (10.4,0){Cooler / condenser\\ \ce{NH3} liquefies};
\node[box, fill=popgreenL, draw=popgreen] (nh3) at (10.4,-2.6) {Liquid \ce{NH3} removed};
\draw[->, very thick, popdark] (mix) -- (conv);
\draw[->, very thick, popdark] (conv) -- (cool);
\draw[->, very thick, popdark] (cool) -- (nh3);
\draw[->, very thick, poporange] (cool.north) ++(0,0.55) -- ++(-5.2,0) node[midway, above, align=center, text width=4.2cm] {unreacted \ce{N2} and \ce{H2} recycled} -- (mix.north);
\end{tikzpicture}
\end{center}
\legende{Flow diagram of the Haber process, with the recycle loop for unreacted gases.}
\end{popfigure}

\textbf{Choosing the conditions — a compromise.} Apply Le Chatelier's principle to \ce{N2(g) + 3H2(g) <=> 2NH3(g)}:

\begin{itemize}
\item \textbf{Pressure:} 4 moles of gas $\rightarrow$ 2 moles of gas, so \emph{high} pressure shifts equilibrium towards \ce{NH3} and also increases the rate. About \SI{200}{atm} is used; higher pressures give better yields but the cost of compression and of building safe plant rises sharply.
\item \textbf{Temperature:} the forward reaction is exothermic, so \emph{low} temperature favours the yield of \ce{NH3}. But at low temperature the reaction is extremely slow (the strong \ce{N#N} bond gives a huge activation energy). A temperature of about \SI{450}{\ensuremath{{}^{\circ}\mathrm{C}}} is chosen as a \cle{compromise} between yield and rate.
\item \textbf{Catalyst:} finely divided \textbf{iron} speeds up attainment of equilibrium without changing the position of equilibrium.
\item \textbf{Removal of product:} the gas mixture leaving the converter contains only about $15\ \%$ \ce{NH3}; it is cooled so that ammonia liquefies and is removed, and the unreacted \ce{N2}/\ce{H2} are \textbf{recycled}, giving an overall conversion of about $98\ \%$.
\end{itemize}

\begin{attention}[The classic Haber trap]
Lowering the temperature would give a \emph{higher equilibrium yield} of \ce{NH3} — but the reaction would be far too slow to be economic. \SI{450}{\ensuremath{{}^{\circ}\mathrm{C}}} is \textbf{not} the temperature of maximum yield; it is a compromise between equilibrium yield and reaction rate. Saying "the catalyst increases the yield" is also wrong: a catalyst increases only the \emph{rate}, never the equilibrium yield.
\end{attention}

\courssub{Properties of Ammonia}

\textbf{Extreme solubility and alkalinity — the fountain experiment.} Ammonia is one of the most soluble gases known: about $700$ volumes of \ce{NH3} dissolve in one volume of water at room temperature. This is dramatically shown by the \cle{fountain experiment}. A round-bottomed flask is filled with dry ammonia and fitted with a jet tube dipping into water containing a little phenolphthalein. When a few drops of water enter the flask, the ammonia dissolves, the pressure inside collapses, and water is driven up through the jet as a spectacular fountain — turning \textbf{pink}, proving that the solution is alkaline.

\begin{popfigure}
\begin{center}
\begin{tikzpicture}[scale=0.9, font=\small]
\draw[thick, popdark] (0,0) circle (1.5);
\node[popblue] at (0,0.6) {dry \ce{NH3} gas};
\draw[thick, popdark] (-0.25,1.5) -- (-0.25,2.1) -- (0.25,2.1) -- (0.25,1.5);
\draw[thick, popdark] (0,2.1) -- (0,3.4);
\draw[thick, popdark] (0,-1.5) -- (0,-2.6);
\draw[thick, popblue] (-0.12,-2.6) -- (-0.12,-1.0) -- (0,-0.7) -- (0.12,-1.0) -- (0.12,-2.6);
\draw[thick, popdark] (-1.3,-2.6) -- (-1.3,-3.4) -- (1.3,-3.4) -- (1.3,-2.6);
\draw[popblue, thick] (-1.15,-2.75) -- (1.15,-2.75);
\node[popblue, align=center, text width=3.4cm] at (3.6,-3.0) {water + phenolphthalein};
\draw[->, thick, poppink] (0,-2.5) -- (0,-1.2);
\node[poppink, align=center, text width=3.2cm] at (4.1,-1.6) {pink fountain $\Rightarrow$ alkaline solution};
\end{tikzpicture}
\end{center}
\legende{The ammonia fountain experiment: dissolution of \ce{NH3} lowers the pressure and water jets in, turning pink with phenolphthalein.}
\end{popfigure}

The solution is alkaline because ammonia is a \textbf{weak base}, partially accepting a proton from water:

\[ \ce{NH3(aq) + H2O(l) <=> NH4+(aq) + OH-(aq)} \]

\begin{propriete}[Ammonia as a Lewis base]
The nitrogen atom in \ce{NH3} carries a \textbf{lone pair of electrons}, which it can donate to a proton or to a metal ion. Ammonia is therefore a \cle{Lewis base} (and a Br\o nsted base). This single fact explains both its alkalinity in water and its ability to form \textbf{complex ions} with transition-metal cations. (No formal proof is examinable; you must be able to \emph{use} the idea.)
\end{propriete}

\textbf{Precipitation of metal hydroxides.} Because aqueous ammonia supplies \ce{OH-} ions, it precipitates insoluble metal hydroxides from salt solutions — a key tool in qualitative analysis:

\begin{align*}
\ce{Cu^{2+}(aq) + 2OH-(aq) &-> Cu(OH)2(s)} \quad \text{pale blue precipitate}\\
\ce{Fe^{3+}(aq) + 3OH-(aq) &-> Fe(OH)3(s)} \quad \text{red-brown precipitate}
\end{align*}

\textbf{Complex formation.} With \emph{excess} ammonia solution, some hydroxide precipitates redissolve as ammine complexes. The classic case is copper(II):

\[ \ce{Cu(OH)2(s) + 4NH3(aq) -> [Cu(NH3)4]^{2+}(aq) + 2OH-(aq)} \]

The deep blue colour of tetraamminecopper(II), \ce{[Cu(NH3)4]^{2+}}, is a confirmatory test for \ce{Cu^{2+}}. (With silver chloride, ammonia forms \ce{[Ag(NH3)2]+}, dissolving the precipitate — another useful test.)

\textbf{Ammonia as a reducing agent.} Nitrogen in \ce{NH3} is in its lowest oxidation state ($-3$), so it can only be oxidised: ammonia is a reducing agent.

\begin{itemize}
\item It reduces heated copper(II) oxide: \ce{2NH3(g) + 3CuO(s) -> 3Cu(s) + N2(g) + 3H2O(l)} (black CuO turns brown-pink).
\item It burns in oxygen (and is oxidised catalytically over platinum): \ce{4NH3(g) + 3O2(g) -> 2N2(g) + 6H2O(g)}; over hot Pt the product is mainly \ce{NO} — the basis of the Ostwald process below.
\end{itemize}

\textbf{Uses of ammonia.} The great majority (about $80\ \%$) goes into \textbf{fertiliser manufacture} (ammonium salts, urea, nitrates). It is also used as a \textbf{refrigerant} gas in large cold stores and in household \textbf{cleaning agents}, and as the raw material for nitric acid.

\courssub{The Ostwald Process: Manufacture of Nitric Acid}

Nitric acid is made by the catalytic oxidation of ammonia — elegantly linking the two compounds of this section.

\begin{popfigure}
\begin{center}
\begin{tikzpicture}[font=\small,
 box/.style={draw=popblue, thick, rounded corners, fill=popblueL, align=center, text width=3.1cm, minimum height=1.1cm}]
\node[box, align=center] (mix) at (0,0){Mixer\\ \ce{NH3} + air\\ (about 1 : 10)};
\node[box, align=center] (conv) at (4.6,0){Converter\\ Pt--Rh gauze\\ $\sim 900\ ^\circ$C};
\node[box, align=center] (cool) at (9.2,0){Cooler / oxidiser\\ \ce{NO -> NO2}};
\node[box, align=center] (abs) at (13.8,0){Absorption tower\\ \ce{NO2} + \ce{H2O}};
\draw[->, very thick, popdark] (mix) -- (conv);
\draw[->, very thick, popdark] (conv) -- (cool);
\draw[->, very thick, popdark] (cool) -- (abs);
\node[box, fill=popgreenL, draw=popgreen] (hno3) at (13.8,-2.6) {\ce{HNO3(aq)} collected};
\draw[->, very thick, popdark] (abs) -- (hno3);
\draw[->, very thick, poporange] (abs.north) ++(0,0.6) -- ++(-9.2,0) node[midway, above, align=center, text width=4.4cm] {\ce{NO} recycled to oxidiser} -- (cool.north);
\end{tikzpicture}
\end{center}
\legende{Flow diagram of the Ostwald process for the manufacture of nitric acid.}
\end{popfigure}

\begin{enumerate}
\item \textbf{Catalytic oxidation of ammonia.} A mixture of ammonia and excess air passes over a platinum--rhodium gauze at about \SI{900}{\ensuremath{{}^{\circ}\mathrm{C}}}:
\[ \ensuremath{\mathrm{4NH_{3}(g) + 5O_{2}(g) \rightarrow [Pt/Rh,\ 900\ ^\circ C] 4NO(g) + 6H_{2}O(g)}} \qquad \Delta H < 0 \]
The gauze glows red-hot: the reaction is exothermic and, once started, maintains its own temperature.
\item \textbf{Oxidation of \ce{NO}.} The gases are cooled and the nitrogen monoxide reacts with more oxygen:
\[ \ce{2NO(g) + O2(g) -> 2NO2(g)} \]
\item \textbf{Absorption in water.} The nitrogen dioxide is absorbed in water in the presence of excess air:
\[ \ce{3NO2(g) + H2O(l) -> 2HNO3(aq) + NO(g)} \]
The \ce{NO} produced is \textbf{recycled} and re-oxidised, so the overall yield is very high. The acid obtained (about $60\ \%$) is concentrated by distillation.
\end{enumerate}

\courssub{Properties of Nitric Acid}

\textbf{A strong acid.} Pure nitric acid is a colourless fuming liquid. In water it is fully ionised:

\[ \ce{HNO3(aq) + H2O(l) -> H3O+(aq) + NO3-(aq)} \]

It shows all the typical reactions of a strong acid:

\begin{align*}
\ce{HNO3(aq) + NaOH(aq) &-> NaNO3(aq) + H2O(l)} && \text{base}\\
\ce{2HNO3(aq) + CuO(s) &-> Cu(NO3)2(aq) + H2O(l)} && \text{basic oxide}\\
\ce{2HNO3(aq) + Na2CO3(aq) &-> 2NaNO3(aq) + H2O(l) + CO2(g)} && \text{carbonate}
\end{align*}

\textbf{An oxidising acid.} Here lies the crucial difference from hydrochloric or dilute sulphuric acid. The nitrate group, with nitrogen in its maximum oxidation state $+5$, is itself a powerful oxidising agent. Consequently, when nitric acid attacks metals, it is the \textbf{nitrate that is reduced}, not the \ce{H+} — so \textbf{hydrogen is generally not evolved}.

With copper, the products depend on concentration:

\begin{align*}
\text{concentrated:}\quad & \ce{Cu(s) + 4HNO3(conc) -> Cu(NO3)2(aq) + 2NO2(g) + 2H2O(l)}\\
& \text{(brown fumes of \ce{NO2}; solution turns blue-green)}\\[4pt]
\text{dilute:}\quad & \ce{3Cu(s) + 8HNO3(dil) -> 3Cu(NO3)2(aq) + 2NO(g) + 4H2O(l)}\\
& \text{(colourless \ce{NO}, which browns in air as it forms \ce{NO2})}
\end{align*}

\begin{attention}[No hydrogen with nitric acid!]
A very common GCE mistake is to write \ce{Cu + 2HNO3 -> Cu(NO3)2 + H2}. This is \textbf{wrong}. Nitric acid is an \emph{oxidising acid}: with dilute acid copper gives \ce{NO}, with concentrated acid it gives \ce{NO2} — never hydrogen (except with very dilute acid and very reactive metals such as Mg or Mn, a rare exception). Learn the two copper equations above; they are set again and again.
\end{attention}

\textbf{Passivation.} Concentrated nitric acid renders aluminium and iron \cle{passive}: a thin, impermeable layer of oxide forms on the metal surface and protects it from further attack. This is why concentrated \ce{HNO3} can be transported in aluminium tanks.

\textbf{Decomposition in light.} Concentrated nitric acid slowly decomposes, especially when warm or exposed to light:

\[ \ce{4HNO3(aq) ->[light] 4NO2(g) + O2(g) + 2H2O(l)} \]

The dissolved brown \ce{NO2} gives old nitric acid its characteristic \textbf{yellow colour}. For this reason the acid is stored in brown bottles.

\courssub{Worked Examples and GCE Examination Focus}

\begin{exoresolu}[Redox half-equations for the reduction of nitric acid]
Write the balanced half-equation, in acid solution, for the reduction of the nitrate ion (a) to \ce{NO2}, (b) to \ce{NO}. State the change in oxidation state of nitrogen in each case.
\tcblower
Balance O with \ce{H2O}, then H with \ce{H+}, then charge with electrons.

(a) \ce{NO3- -> NO2}. Add one \ce{H2O} on the right for the extra O, then $2\ce{H+}$ on the left:
\[ \ce{NO3-(aq) + 2H+(aq) + e- -> NO2(g) + H2O(l)} \]
Nitrogen goes from $+5$ to $+4$: a gain of \textbf{one} electron.

(b) \ce{NO3- -> NO}. Add $2\ce{H2O}$ on the right, then $4\ce{H+}$ on the left:
\[ \ce{NO3-(aq) + 4H+(aq) + 3e- -> NO(g) + 2H2O(l)} \]
Nitrogen goes from $+5$ to $+2$: a gain of \textbf{three} electrons.

Combining (b) with the oxidation of copper, \ce{Cu -> Cu^{2+} + 2e-} (multiplied by 3, while (b) is multiplied by 2), regenerates the full equation \ce{3Cu + 8HNO3 -> 3Cu(NO3)2 + 2NO + 4H2O} — an excellent check.
\end{exoresolu}

\begin{exoresolu}[Yield calculation on the Haber process]
A converter is fed with \SI{100}{mol} of \ce{N2} and \SI{300}{mol} of \ce{H2}. At equilibrium the mixture contains $20\ \%$ conversion of the nitrogen. Calculate (a) the amount of ammonia formed, (b) the total number of moles of gas at equilibrium.
\tcblower
(a) \ce{N2} reacted $= 0.20 \times 100 = \SI{20}{mol}$. From \ce{N2 + 3H2 -> 2NH3}, ammonia formed $= 2 \times 20 = \boxed{\SI{40}{mol}}$.

(b) \ce{H2} reacted $= 3 \times 20 = \SI{60}{mol}$. Remaining: \ce{N2}: $100 - 20 = \SI{80}{mol}$; \ce{H2}: $300 - 60 = \SI{240}{mol}$; \ce{NH3}: \SI{40}{mol}.
\[ n_{\text{total}} = 80 + 240 + 40 = \boxed{\SI{360}{mol}} \]
Note that the total has decreased from $400$ to $360$ mol — consistent with 4 volumes of reactants giving only 2 volumes of product.
\end{exoresolu}

\begin{exoresolu}[The Ostwald process from ammonia to nitric acid]
Calculate the maximum mass of nitric acid that can be obtained from \SI{1.70}{kg} of ammonia, assuming quantitative conversion. ($M(\ce{NH3}) = 17$, $M(\ce{HNO3}) = 63$.)
\tcblower
Combining the three steps, each nitrogen atom of \ce{NH3} ends up in one molecule of \ce{HNO3} (the \ce{NO} by-product is recycled), so the overall mole ratio is $\ce{NH3} : \ce{HNO3} = 1 : 1$.
\[ n(\ce{NH3}) = \frac{1700}{17} = \SI{100}{mol} \quad\Rightarrow\quad n(\ce{HNO3}) = \SI{100}{mol} \]
\[ m(\ce{HNO3}) = 100 \times 63 = \SI{6300}{g} = \boxed{\SI{6.30}{kg}} \]
In practice, yields of about $96\ \%$ are achieved thanks to the recycling of \ce{NO}.
\end{exoresolu}

\begin{exercice}[Test yourself]
\begin{enumerate}
\item Assign the oxidation state of nitrogen in \ce{N2O}, \ce{NO}, \ce{N2O3}, \ce{NO2} and \ce{HNO3}.
\item Explain, with equations, why ammonia cannot be dried with concentrated sulphuric acid or anhydrous calcium chloride.
\item State the conditions used in the Haber process and justify each one using Le Chatelier's principle and kinetic arguments.
\item Write balanced equations for the three steps of the Ostwald process.
\item A blue solution of copper(II) sulphate is treated with a few drops of ammonia solution, then with excess. Describe and explain all observations, with equations.
\end{enumerate}
\end{exercice}

\begin{corrige}[Exercise n° 1]
\begin{enumerate}
\item \ce{N2O}: $+1$; \ce{NO}: $+2$; \ce{N2O3}: $+3$; \ce{NO2}: $+4$; \ce{HNO3}: $+5$.
\item \ce{H2SO4} is acidic: \ce{2NH3 + H2SO4 -> (NH4)2SO4}. \ce{CaCl2} forms an addition compound: \ce{CaCl2 + 8NH3 -> CaCl2.8NH3}. Only a basic dryer such as \ce{CaO} is suitable.
\item Fe catalyst (increases rate, not yield); $\sim \SI{450}{\ensuremath{{}^{\circ}\mathrm{C}}}$ (compromise: lower T favours yield of the exothermic reaction but is too slow); $\sim \SI{200}{atm}$ (fewer moles of gas on the product side, so high pressure increases both rate and yield); \ce{NH3} liquefied off and unreacted gases recycled.
\item \ce{4NH3 + 5O2 ->[Pt/Rh] 4NO + 6H2O}; \ce{2NO + O2 -> 2NO2}; \ce{3NO2 + H2O -> 2HNO3 + NO}.
\item A few drops: pale blue precipitate of \ce{Cu(OH)2}: \ce{Cu^{2+} + 2OH- -> Cu(OH)2}. Excess ammonia: precipitate dissolves to a deep blue solution of the complex \ce{[Cu(NH3)4]^{2+}}: \ce{Cu(OH)2 + 4NH3 -> [Cu(NH3)4]^{2+} + 2OH-}.
\end{enumerate}
\end{corrige}

\begin{aretenir}[Key points of this section]
\begin{itemize}
\item \ce{N2} is inert because of the strong \ce{N#N} triple bond; nitrogen shows every oxidation state from $-3$ (\ce{NH3}) to $+5$ (\ce{HNO3}).
\item Laboratory \ce{NH3}: heat \ce{NH4Cl} with \ce{Ca(OH)2}; dry over \ce{CaO}; collect by upward delivery.
\item Haber process: \ce{N2 + 3H2 <=> 2NH3}, $\Delta H = \SI{-92}{kJ.mol^{-1}}$; Fe catalyst, $\sim \SI{450}{\ensuremath{{}^{\circ}\mathrm{C}}}$, $\sim \SI{200}{atm}$ — a compromise between rate and yield.
\item \ce{NH3} is a weak base and a Lewis base: it precipitates metal hydroxides, forms complexes such as deep blue \ce{[Cu(NH3)4]^{2+}}, and acts as a reducing agent.
\item Ostwald process: \ce{NH3 -> NO -> NO2 -> HNO3} over Pt--Rh at $\sim \SI{900}{\ensuremath{{}^{\circ}\mathrm{C}}}$, with \ce{NO} recycled.
\item \ce{HNO3} is a strong acid \emph{and} an oxidising agent: with Cu it gives \ce{NO2} (conc.) or \ce{NO} (dil.), never \ce{H2}; it passivates Al and Fe and turns yellow in light due to dissolved \ce{NO2}.
\end{itemize}
\end{aretenir}

\begin{savaistu}[Fritz Haber: feeding the world]
Fritz Haber developed the ammonia synthesis in Germany in the early 20th century and received the Nobel Prize in Chemistry for it. It has been estimated that a large fraction of the nitrogen atoms in the human body today passed through a Haber plant — the fertilisers that grew our food were made from ammonia. Every bag of NPK or urea fertiliser used on the maize farms of the North-West and West Regions of Cameroon carries a little of this chemistry. In the next section we look closely at these fertilisers — and at what happens when too much of them washes into our rivers and lakes.
\end{savaistu}

\coursec{Fertilizers and their Environmental Effects}

In the previous section we followed nitrogen from the inert gas in the atmosphere through the Haber process to ammonia, and on through the Ostwald process to nitric acid. We now ask the obvious question: \emph{why} does industry make millions of tonnes of ammonia and nitric acid every year? The answer, in one word, is \cle{fertilizers}. The same chemistry that feeds the world, however, can also poison rivers, lakes and drinking water if it is misused. This final section closes the loop: we see why crops need fertilizers, how the main fertilizers are made from the very chemicals of this chapter, how to compare them quantitatively, and what happens when they escape into the environment.

\courssub{Why Fertilizers Are Needed}

Walk through any intensively farmed zone of Cameroon — the banana plantations of the Littoral, the maize fields of the West Region — and you will see the same story. Natural soil contains the mineral nutrients plants need, but every harvest \emph{removes} nutrients in the crop and carries them away. Unlike a natural ecosystem, where dead leaves and organisms decompose and return their minerals to the soil, a farm exports its nutrients to the market. After a few seasons of continuous cropping, yields fall dramatically unless the lost nutrients are replaced.

Three elements are needed in the largest amounts and are therefore the most commonly deficient. They are remembered by the letters \cle{N}, \cle{P} and \cle{K}.

\begin{center}
\begin{tabularx}{\linewidth}{|l|X|X|}
\hline
\rowcolor{popblueL} \textbf{Nutrient} & \textbf{Role in the plant} & \textbf{Deficiency symptom} \\
\hline
Nitrogen (N) & Component of proteins, chlorophyll and nucleic acids; drives leaf and stem growth & Yellowing of older leaves (chlorosis), stunted growth \\
\hline
Phosphorus (P) & Part of ATP (energy transfer), DNA and cell membranes; promotes root development and flowering & Poor root growth, purplish leaves, delayed maturity \\
\hline
Potassium (K) & Regulates water balance and stomata; activates enzymes; improves disease and drought resistance & Scorched leaf edges, weak stems, poor fruit quality \\
\hline
\end{tabularx}
\end{center}

\begin{definition}[Fertilizer]
A \cle{fertilizer} is any substance, natural or manufactured, added to soil or water to supply one or more of the nutrient elements essential for plant growth, principally nitrogen, phosphorus and potassium.
\end{definition}

\courssub{Types of Fertilizers}

\textbf{Straight fertilizers} supply a single major nutrient. The important examples for this syllabus are:

\begin{itemize}
\item \cle{Ammonium nitrate}, \ce{NH4NO3} — a nitrogen fertilizer, highly soluble, fast acting.
\item \cle{Ammonium sulphate}, \ce{(NH4)2SO4} — supplies nitrogen and sulphur; useful on sulphur-deficient soils.
\item \cle{Urea}, \ce{CO(NH2)2} — the most concentrated solid nitrogen fertilizer in common use.
\item \cle{Superphosphate} — a mixture of \ce{Ca(H2PO4)2} and \ce{CaSO4}, the main phosphorus fertilizer; the calcium dihydrogenphosphate is water-soluble and therefore available to plants, whereas the original phosphate rock is not.
\end{itemize}

\textbf{Compound (NPK) fertilizers} contain two or three of the major nutrients blended in one product, e.g.\ a ``15--15--15'' or ``20--10--10'' grade.

\begin{attention}[What the NPK numbers really mean]
The three numbers on a fertilizer bag give the percentages by mass of \textbf{N}, \textbf{\ce{P2O5}} and \textbf{\ce{K2O}} \emph{equivalents} — \emph{not} the percentages of elemental phosphorus and potassium. A 20--10--10 fertilizer contains \SI{20}{\percent} N, \SI{10}{\percent} \ce{P2O5} equivalent and \SI{10}{\percent} \ce{K2O} equivalent. This is a classic GCE trap: never convert the label directly into elemental P or K without being asked to.
\end{attention}

\courssub{Manufacture: Linking the Industrial Chemistry of this Chapter}

Notice how the fertilizers are made directly from the products of the Haber, Ostwald and Contact processes — this is why those three processes dominate the chemical industry.

\textbf{Ammonium nitrate} is made by neutralising ammonia with nitric acid (the Ostwald product):
\[ \ce{NH3(g) + HNO3(aq) -> NH4NO3(aq)} \]
followed by evaporation of the water and prilling (forming small pellets from droplets of the melt) to give the solid fertilizer.

\textbf{Ammonium sulphate} is made by neutralising ammonia with sulphuric acid (the Contact product):
\[ \ce{2NH3(g) + H2SO4(aq) -> (NH4)2SO4(aq)} \]

\textbf{Superphosphate} is made by treating insoluble phosphate rock, \ce{Ca3(PO4)2}, with sulphuric acid:
\[ \ce{Ca3(PO4)2(s) + 2H2SO4(aq) -> Ca(H2PO4)2(s) + 2CaSO4(s)} \]
The product mixture (soluble \ce{Ca(H2PO4)2} + gypsum, \ce{CaSO4}) is sold directly as ``single superphosphate''. The chemistry is simply an acid--base displacement: the stronger acid (\ce{H2SO4}) protonates the phosphate to give the soluble dihydrogenphosphate.

\courssub{Comparing Fertilizers: Percentage Nitrogen}

A farmer (and an examiner!) wants to know which fertilizer delivers the most nitrogen per franc. The measure is the \cle{percentage by mass of nitrogen} in the compound.

\begin{methode}[Calculating the percentage of nitrogen in a fertilizer]
\begin{enumerate}
\item Write the formula and count the number of nitrogen atoms per formula unit.
\item Calculate the mass of nitrogen in one mole: $n(\text{N atoms}) \times 14$.
\item Calculate the molar mass $M$ of the compound.
\item Apply: $\%N = \dfrac{\text{mass of N per mole}}{M} \times 100$.
\end{enumerate}
\end{methode}

\begin{exoresolu}[Percentage nitrogen in ammonium nitrate and urea]
Calculate the percentage by mass of nitrogen in (a) ammonium nitrate, \ce{NH4NO3}; (b) urea, \ce{CO(NH2)2}. Hence state which supplies more nitrogen per tonne. ($A_r$: H = 1, C = 12, N = 14, O = 16.)
\tcblower
\textbf{(a) Ammonium nitrate, \ce{NH4NO3}.} The formula contains \textbf{2} nitrogen atoms (one in \ce{NH4+}, one in \ce{NO3-}).
\begin{align*}
M(\ce{NH4NO3}) &= 14 + 4(1) + 14 + 3(16) = 80\ \mathrm{g\,mol^{-1}}\\
\%N &= \frac{2 \times 14}{80} \times 100 = \frac{28}{80}\times 100 = \boxed{35\ \%}
\end{align*}
\textbf{(b) Urea, \ce{CO(NH2)2}.} The formula contains \textbf{2} nitrogen atoms.
\begin{align*}
M(\ce{CO(NH2)2}) &= 12 + 16 + 2\,[14 + 2(1)] = 60\ \mathrm{g\,mol^{-1}}\\
\%N &= \frac{2 \times 14}{60} \times 100 = \frac{28}{60}\times 100 \approx \boxed{46.7\ \%}
\end{align*}
\textbf{Conclusion:} urea is the more concentrated source of nitrogen — one tonne of urea supplies about \SI{467}{kg} of N against \SI{350}{kg} for one tonne of ammonium nitrate. (For comparison, \ce{(NH4)2SO4} gives $\frac{28}{132}\times 100 \approx 21.2\ \%$.)
\end{exoresolu}

\begin{attention}[Count the nitrogen atoms first!]
The most common error is forgetting that \ce{NH4NO3} contains \emph{two} nitrogen atoms and writing $\%N = \frac{14}{80}\times 100 = 17.5\ \%$. Always count atoms before computing.
\end{attention}

\begin{popfigure}
\begin{center}
\begin{tikzpicture}
\begin{axis}[
    width=0.72\linewidth, height=5.2cm,
    ybar, bar width=0.9cm,
    ymin=0, ymax=55,
    ylabel={\% N by mass},
    symbolic x coords={\ce{(NH4)2SO4}, \ce{NH4NO3}, \ce{CO(NH2)2}},
    xtick=data,
    x tick label style={font=\small},
    nodes near coords, nodes near coords align={vertical},
    every node near coord/.append style={font=\small\bfseries},
    axis lines=left,
    ymajorgrids=true, grid style={popline!40},
]
\addplot[fill=popblue, draw=popdark] coordinates {(\ce{(NH4)2SO4},21.2) (\ce{NH4NO3},35) (\ce{CO(NH2)2},46.7)};
\end{axis}
\end{tikzpicture}
\end{center}
\legende{Percentage of nitrogen by mass in the three common straight nitrogen fertilizers.}
\end{popfigure}

\courssub{Eutrophication: When Fertilizers Reach the Water}

Fertilizer spread on a field does not all stay there. Nitrates, \ce{NO3-}, are extremely soluble and are not held by soil particles, so rain \emph{leaches} them through the soil into streams; phosphates are carried mainly in surface \emph{runoff} with eroded soil. The nutrients end up in rivers and lakes — exactly the wrong place.

\begin{definition}[Eutrophication]
\cle{Eutrophication} is the excessive enrichment of a body of water with plant nutrients, especially nitrates and phosphates, leading to explosive algal growth, depletion of dissolved oxygen and the death of aquatic life.
\end{definition}

The chain of events is a beautiful — and tragic — piece of applied biology and chemistry:

\begin{enumerate}
\item \textbf{Nutrient input:} nitrates and phosphates from fertilizer runoff (also from sewage and detergents) enter the water.
\item \textbf{Algal bloom:} algae, suddenly well fed, multiply explosively and cover the surface as a green carpet.
\item \textbf{Light blockage and death:} the dense surface layer blocks sunlight; algae and plants below die.
\item \textbf{Aerobic decomposition:} bacteria decompose the dead organic matter, \emph{consuming dissolved oxygen} as they respire. The water develops a high \cle{biochemical oxygen demand} (BOD).
\item \textbf{Oxygen depletion:} dissolved \ce{O2} falls so low that fish and other aquatic animals suffocate and die. Anaerobic decay then releases foul gases such as \ce{H2S} and \ce{CH4}.
\end{enumerate}

\begin{popfigure}
\begin{center}
\begin{tikzpicture}[
    box/.style={draw=popdark, fill=popblueL, rounded corners=2pt, text width=2.6cm, align=center, minimum height=1.1cm, font=\small},
    arr/.style={-{Stealth[length=2.5mm]}, thick, poporange}
]
\node[box] (a) at (0,0) {Fertilizer runoff: \ce{NO3-}, \ce{PO4^{3-}}};
\node[box, fill=popgreenL] (b) at (3.6,0) {Algal bloom on surface};
\node[box] (c) at (7.2,0) {Bacterial decay of dead algae};
\node[box] (d) at (10.8,0) {Dissolved \ce{O2} depleted (high BOD)};
\node[box, fill=poppink!25] (e) at (14.4,0) {Fish and aquatic life die};
\draw[arr] (a) -- (b);
\draw[arr] (b) -- (c);
\draw[arr] (c) -- (d);
\draw[arr] (d) -- (e);
\end{tikzpicture}

\vspace{0.4cm}
\begin{tikzpicture}
\begin{axis}[
    width=0.72\linewidth, height=4.6cm,
    xlabel={Distance downstream of fertilizer input / km},
    ylabel={Dissolved \ce{O2} / mg\,dm$^{-3}$},
    xmin=0, xmax=10, ymin=0, ymax=10,
    axis lines=left,
    ymajorgrids=true, grid style={popline!40},
    legend style={font=\small, at={(0.97,0.95)}, anchor=north east, draw=none, fill=none},
]
\addplot[very thick, popblue, smooth, domain=0:10, samples=100]
    {9 - 7.5*exp(-((x-4)^2)/3.5)};
\addlegendentry{dissolved \ce{O2} (``oxygen sag'' curve)}
\draw[dashed, poppink, thick] (axis cs:0,4) -- (axis cs:10,4);
\addlegendentry{minimum for fish survival}
\end{axis}
\end{tikzpicture}
\end{center}
\legende{The eutrophication sequence (top) and the characteristic ``oxygen sag'': dissolved \ce{O2} plunges just downstream of the nutrient input, then slowly recovers as the organic load is used up.}
\end{popfigure}

\begin{savaistu}[Eutrophication in Cameroon]
The Municipal Lake in Yaoundé has repeatedly suffered algal blooms and fish deaths linked to nutrient-rich runoff from surrounding farms, gardens and untreated waste water. Similar problems affect streams draining intensively fertilised plantations and market-gardening zones. Eutrophication is not a distant, textbook problem — it is visible in our own cities.
\end{savaistu}

\courssub{Other Environmental Effects of Fertilizers}

\textbf{Nitrate in drinking water.} Leached nitrate can contaminate wells and boreholes. In infants, gut bacteria reduce \ce{NO3-} to nitrite, \ce{NO2-}, which oxidises the iron(II) of haemoglobin to iron(III), forming methaemoglobin, which cannot carry oxygen. The result is \cle{methaemoglobinaemia}, the ``blue baby syndrome''. This is why drinking-water standards set a strict limit on nitrate concentration.

\textbf{Soil acidification.} Long-term use of ammonium-based fertilizers slowly acidifies the soil.

\begin{propriete}[Acidifying action of ammonium fertilizers]
Nitrifying bacteria in the soil oxidise ammonium ions to nitrate, releasing hydrogen ions:
\[ \ce{NH4+(aq) + 2O2(g) ->[nitrifying bacteria] NO3-(aq) + 2H+(aq) + H2O(l)} \]
The released \ce{H+} lowers the soil pH over the years, so farmers must periodically add lime (\ce{CaO} or \ce{CaCO3}) to neutralise the acidity. (The overall process is examinable as an explanation, not a proof.)
\end{propriete}

\textbf{Gaseous emissions from soils.} Fertilised soils emit nitrogen oxides. Of particular concern is \cle{dinitrogen monoxide}, \ce{N2O} (nitrogen oxidation state $+1$), released by microbial nitrification and denitrification. \ce{N2O} is a powerful \cle{greenhouse gas} — far more effective per molecule than \ce{CO2} — and it also contributes to the destruction of stratospheric ozone.

\courssub{Oxides of Nitrogen and Photochemical Smog}

Fertilizers are not the only nitrogen pollution. This completes the picture begun with the oxides of sulphur in Section 3.

\textbf{Sources of \ce{NO} and \ce{NO2} (together called \ce{NO_x}):}
\begin{itemize}
\item \textbf{Vehicle engines:} at the very high temperatures of combustion, atmospheric nitrogen and oxygen combine:
\[ \ensuremath{\mathrm{N_{2}(g) + O_{2}(g) \rightarrow [high\ T] 2NO(g)}} \qquad \ce{2NO(g) + O2(g) -> 2NO2(g)} \]
\item \textbf{Lightning:} the enormous energy of a lightning flash drives the same \ce{N2 + O2} reaction — nature's own nitrogen fixation.
\end{itemize}

\textbf{Photochemical smog.} In strong sunlight over traffic-choked cities, \ce{NO2} absorbs light and splits:
\[ \ce{NO2(g) ->[sunlight] NO(g) + O(g)} \qquad \ce{O(g) + O2(g) -> O3(g)} \]
The ozone and atomic oxygen react with unburnt hydrocarbons to give a brown, irritating haze — \cle{photochemical smog} — containing harmful \cle{tropospheric ozone} and compounds that attack eyes and lungs. Note the irony: ozone high in the stratosphere protects us, but ozone at ground level is a pollutant.

\textbf{Control: catalytic converters.} Modern exhaust systems pass the gases over a platinum/rhodium catalyst, which reduces the nitrogen oxides back to harmless nitrogen while oxidising carbon monoxide and unburnt hydrocarbons:
\[ \ce{2NO(g) + 2CO(g) ->[Pt/Rh] N2(g) + 2CO2(g)} \]

\begin{attention}[Eutrophication is not acid rain]
Do not confuse these two pollution problems. \textbf{Acid rain} is caused by \ce{SO2} and \ce{NO_x} gases dissolving in rain to form acids; it is controlled by flue-gas desulphurisation and catalytic converters. \textbf{Eutrophication} is caused by \emph{dissolved nutrients} (\ce{NO3-}, \ce{PO4^{3-}}) in water; it is controlled by careful fertilizer use and sewage treatment. Different causes, different victims, different remedies — examiners reward candidates who keep them apart.
\end{attention}

\courssub{Sustainable Use of Fertilizers}

The goal is not to abandon fertilizers — without them, Cameroon could not feed its growing population — but to use them wisely:

\begin{itemize}
\item \textbf{Controlled application rates:} apply only what the crop can absorb, guided by soil testing; excess fertilizer is wasted money and guaranteed pollution.
\item \textbf{Correct timing:} avoid spreading just before heavy rain, which washes nutrients straight into streams.
\item \textbf{Organic manure and compost:} release nutrients slowly, improve soil structure, and recycle farm waste.
\item \textbf{Legume crop rotation:} beans, groundnuts and soya host \emph{Rhizobium} bacteria in root nodules that fix atmospheric \ce{N2} naturally, reducing the need for nitrogen fertilizer in the following season.
\item \textbf{Buffer strips:} bands of unfertilised grass or trees along rivers and lakes trap runoff nutrients and sediment before they reach the water.
\end{itemize}

\courssub{Chapter Synthesis: Sulphur and Nitrogen in Perspective}

We close the chapter by gathering its two great themes.

\textbf{Oxidation-state maps.} Both elements display a wide range of oxidation states, and the chemistry of this chapter is largely the chemistry of moving between them:

\begin{popfigure}
\begin{center}
\begin{tikzpicture}[scale=1.0, font=\small]
% Sulphur line
\draw[thick, popdark, -{Stealth}] (-0.4,1.6) -- (12.6,1.6);
\node[left, popdark, font=\small\bfseries] at (-0.5,1.6) {S:};
\fill[popblue] (0.8,1.6) circle (2.2pt);
\node[above, popblue] at (0.8,1.7) {$-2$};
\node[below] at (0.8,1.45) {\ce{H2S}};
\fill[poporange] (4.0,1.6) circle (2.2pt);
\node[above, poporange] at (4.0,1.7) {$0$};
\node[below] at (4.0,1.45) {\ce{S8}};
\fill[popblue] (7.2,1.6) circle (2.2pt);
\node[above, popblue] at (7.2,1.7) {$+4$};
\node[below] at (7.2,1.45) {\ce{SO2}};
\fill[popblue] (10.4,1.6) circle (2.2pt);
\node[above, popblue] at (10.4,1.7) {$+6$};
\node[below] at (10.4,1.45) {\ce{SO3}, \ce{H2SO4}};
% Nitrogen line
\draw[thick, popdark, -{Stealth}] (-0.4,0) -- (12.6,0);
\node[left, popdark, font=\small\bfseries] at (-0.5,0) {N:};
\fill[popgreen] (0.8,0) circle (2.2pt);
\node[above, popgreen] at (0.8,0.1) {$-3$};
\node[below] at (0.8,-0.15) {\ce{NH3}};
\fill[poporange] (2.4,0) circle (2.2pt);
\node[above, poporange] at (2.4,0.1) {$0$};
\node[below] at (2.4,-0.15) {\ce{N2}};
\fill[popgreen] (4.0,0) circle (2.2pt);
\node[above, popgreen] at (4.0,0.1) {$+1$};
\node[below] at (4.0,-0.15) {\ce{N2O}};
\fill[popgreen] (5.6,0) circle (2.2pt);
\node[above, popgreen] at (5.6,0.1) {$+2$};
\node[below] at (5.6,-0.15) {\ce{NO}};
\fill[popgreen] (7.2,0) circle (2.2pt);
\node[above, popgreen] at (7.2,0.1) {$+4$};
\node[below] at (7.2,-0.15) {\ce{NO2}};
\fill[popgreen] (10.4,0) circle (2.2pt);
\node[above, popgreen] at (10.4,0.1) {$+5$};
\node[below] at (10.4,-0.15) {\ce{HNO3}};
\end{tikzpicture}
\end{center}
\legende{Oxidation-state maps for the key sulphur and nitrogen species of this chapter. Oxidation (loss of electrons) moves to the right; reduction moves to the left.}
\end{popfigure}

\textbf{The nitrogen cycle} ties fixation, life, decay and return to the atmosphere into one grand loop:

\begin{popfigure}
\begin{center}
\begin{tikzpicture}[
    box/.style={draw=popdark, fill=popblueL, rounded corners=2pt, text width=2.5cm, align=center, minimum height=0.95cm, font=\small},
    arr/.style={-{Stealth[length=2.5mm]}, thick, poppurple}
]
\node[box, fill=popgreenL] (atm) at (0,3.4) {Atmospheric \ce{N2}};
\node[box] (fix) at (-4.2,1.6) {Fixation: lightning, Haber, \emph{Rhizobium}};
\node[box] (plant) at (-4.2,-1.2) {Nitrates absorbed by plants (protein)};
\node[box] (anim) at (0,-2.6) {Animals eat plants};
\node[box] (dec) at (4.2,-1.2) {Decay \& waste: \ce{NH3}/\ce{NH4+}};
\node[box] (nit) at (4.2,1.6) {Nitrification to \ce{NO3-}};
\node[box, fill=poppink!20] (den) at (0,1.6) {Denitrifying bacteria};
\draw[arr] (atm) -- (fix);
\draw[arr] (fix) -- (plant);
\draw[arr] (plant) -- (anim);
\draw[arr] (anim) -- (dec);
\draw[arr] (dec) -- (nit);
\draw[arr] (nit) -- (den);
\draw[arr] (den) -- (atm);
\draw[arr, dashed] (nit.west) to[bend right=15] (plant.east);
\end{tikzpicture}
\end{center}
\legende{The nitrogen cycle: fixation converts \ce{N2} into usable forms; nitrification regenerates nitrate; denitrification returns \ce{N2} to the atmosphere.}
\end{popfigure}

\textbf{The three industrial processes} that anchor the chapter:

\begin{center}
\begin{tabularx}{\linewidth}{|l|X|X|X|}
\hline
\rowcolor{popblueL} \textbf{Process} & \textbf{Product} & \textbf{Key equation} & \textbf{Conditions} \\
\hline
Haber & \ce{NH3} & \ce{N2(g) + 3H2(g) <=> 2NH3(g)} & $\sim$450\,\textdegree C, $\sim$200 atm, Fe catalyst \\
\hline
Ostwald & \ce{HNO3} & \ce{4NH3(g) + 5O2(g) -> 4NO(g) + 6H2O(g)}, then \ce{NO -> NO2 -> HNO3} & $\sim$900\,\textdegree C, Pt--Rh catalyst \\
\hline
Contact & \ce{H2SO4} & \ce{2SO2(g) + O2(g) <=> 2SO3(g)}; \ce{SO3 + H2SO4 -> H2S2O7}; \ce{H2S2O7 + H2O -> 2H2SO4} & $\sim$450\,\textdegree C, \ce{V2O5} catalyst \\
\hline
\end{tabularx}
\end{center}

\begin{aretenir}[Key points of this section]
\begin{itemize}
\item Crops remove N, P and K from the soil; fertilizers replace them. N $\to$ leaves/protein, P $\to$ roots/energy, K $\to$ water balance/disease resistance.
\item $\%N = \dfrac{\text{mass of N in formula}}{M}\times 100$: \ce{NH4NO3} = 35\,\%, urea $\approx$ 46.7\,\%, \ce{(NH4)2SO4} $\approx$ 21.2\,\%.
\item NPK labels give \%N, \%\ce{P2O5} and \%\ce{K2O} \emph{equivalents}.
\item Eutrophication: runoff nutrients $\to$ algal bloom $\to$ aerobic decay $\to$ \ce{O2} depletion (high BOD) $\to$ fish death.
\item Nitrate in drinking water causes blue-baby syndrome; ammonium fertilizers acidify soil (\ce{NH4+ -> NO3- + 2H+}); soils emit the greenhouse gas \ce{N2O}.
\item \ce{NO_x} from engines and lightning cause photochemical smog and tropospheric ozone; catalytic converters reduce \ce{NO} back to \ce{N2}.
\item Sustainable practice: right rate, right timing, manure, legume rotation, buffer strips.
\end{itemize}
\end{aretenir}

\begin{exercice}[Fertilizer value and the environment]
(a) Calculate the percentage by mass of nitrogen in ammonium sulphate, \ce{(NH4)2SO4}. ($A_r$: H = 1, N = 14, O = 16, S = 32.)

(b) A farmer applies ammonium nitrate at \SI{200}{kg} per hectare. What mass of nitrogen is applied per hectare?

(c) Explain, with the aid of one equation, why years of ammonium sulphate use make the soil acidic, and name a substance used to correct this.

(d) Describe the sequence of events by which fertilizer runoff can kill fish in a lake, and state two farming practices that reduce this risk.
\end{exercice}

\begin{corrige}[Exercise 1]
\textbf{(a)} $M(\ce{(NH4)2SO4}) = 2(14) + 8(1) + 32 + 4(16) = 132\ \mathrm{g\,mol^{-1}}$; mass of N $= 2\times 14 = 28$.
\[ \%N = \frac{28}{132}\times 100 \approx 21.2\ \% \]
\textbf{(b)} \ce{NH4NO3} is 35\,\% N, so mass of N $= 0.35 \times 200 = \SI{70}{kg}$ per hectare.

\textbf{(c)} Nitrifying bacteria oxidise ammonium to nitrate, releasing hydrogen ions:
\[ \ce{NH4+(aq) + 2O2(g) -> NO3-(aq) + 2H+(aq) + H2O(l)} \]
The accumulating \ce{H+} lowers the pH. It is corrected by adding lime (\ce{CaO} or \ce{CaCO3}), a base that neutralises the acid.

\textbf{(d)} Nitrates and phosphates leach or run off into the lake; the nutrient enrichment causes an algal bloom; when the algae die, aerobic bacteria decompose them and consume dissolved oxygen (high BOD); the oxygen concentration falls below the level fish need, and they suffocate. Risk is reduced by: applying only the amount of fertilizer the crop needs; timing application to avoid heavy rain; keeping buffer strips of vegetation along waterways; using organic manure or legume rotation. (Any two.)
\end{corrige}

\newpage

\coursec{Integration activity}

\courssub{Aims of the activity}

This integration activity brings together the whole chapter in a single, realistic decision-making problem. By the end of the activity you should be able to:

\begin{itemize}
\item \cle{Calculate} the percentage by mass of nitrogen in common fertilizers and compare their \cle{cost per kilogram of nitrogen};
\item \cle{Write balanced equations} for the manufacture of urea, ammonium sulphate and an NPK fertilizer, tracing each back to the \cle{Haber process}, the \cle{Ostwald process} and the \cle{Contact process}, with their operating conditions;
\item \cle{Explain chemically} how fertilizer runoff causes \cle{eutrophication} and fish deaths, and why ammonium sulphate can \cle{acidify} a soil;
\item \cle{Construct and justify} a costed fertilization plan together with realistic pollution-control measures, and defend your recommendation orally.
\end{itemize}

The activity is designed for group work (3--4 students), followed by a 5-minute presentation and class debate. It reproduces the kind of evidence-based reasoning expected in the GCE Advanced Level structured and essay questions.

\begin{methode}[How to compare fertilizer costs]
\begin{enumerate}
\item Calculate the molar mass of the fertilizer.
\item Find the mass of nitrogen per mole (or per 100 g).
\item Compute \%N = (mass of N / molar mass) $\times$ 100.
\item Mass of N per bag = bag mass $\times$ \%N/100.
\item Cost per kg N = bag price / mass of N per bag.
\item Always compare \emph{per kg of nutrient}, never per bag.
\end{enumerate}
\end{methode}

\courssub{Situation}

\textbf{Case Study: Fertilizer Choice and Pollution Control for a Ndop Plain Rice Cooperative.}

The Ndop Plain in the North West Region of Cameroon is one of the country's major rice-producing areas. The \emph{Mboghaa Rice Farmers' Cooperative} cultivates \SI{50}{ha} of paddy rice. Yields have fallen over recent seasons, so the cooperative commissioned a soil analysis and collected market prices for three fertilizers available in Bamenda. At the same time, fishermen on a stream draining the plain reported sudden \cle{fish deaths} in August, shortly after heavy rains that followed the main fertilizer application.

\textbf{Data pack 1 -- Soil analysis (per hectare, top 20 cm).}

\begin{center}
\begin{tabularx}{\linewidth}{|l|X|X|}
\hline
\rowcolor{popblueL} \textbf{Parameter} & \textbf{Measured value} & \textbf{Agronomist's comment} \\
\hline
Available nitrogen & Low & Deficient for rice; supplement needed \\
\hline
Available phosphorus (P) & Low & Deficient; supplement needed \\
\hline
Available potassium (K) & Adequate & No supplement needed \\
\hline
Soil pH & 5.6 & Slightly acidic; avoid further acidification \\
\hline
\end{tabularx}
\end{center}

The agronomist recommends applying \SI{80}{kg} of nitrogen and \SI{40}{kg} of $\ce{P2O5}$ equivalent per hectare for the season.

\textbf{Data pack 2 -- Fertilizers on the Bamenda market (50 kg bags).}

\begin{center}
\begin{tabularx}{\linewidth}{|l|X|X|}
\hline
\rowcolor{popblueL} \textbf{Fertilizer} & \textbf{Composition} & \textbf{Price per 50 kg bag (FCFA)} \\
\hline
Urea & $\ce{CO(NH2)2}$, assumed pure & 18\,000 \\
\hline
Ammonium sulphate & $\ce{(NH4)2SO4}$, assumed pure & 15\,000 \\
\hline
NPK 20-10-10 & 20\,\% N, 10\,\% $\ce{P2O5}$, 10\,\% $\ce{K2O}$ by mass & 22\,000 \\
\hline
\end{tabularx}
\end{center}

\textbf{Data pack 3 -- Rainfall and stream report.} August rainfall was \SI{310}{mm}, with three storms in the week after top-dressing. The stream draining the plain showed a dense green algal mat, a strong smell of decay, and dissolved oxygen below \SI{2}{mg.L^{-1}} at dawn (healthy streams usually hold above \SI{6}{mg.L^{-1}}). Dead tilapia were found along the banks.

The cooperative asks your group of student chemists to prepare a \textbf{one-page advisory note}: which fertilizer(s) to buy, at what cost, and how to prevent a repeat of the fish kill.

\courssub{Tasks}

Work through the following tasks in order. Show all working; give masses to 3 significant figures.

\begin{enumerate}
\item \textbf{Fertilizer analysis.} Using relative atomic masses ($H = 1$, $C = 12$, $N = 14$, $O = 16$, $S = 32$), calculate the percentage by mass of nitrogen in (a) urea $\ce{CO(NH2)2}$ and (b) ammonium sulphate $\ce{(NH4)2SO4}$. Hence calculate, for each of the three fertilizers, the \cle{cost per kilogram of nitrogen} delivered.
\item \textbf{Industrial chemistry.} Write balanced equations, with conditions, for the manufacture of:
\begin{enumerate}
\item ammonia by the Haber process from $\ce{N2}$ and $\ce{H2}$;
\item nitric acid from ammonia by the Ostwald process (three steps);
\item sulphuric acid by the Contact process (three steps, starting from sulphur);
\item urea from ammonia and carbon dioxide;
\item ammonium sulphate from ammonia and sulphuric acid.
\end{enumerate}
State the oxidation number of nitrogen or sulphur in each reactant and product, and identify which steps are redox.
\item \textbf{The NPK question.} The NPK 20-10-10 supplies phosphorus as well as nitrogen. Given that the soil needs \SI{80}{kg} N and \SI{40}{kg} $\ce{P2O5}$ per hectare, calculate the mass of NPK 20-10-10 needed per hectare to supply all the phosphorus, and the mass of nitrogen this provides. How much additional nitrogen must come from another fertilizer?
\item \textbf{Environmental chemistry.} Using the data pack, explain chemically the sequence of events from fertilizer application to fish deaths. Your answer must mention: dissolution and runoff of nitrate and ammonium ions, \cle{eutrophication}, algal bloom, decomposition by aerobic bacteria, and the fall in dissolved oxygen. Include the equation for the oxidation of ammonium ions by nitrifying bacteria. Also explain, with an equation, why long-term use of ammonium sulphate can lower the soil pH, and why this matters for a soil already at pH 5.6.
\item \textbf{Recommendation.} Propose a costed fertilization plan for the whole \SI{50}{ha} (choice of fertilizer(s), masses, number of bags, total cost in FCFA) and \textbf{two} pollution-control measures, each justified chemically or physically. Prepare your one-page advisory note and defend it in a 5-minute class debate.
\end{enumerate}

\courssub{Model answer}

\textbf{Task 1 -- Percentage nitrogen and cost per kg N.}

\begin{exoresolu}[Urea $\ce{CO(NH2)2}$]
Molar mass $M = 12 + 16 + 2\times(14 + 2\times1) = \SI{60}{g.mol^{-1}}$.
Mass of N per mole = $2 \times 14 = \SI{28}{g}$.
\[ \%N = \frac{28}{60}\times 100 = 46.7\,\%. \]
A 50 kg bag contains $50 \times 0.467 = \SI{23.3}{kg}$ N.
Cost per kg N = $18\,000 / 23.3 = \SI{772}{FCFA.kg^{-1}}$.
\tcblower
\textbf{Answer:} 46.7\% N; \SI{772}{FCFA.kg^{-1}} N.
\end{exoresolu}

\begin{exoresolu}[Ammonium sulphate $\ce{(NH4)2SO4}$]
Molar mass $M = 2\times(14+4) + 32 + 4\times16 = \SI{132}{g.mol^{-1}}$.
Mass of N per mole = $2 \times 14 = \SI{28}{g}$.
\[ \%N = \frac{28}{132}\times 100 = 21.2\,\%. \]
A 50 kg bag contains $50 \times 0.212 = \SI{10.6}{kg}$ N.
Cost per kg N = $15\,000 / 10.6 = \SI{1\,415}{FCFA.kg^{-1}}$.
\tcblower
\textbf{Answer:} 21.2\% N; \SI{1\,415}{FCFA.kg^{-1}} N.
\end{exoresolu}

\begin{exoresolu}[NPK 20-10-10]
Grade means 20\% N by mass. A 50 kg bag contains $50 \times 0.20 = \SI{10.0}{kg}$ N.
Cost per kg N = $22\,000 / 10.0 = \SI{2\,200}{FCFA.kg^{-1}}$ N.
\tcblower
\textbf{Answer:} 20.0\% N; \SI{2\,200}{FCFA.kg^{-1}} N.
\end{exoresolu}

\begin{aretenir}[Cheapest nitrogen source]
Per kilogram of nitrogen, urea (\SI{772}{FCFA}) is by far the cheapest source, then ammonium sulphate (\SI{1\,415}{FCFA}), then NPK (\SI{2\,200}{FCFA}). However, only the NPK supplies the phosphorus the soil also needs.
\end{aretenir}

\textbf{Task 2 -- Industrial equations with oxidation states and redox identification.}

\begin{methode}[Writing industrial equations for GCE]
\begin{enumerate}
\item Balance atoms and charge; add state symbols (g), (l), (aq), (s).
\item Write conditions above the arrow: catalyst, temperature, pressure.
\item Assign oxidation numbers to the key element (N or S) in every species.
\item A step is \cle{redox} if oxidation numbers change.
\end{enumerate}
\end{methode}

\paragraph{Haber process (ammonia synthesis).}
\begin{align*}
\ensuremath{\mathrm{N_{2}(g) + 3H_{2}(g) }}&\ensuremath{\mathrm{\rightleftharpoons [Fe,\ 450\ ^{\circ}C,\ 200\ atm] 2NH_{3}(g)}} \qquad \Delta H = -\SI{92}{kJ.mol^{-1}}
\end{align*}
Oxidation numbers: $\ce{N2}$: 0; $\ce{H2}$: 0; $\ce{NH3}$: N = -3, H = +1.
\cle{Redox}: Yes (N reduced 0 $\to$ -3; H oxidised 0 $\to$ +1).

\paragraph{Ostwald process (nitric acid from ammonia).}
\begin{align*}
\text{Step 1:}\quad & \ensuremath{\mathrm{4NH_{3}(g) + 5O_{2}(g) \rightarrow [Pt/Rh,\ 900\ ^{\circ}C] 4NO(g) + 6H_{2}O(g)}} \\
\text{Step 2:}\quad & \ce{2NO(g) + O2(g) -> 2NO2(g)} \\
\text{Step 3:}\quad & \ce{3NO2(g) + H2O(l) -> 2HNO3(aq) + NO(g)}
\end{align*}
Oxidation numbers of N:
Step 1: $\ce{NH3}$ (-3) $\to$ $\ce{NO}$ (+2) \emph{oxidation}; $\ce{O2}$ (0) $\to$ $\ce{H2O}$ (-2) \emph{reduction}. \cle{Redox}.
Step 2: $\ce{NO}$ (+2) $\to$ $\ce{NO2}$ (+4) \emph{oxidation}; $\ce{O2}$ (0) $\to$ $\ce{NO2}$ (-2) \emph{reduction}. \cle{Redox}.
Step 3: $\ce{NO2}$ (+4) $\to$ $\ce{HNO3}$ (+5) \emph{oxidation} and $\to$ $\ce{NO}$ (+2) \emph{reduction} (disproportionation). \cle{Redox}.

\paragraph{Contact process (sulphuric acid from sulphur).}
\begin{align*}
\text{Step 1:}\quad & \ce{S(s) + O2(g) -> SO2(g)} \\
\text{Step 2:}\quad & \ensuremath{\mathrm{2SO_{2}(g) + O_{2}(g) \rightleftharpoons [V_{2}O_{5},\ 450\ ^{\circ}C] 2SO_{3}(g)}} \qquad \Delta H < 0 \\
\text{Step 3:}\quad & \ensuremath{\mathrm{SO_{3}(g) + H_{2}SO_{4}(l) \rightarrow  H_{2}S_{2}O_{7}(l) \xrightarrow{H2O} 2H_{2}SO_{4}(aq)}}
\end{align*}
Oxidation numbers of S:
Step 1: $\ce{S}$ (0) $\to$ $\ce{SO2}$ (+4) \emph{oxidation}; $\ce{O2}$ (0) $\to$ $\ce{SO2}$ (-2) \emph{reduction}. \cle{Redox}.
Step 2: $\ce{SO2}$ (+4) $\to$ $\ce{SO3}$ (+6) \emph{oxidation}; $\ce{O2}$ (0) $\to$ $\ce{SO3}$ (-2) \emph{reduction}. \cle{Redox}.
Step 3: S remains +6 throughout. \cle{Not redox} (absorption/hydrolysis).

\paragraph{Urea manufacture.}
\[ \ensuremath{\mathrm{2NH_{3}(g) + CO_{2}(g) \rightarrow [150-200\ ^{\circ}C,\ 150-250\ atm] CO(NH_{2})_{2}(s) + H_{2}O(l)}} \]
Oxidation numbers: N in $\ce{NH3}$ and urea = -3; C in $\ce{CO2}$ and urea = +4. \cle{Not redox}.

\paragraph{Ammonium sulphate manufacture.}
\[ \ce{2NH3(g) + H2SO4(aq) -> (NH4)2SO4(aq)} \]
Oxidation numbers: N = -3 in both; S = +6 in both. \cle{Not redox} (acid--base neutralisation).

\begin{attention}[Common errors in industrial equations]
\begin{itemize}
\item Forgetting state symbols or conditions (catalyst, T, P) -- these earn marks in GCE.
\item Writing the Ostwald absorption as $\ce{NO2 + H2O -> HNO3}$ (unbalanced). The correct stoichiometry is $\ce{3NO2 + H2O -> 2HNO3 + NO}$.
\item In the Contact process, writing $\ce{SO3 + H2O -> H2SO4}$ as the industrial step. Industrially, $\ce{SO3}$ is absorbed in concentrated $\ce{H2SO4}$ to form oleum ($\ce{H2S2O7}$), then diluted.
\item Confusing oxidation numbers: in $\ce{NH4+}$, N is -3; in $\ce{NO3-}$, N is +5.
\end{itemize}
\end{attention}

\textbf{Task 3 -- The NPK calculation.}

\begin{exoresolu}[Mass of NPK 20-10-10 required per hectare]
The soil needs \SI{40}{kg} $\ce{P2O5}$ per hectare. NPK 20-10-10 contains 10\% $\ce{P2O5}$.
\[ m(\text{NPK}) = \frac{\SI{40}{kg}}{0.10} = \SI{400}{kg.ha^{-1}}. \]
This delivers $400 \times 0.20 = \SI{80}{kg}$ N per hectare -- exactly the nitrogen requirement.
Number of 50 kg bags per hectare = $400 / 50 = 8$ bags.
\tcblower
\textbf{Answer:} \SI{400}{kg.ha^{-1}} (8 bags/ha) supplies both N and P needs exactly. No additional nitrogen needed.
\end{exoresolu}

\begin{exoresolu}[Alternative blend: NPK + urea]
If cash is limited, supply part of the P from NPK and the rest of N from cheaper urea.
Example: Use \SI{200}{kg} NPK (4 bags) $\to$ \SI{20}{kg} $\ce{P2O5}$ and \SI{40}{kg} N.
Remaining N needed = \SI{40}{kg}. Urea (46.7\% N) required = $40 / 0.467 = \SI{85.7}{kg}$ ($\approx$ 1.7 bags).
Cost per hectare: $4 \times 22\,000 + 1.7 \times 18\,000 = \SI{118\,600}{FCFA}$ (vs \SI{176\,000}{FCFA} for NPK alone).
\tcblower
\textbf{Note:} The blend is cheaper but requires handling two products and still leaves a P deficit (only 20 of 40 kg $\ce{P2O5}$ supplied). For full P requirement, NPK alone is simpler.
\end{exoresolu}

\textbf{Task 4 -- From fertilizer to fish deaths (eutrophication).}

\begin{popfigure}
\centering
\begin{tikzpicture}[
  node distance=1.8cm and 1.2cm,
  box/.style={draw=popdark, fill=popblueL, rounded corners, minimum width=2.5cm, minimum height=1cm, align=center, font=\small},
  arrow/.style={-stealth, thick, popdark}
]
\node[box, align=center] (app){Fertilizer application\\(urea, (NH$_4$)$_2$SO$_4$, NPK)};
\node[box, right=of app, align=center] (diss){Dissolution in rainwater\\(urea $\xrightarrow{\text{urease}}$ 2NH$_4^+$ + CO$_3^{2-}$)};
\node[box, right=of diss, align=center] (runoff){Runoff / leaching\\(NO$_3^-$, NH$_4^+$, PO$_4^{3-}$ enter stream)};
\node[box, below=of runoff] (eut) {Eutrophication: nutrient enrichment $\to$ algal bloom};
\node[box, right=of eut, align=center] (decomp){Algae die $\to$ aerobic decomposition\\(organic matter + O$_2$ $\xrightarrow{\text{bacteria}}$ CO$_2$ + H$_2$O)};
\node[box, right=of decomp, align=center] (nitrif){Nitrification:\\NH$_4^+$ + 2O$_2$ $\to$ NO$_3^-$ + 2H$^+$ + H$_2$O};
\node[box, below=of nitrif, align=center] (hypox){Dissolved O$_2$ falls below \SI{2}{mg.L^{-1}}\\(esp. at dawn)};
\node[box, right=of hypox] (fish) {Fish suffocation (tilapia death)};
\draw[arrow] (app) -- (diss);
\draw[arrow] (diss) -- (runoff);
\draw[arrow] (runoff) -- (eut);
\draw[arrow] (eut) -- (decomp);
\draw[arrow] (decomp) -- (hypox);
\draw[arrow] (nitrif) -- (hypox);
\draw[arrow] (hypox) -- (fish);
\end{tikzpicture}
\legende{Sequence from fertilizer application to fish kill (eutrophication).}
\end{popfigure}

\textbf{Chemical explanation:}
\begin{enumerate}
\item \textbf{Dissolution and runoff:} Urea hydrolyses in soil (enzyme urease): $\ce{CO(NH2)2 + 2H2O -> 2NH4+ + CO3^{2-}}$. Ammonium sulphate dissociates: $\ce{(NH4)2SO4 -> 2NH4+ + SO4^{2-}}$. Nitrate ($\ce{NO3-}$) from nitrification or NPK is highly soluble and not retained by soil. Heavy rains (\SI{310}{mm}) cause runoff/leaching, carrying $\ce{NO3-}$, $\ce{NH4+}$, and $\ce{PO4^{3-}}$ (from NPK) into the stream.
\item \textbf{Eutrophication and algal bloom:} The sudden influx of N and P (limiting nutrients) causes explosive growth of algae -- the observed green mat.
\item \textbf{Decomposition and oxygen depletion:} When algae die, aerobic bacteria decompose them: $\ce{C_xH_yO_z + O2 ->[bacteria] CO2 + H2O}$, consuming dissolved oxygen.
\item \textbf{Nitrification in the stream:} Nitrifying bacteria ($\ce{Nitrosomonas}$, $\ce{Nitrobacter}$) oxidise ammonium, further consuming oxygen:
\[ \ce{NH4+(aq) + 2O2(g) -> NO3-(aq) + 2H+(aq) + H2O(l)} \]
\item \textbf{Fish kill:} Dissolved oxygen drops below \SI{2}{mg.L^{-1}} at dawn (respiration continues at night, no photosynthesis). Tilapia suffocate.
\end{enumerate}

\textbf{Soil acidification by ammonium sulphate:}
The nitrification of ammonium ions in soil releases $\ce{H+}$:
\[ \ce{NH4+ + 2O2 -> NO3- + 2H+ + H2O} \]
Additionally, sulphate ($\ce{SO4^{2-}}$) is not taken up by rice roots in significant amounts, so the accompanying acidity remains in the soil. On a soil already at pH 5.6 (slightly acidic), continued use of ammonium sulphate would lower the pH further, reducing availability of P, Mo, and beneficial microbial activity, ultimately depressing rice yields. Urea or NPK (which contains nitrate and ammonium but less acidifying potential per kg N) are safer choices here.

\begin{aretenir}[Eutrophication key points]
\begin{itemize}
\item Cause: Excess N and P (nutrient enrichment) from fertilizer runoff.
\item Effect: Algal bloom $\to$ death $\to$ aerobic decomposition $\to$ oxygen depletion (hypoxia/anoxia).
\item Fish die from \emph{lack of oxygen}, not direct nitrate toxicity.
\item Nitrification ($\ce{NH4+ -> NO3-}$) consumes $\ce{O2}$ and produces $\ce{H+}$, acidifying water and soil.
\end{itemize}
\end{aretenir}

\textbf{Task 5 -- Recommendation (advisory note for the cooperative).}

\begin{aretenir}[Recommended plan: NPK 20-10-10 only]
\begin{itemize}
\item \textbf{Fertilizer:} NPK 20-10-10.
\item \textbf{Rate:} \SI{400}{kg.ha^{-1}} (8 bags/ha).
\item \textbf{Total for 50 ha:} 400 bags = \SI{20\,000}{kg}.
\item \textbf{Total cost:} $400 \times \SI{22\,000}{FCFA} = \SI{8\,800\,000}{FCFA}$.
\item \textbf{Justification:} Meets both N (\SI{80}{kg/ha}) and $\ce{P2O5}$ (\SI{40}{kg/ha}) requirements exactly with a single product; soil K is adequate so the 10\% $\ce{K2O}$ is harmless. Avoids acidifying ammonium sulphate. Simpler logistics (one product).
\end{itemize}
\end{aretenir}

\begin{aretenir}[Alternative blend (if budget forces)]
\begin{itemize}
\item \SI{200}{kg/ha} NPK (4 bags) + \SI{85.7}{kg/ha} urea (1.7 bags).
\item Supplies \SI{20}{kg} $\ce{P2O5}$ (half requirement) and \SI{80}{kg} N.
\item Cost: $\approx \SI{118\,600}{FCFA.ha^{-1}}$ $\to$ \SI{5\,930\,000}{FCFA} total.
\item \textbf{Drawback:} Only half the P requirement met; yields may still be P-limited. Two products to transport and apply.
\end{itemize}
\end{aretenir}

\begin{propriete}[Two pollution-control measures (chemically justified)]
\begin{enumerate}
\item \textbf{Split application and weather forecasting:} Apply fertilizer in three splits (basal, early tillering, panicle initiation) and \emph{never} within 48 h of forecast heavy rain. This reduces the pool of soluble $\ce{NO3-}$ and $\ce{NH4+}$ present at any one time, minimising the mass available for runoff during storms.
\item \textbf{Vegetated buffer strips and paddy bunding:} Maintain an unfertilized grass strip $\ge \SI{5}{m}$ wide along the stream, and keep paddy bunds intact. Plant roots and soil microbes in the buffer take up $\ce{NO3-}$ and $\ce{NH4+}$ (phyto-remediation) and trap sediment carrying particulate $\ce{P}$ before water reaches the stream. Bunds increase residence time, promoting denitrification ($\ce{NO3- -> N2}$) in anaerobic zones.
\end{enumerate}
\end{propriete}

\begin{attention}[Common errors in this activity]
\begin{itemize}
\item Comparing fertilizer prices \emph{per bag} instead of \cle{per kg of nutrient}.
\item Confusing NPK grade: 20-10-10 means \%N - \%\ce{P2O5} - \%\ce{K2O}, \textbf{not} elemental P and K.
\item In the Ostwald process, writing the absorption step as $\ce{NO2 + H2O -> HNO3}$ (unbalanced, misses NO recycle).
\item Stating that eutrophication kills fish by \emph{nitrate toxicity}. The cause is \emph{oxygen depletion} from bacterial decomposition and nitrification.
\item Forgetting that urea hydrolysis initially \emph{raises} pH (produces $\ce{OH-}$ via $\ce{NH3 + H2O <=> NH4+ + OH-}$), but subsequent nitrification acidifies.
\end{itemize}
\end{attention}

\begin{savaistu}[Ndop Plain and rice in Cameroon]
The Ndop Plain (North West Region) is a volcanic floodplain with heavy clay soils (vertisols) ideal for irrigated rice. It produces over 40\% of Cameroon's rice. Fertilizer use has intensified since the 2000s, but soil organic matter is declining and stream pollution is a growing concern. The \emph{Upper Noun Valley Development Authority (UNVDA)} promotes integrated soil fertility management (ISFM) combining mineral fertilizer, organic matter, and water control -- exactly the approach your advisory note should advocate.
\end{savaistu}

\newpage

\coursec{Methods and worked examples}

\courssub{Skill 1: Assigning Oxidation Numbers to Sulphur and Nitrogen}

\begin{methode}[Assigning oxidation numbers]
To determine the oxidation number (o.n.) of S or N in any species:
\begin{enumerate}
\item Recall the fixed values: \cle{oxygen} is $-2$ (except in peroxides, $-1$), \cle{hydrogen} is $+1$ (except in metal hydrides, $-1$), Group 1 metals $+1$, Group 2 metals $+2$.
\item Apply the rule: the \cle{sum of the oxidation numbers} equals the overall charge of the species (0 for a neutral molecule).
\item Set up a simple equation with $x$ = unknown oxidation number and solve.
\item Check that the answer lies within the expected range: S from $-2$ to $+6$; N from $-3$ to $+5$.
\end{enumerate}
\textbf{Reflex:} for oxoanions such as \ce{SO4^2-} or \ce{NO3-}, the sum must equal the \emph{charge of the ion}, not zero.
\end{methode}

\begin{exoresolu}[Oxidation numbers of S and N]
Determine the oxidation number of sulphur in \ce{H2S}, \ce{SO2}, \ce{SO3}, \ce{H2SO4} and \ce{S2O3^2-}, and of nitrogen in \ce{NH3}, \ce{N2}, \ce{NO}, \ce{NO2} and \ce{NO3-}.
\tcblower
Apply the sum rule in each case, taking o.n.(O) $= -2$ and o.n.(H) $= +1$.
\begin{itemize}
\item \ce{H2S}: $2(+1) + x = 0 \Rightarrow x = -2$. Sulphur is in its lowest state, $-2$.
\item \ce{SO2}: $x + 2(-2) = 0 \Rightarrow x = +4$.
\item \ce{SO3}: $x + 3(-2) = 0 \Rightarrow x = +6$, the maximum for sulphur.
\item \ce{H2SO4}: $2(+1) + x + 4(-2) = 0 \Rightarrow x = +6$.
\item \ce{S2O3^2-} (thiosulphate): $2x + 3(-2) = -2 \Rightarrow 2x = +4 \Rightarrow x = +2$ (average value).
\item \ce{NH3}: $x + 3(+1) = 0 \Rightarrow x = -3$, the minimum for nitrogen.
\item \ce{N2}: elemental form, $x = 0$.
\item \ce{NO}: $x + (-2) = 0 \Rightarrow x = +2$.
\item \ce{NO2}: $x + 2(-2) = 0 \Rightarrow x = +4$.
\item \ce{NO3-}: $x + 3(-2) = -1 \Rightarrow x = +5$, the maximum for nitrogen.
\end{itemize}
\textbf{Conclusion:} sulphur shows the states $-2, +2, +4, +6$; nitrogen shows every integer state from $-3$ to $+5$. Species in intermediate states (\ce{SO2}, \ce{NO}) can be both oxidised and reduced: they can act as reducing agents \emph{or} as oxidising agents depending on the partner.
\end{exoresolu}

\begin{attention}[Frequent traps with oxidation numbers]
\begin{itemize}
\item In \ce{S2O3^2-} the value $+2$ is an \emph{average}: the two sulphur atoms are not equivalent (one is in state $+5$-like, the other $-1$-like in the structural picture), but at GCE level the average value obtained from the sum rule is what is required.
\item Do not forget that the sum equals the \emph{charge} for ions: writing $x + 3(-2) = 0$ for \ce{NO3-} instead of $= -1$ is the commonest error.
\item An element in its \emph{maximum} state (\ce{SO3}, \ce{H2SO4}, \ce{HNO3}) can only be \emph{reduced}: such species are oxidising agents only.
\end{itemize}
\end{attention}

\courssub{Skill 2: Writing and Balancing Redox Half-Equations}

\begin{methode}[Balancing a redox equation by half-equations]
\begin{enumerate}
\item Identify the species oxidised (o.n. \emph{increases}) and the species reduced (o.n. \emph{decreases}).
\item Balance the atoms of the element changing oxidation state.
\item Balance O by adding \ce{H2O}, then H by adding \ce{H+} (acidic medium).
\item Balance the charge by adding electrons \ce{e-}.
\item Multiply the half-equations so that the electrons cancel, then add and simplify.
\item \textbf{Check:} atoms \emph{and} total charge must balance on both sides.
\end{enumerate}
\textbf{Reflex:} \ce{SO2} is usually a \emph{reducing agent} (S: $+4 \to +6$); concentrated \ce{H2SO4} and \ce{HNO3} are \emph{oxidising agents}.
\end{methode}

\begin{exoresolu}[Oxidation of sulphur dioxide by acidified dichromate]
In the laboratory, \ce{SO2(g)} decolourises acidified potassium dichromate solution, which turns from orange to green. Write the two half-equations and the overall ionic equation. (\ce{Cr2O7^2-} is reduced to \ce{Cr^3+}; \ce{SO2} is oxidised to \ce{SO4^2-}.)
\tcblower
\textbf{Reduction half-equation} (dichromate to chromium(III), Cr: $+6 \to +3$, gain of $3e^-$ per Cr, i.e. $6e^-$ for \ce{Cr2O7^2-}):
\[\ce{Cr2O7^2- + 14H+ + 6e- -> 2Cr^3+ + 7H2O}\]
Check: Cr balanced (2 each side); O balanced (7 in \ce{Cr2O7^2-}, 7 in \ce{H2O}); charge: left $-2 + 14 - 6 = +6$; right $2 \times (+3) = +6$. \checkmark

\textbf{Oxidation half-equation} (sulphur dioxide to sulphate, S: $+4 \to +6$, loss of $2e^-$):
\[\ce{SO2 + 2H2O -> SO4^2- + 4H+ + 2e-}\]
Check: S balanced; O: $2 + 2 = 4$ each side; H: 4 each side; charge: left $0$; right $-2 + 4 - 2 = 0$. \checkmark

\textbf{Combine:} multiply the oxidation half-equation by 3 so that both involve $6e^-$:
\[\ce{3SO2 + 6H2O -> 3SO4^2- + 12H+ + 6e-}\]
Add the two half-equations and cancel \ce{e-}, then simplify \ce{H+} ($14 - 12 = 2$ on the left) and \ce{H2O} ($6 - 7$: one water remains on the right):
\[\ce{Cr2O7^2- + 3SO2 + 2H+ -> 2Cr^3+ + 3SO4^2- + H2O}\]
Final check: atoms balanced (Cr 2, S 3, O 11, H 2 each side); charge: left $-2 + 0 + 2 = 0$; right $2(+3) + 3(-2) + 0 = 0$. \checkmark

The colour change orange $\to$ green confirms that \ce{SO2} acted as a \cle{reducing agent}. This reaction is also a standard \emph{test} for \ce{SO2} in qualitative analysis.
\end{exoresolu}

\courssub{Skill 3: Calculations on the Contact Process}

\begin{methode}[Yield and mass calculations in the Contact process]
\begin{enumerate}
\item Write the three key steps: burning of sulphur $\ce{S(s) + O2(g) -> SO2(g)}$; catalytic oxidation $\ce{2SO2(g) + O2(g) <=> 2SO3(g)}$ ($\ce{V2O5}$ catalyst, about $450\ ^\circ\mathrm{C}$, 1--2 atm); absorption of \ce{SO3} in concentrated acid $\ce{SO3 + H2SO4 -> H2S2O7}$ (oleum), then dilution $\ce{H2S2O7 + H2O -> 2H2SO4}$.
\item Overall: 1 mol S gives 1 mol \ce{H2SO4} (sulphur is conserved through the chain).
\item Use $n = \dfrac{m}{M}$, scale by the mole ratio, then apply the percentage yield:
\[\text{actual mass} = \text{theoretical mass} \times \frac{\%\text{ yield}}{100}\]
\item Quote the conditions and the \emph{reason} for them (compromise between rate and equilibrium position).
\end{enumerate}
\end{methode}

\begin{attention}[Why \ce{SO3} is not dissolved directly in water]
\ce{SO3} reacts with water extremely violently, producing a corrosive \emph{mist} of \ce{H2SO4} droplets that is very difficult to condense. Industrially, \ce{SO3} is therefore absorbed in \emph{concentrated} sulphuric acid to form oleum \ce{H2S2O7}, which is then diluted in a controlled way. Quoting this detail earns marks in ``describe the Contact process'' questions.
\end{attention}

\begin{exoresolu}[Mass of sulphuric acid from sulphur]
A plant burns 64 tonnes of sulphur. Assuming the conversion to \ce{SO3} is 98\% efficient and absorption is complete, calculate the mass of pure sulphuric acid produced. (S = 32, O = 16, H = 1.)
\tcblower
\textbf{Step 1 — Moles of sulphur:}
\[n(\ce{S}) = \frac{64 \times 10^{6}\ \mathrm{g}}{32\ \mathrm{g\,mol^{-1}}} = 2.0 \times 10^{6}\ \mathrm{mol}\]
\textbf{Step 2 — Mole ratio:} from the chain $\ce{S -> SO2 -> SO3 -> H2SO4}$, one mole of S gives one mole of \ce{H2SO4}:
\[n(\ce{H2SO4})_{\text{theoretical}} = 2.0 \times 10^{6}\ \mathrm{mol}\]
\textbf{Step 3 — Apply the 98\% conversion:}
\[n(\ce{H2SO4})_{\text{actual}} = 0.98 \times 2.0 \times 10^{6} = 1.96 \times 10^{6}\ \mathrm{mol}\]
\textbf{Step 4 — Convert to mass:} $M(\ce{H2SO4}) = 2(1) + 32 + 4(16) = 98\ \mathrm{g\,mol^{-1}}$
\[m = 1.96 \times 10^{6} \times 98 = 1.9208 \times 10^{8}\ \mathrm{g} \approx 192\ \text{tonnes}\]
\textbf{Answer:} about $\boxed{192\ \text{tonnes}}$ of pure \ce{H2SO4}.

\textbf{Remark:} the oxidation of \ce{SO2} is exothermic; a moderate temperature ($\sim 450\ ^\circ$C) is a \cle{compromise}: lower temperatures favour the equilibrium yield but the reaction would be too slow even with the \ce{V2O5} catalyst. The heat released is recovered to maintain the catalyst beds at working temperature.
\end{exoresolu}

\courssub{Skill 4: Explaining Acid Rain: Equations and pH}

\begin{methode}[Writing the chemistry of sulphur pollution]
\begin{enumerate}
\item Identify the source: burning of fossil fuels containing sulphur impurities, roasting of sulphide ores, volcanic emissions.
\item Write the sequence of atmospheric reactions:
\[\ce{S + O2 -> SO2} \qquad \ce{2SO2 + O2 -> 2SO3}\ (\text{slow, catalysed by dust and NO}_x)\]
\[\ce{SO3 + H2O -> H2SO4} \qquad \ce{SO2 + H2O <=> H2SO3}\]
\item Conclude: rainwater contains \ce{H2SO4} and \ce{H2SO3}, so its pH falls below 5 (unpolluted rain has pH $\approx 5.6$ because of dissolved \ce{CO2} forming weak carbonic acid).
\item Link to effects: corrosion of limestone/marble ($\ce{CaCO3 + H2SO4 -> CaSO4 + H2O + CO2}$), acidification of lakes and soils, leaching of toxic \ce{Al^3+} ions, damage to leaves and reduced crop yields.
\end{enumerate}
\textbf{Reflex:} always mention both the \emph{dry} pathway (direct deposition of \ce{SO2} gas on surfaces and lungs) and the \emph{wet} pathway (acid rain).
\end{methode}

\begin{exoresolu}[Acid rain attacking a limestone building]
A limestone (\ce{CaCO3}) statue in an industrial town loses 0.50 mol of \ce{CaCO3} per year through reaction with acid rain. (a) Write the equation for the reaction with sulphuric acid. (b) Calculate the mass of limestone eroded per year. (c) State two other environmental effects of \ce{SO2} pollution. (Ca = 40, C = 12, O = 16.)
\tcblower
\textbf{(a)} Sulphuric acid reacts with calcium carbonate:
\[\ce{CaCO3(s) + H2SO4(aq) -> CaSO4(s) + H2O(l) + CO2(g)}\]
\textbf{(b)} $M(\ce{CaCO3}) = 40 + 12 + 3(16) = 100\ \mathrm{g\,mol^{-1}}$
\[m = n \times M = 0.50 \times 100 = 50\ \mathrm{g}\]
The statue loses $\boxed{50\ \mathrm{g}}$ of limestone per year.
\textbf{(c)} Two further effects:
\begin{itemize}
\item \cle{Acidification of lakes and rivers}: fish eggs fail to hatch and aquatic life dies when the pH falls below about 5.
\item \cle{Damage to vegetation}: \ce{SO2} destroys chlorophyll, yellows leaves and reduces crop yields; soils are acidified and aluminium ions toxic to roots are released.
\end{itemize}
\textbf{Control measures (to quote):} flue-gas desulphurisation with limestone ($\ce{CaCO3 + SO2 -> CaSO3 + CO2}$), use of low-sulphur fuels, and catalytic converters on vehicles (which also reduce the \ce{NO_x} that catalyses \ce{SO2} oxidation).
\end{exoresolu}

\courssub{Skill 5: Calculations on Ammonia and the Haber Process}

\begin{methode}[Equilibrium/yield reasoning for the Haber process]
For $\ce{N2(g) + 3H2(g) <=> 2NH3(g)}$, $\Delta H < 0$:
\begin{enumerate}
\item Predict the effect of a change using \cle{Le Chatelier's principle}: high pressure favours the side with fewer gas moles (4 $\to$ 2); low temperature favours the exothermic (forward) direction.
\item Explain the industrial compromise: $\sim 200$ atm, $\sim 450\ ^\circ$C, iron catalyst (to increase the rate), with continuous recycling of unreacted \ce{N2} and \ce{H2}.
\item For calculations: use the mole ratio from the balanced equation, $n = m/M$ or $n = V/V_m$ for gases, then apply the percentage conversion.
\end{enumerate}
\end{methode}

\begin{exoresolu}[Ammonia production from hydrogen]
In a Haber plant, 600 kg of hydrogen react with excess nitrogen and the measured yield of ammonia is 30\%. Calculate the mass of ammonia obtained. (N = 14, H = 1.)
\tcblower
\textbf{Step 1 — Moles of hydrogen:}
\[n(\ce{H2}) = \frac{600 \times 10^{3}\ \mathrm{g}}{2\ \mathrm{g\,mol^{-1}}} = 3.0 \times 10^{5}\ \mathrm{mol}\]
\textbf{Step 2 — Mole ratio:} from $\ce{N2 + 3H2 -> 2NH3}$, 3 mol \ce{H2} give 2 mol \ce{NH3}:
\[n(\ce{NH3})_{\text{theoretical}} = \frac{2}{3} \times 3.0 \times 10^{5} = 2.0 \times 10^{5}\ \mathrm{mol}\]
\textbf{Step 3 — Apply the 30\% yield:}
\[n(\ce{NH3})_{\text{actual}} = 0.30 \times 2.0 \times 10^{5} = 6.0 \times 10^{4}\ \mathrm{mol}\]
\textbf{Step 4 — Mass:} $M(\ce{NH3}) = 14 + 3(1) = 17\ \mathrm{g\,mol^{-1}}$
\[m = 6.0 \times 10^{4} \times 17 = 1.02 \times 10^{6}\ \mathrm{g} \approx 1.0\ \text{tonne}\]
\textbf{Answer:} $\boxed{m(\ce{NH3}) \approx 1.0\ \text{tonne}}$.

\textbf{Comment:} the low single-pass yield is acceptable industrially because unreacted \ce{N2} and \ce{H2} are \cle{recycled} continuously after the \ce{NH3} is liquefied out, so the overall conversion approaches completion.
\end{exoresolu}

\courssub{Skill 6: Reactions of Nitric Acid as an Oxidising Agent}

\begin{methode}[Predicting products of \ce{HNO3} with metals]
\begin{enumerate}
\item Concentrated \ce{HNO3} with a metal (e.g. Cu) gives the metal nitrate, \cle{brown \ce{NO2}} gas and water.
\item Dilute \ce{HNO3} gives the metal nitrate, \cle{colourless \ce{NO}} (which browns in air: $\ce{2NO + O2 -> 2NO2}$) and water.
\item Balance using half-equations: $\ce{NO3- + 2H+ + e- -> NO2 + H2O}$ (conc.) or $\ce{NO3- + 4H+ + 3e- -> NO + 2H2O}$ (dil.).
\item Remember: \ce{HNO3} never gives \ce{H2} with metals (it is an oxidising acid); N goes from $+5$ to $+4$ (\ce{NO2}) or to $+2$ (\ce{NO}).
\end{enumerate}
\end{methode}

\begin{exoresolu}[Copper and concentrated nitric acid]
A piece of copper foil is dropped into concentrated nitric acid. A brown gas is evolved and the solution turns blue-green. (a) Identify the gas and the blue-green species. (b) Write the balanced equation. (c) Calculate the volume of gas, measured at s.t.p., produced by 3.175 g of copper. (Cu = 63.5; molar volume at s.t.p. = $22.4\ \mathrm{dm^3\,mol^{-1}}$.)
\tcblower
\textbf{(a)} The brown gas is nitrogen dioxide \ce{NO2}; the blue-green solution contains copper(II) ions \ce{Cu^2+} (as \ce{Cu(NO3)2}).
\textbf{(b)} Half-equations:
\[\ce{Cu -> Cu^2+ + 2e-} \qquad \ce{NO3- + 2H+ + e- -> NO2 + H2O}\]
Double the second half-equation so that the electrons cancel, then add:
\[\ce{Cu(s) + 4HNO3(aq) -> Cu(NO3)2(aq) + 2NO2(g) + 2H2O(l)}\]
\textbf{(c)} Moles of copper:
\[n(\ce{Cu}) = \frac{3.175}{63.5} = 0.0500\ \mathrm{mol}\]
From the equation, $n(\ce{NO2}) = 2\,n(\ce{Cu}) = 0.100\ \mathrm{mol}$.
\[V = n \times V_m = 0.100 \times 22.4 = 2.24\ \mathrm{dm^3}\]
\textbf{Answer:} $\boxed{V(\ce{NO2}) = 2.24\ \mathrm{dm^3}}$ at s.t.p.

\textbf{Safety note:} \ce{NO2} is toxic and contributes to photochemical smog and acid rain ($\ce{4NO2 + 2H2O + O2 -> 4HNO3}$); the experiment must be carried out in a fume cupboard.
\end{exoresolu}

\courssub{Skill 7: Percentage of Nutrient in a Fertiliser}

\begin{methode}[Computing \%N, \%P2O5 or \%K2O in a fertiliser]
\begin{enumerate}
\item Write the formula and compute its molar mass $M$.
\item Count the mass of the nutrient element in \emph{one mole} of the compound.
\item Apply: $\%\text{nutrient} = \dfrac{\text{mass of nutrient in 1 mol}}{M} \times 100$.
\item For mixtures (NPK fertilisers), multiply by the mass fraction of each component.
\end{enumerate}
\textbf{Reflex:} the \emph{quality} of a nitrogen fertiliser increases with its \%N: compare \ce{NH4NO3} (35\%), \ce{(NH4)2SO4} (21\%), \ce{CO(NH2)2} (urea, about 47\%).
\end{methode}

\begin{exoresolu}[Comparing two nitrogen fertilisers]
A farmer near Bafoussam can buy either ammonium nitrate \ce{NH4NO3} or ammonium sulphate \ce{(NH4)2SO4} at the same price per bag. (a) Calculate the percentage by mass of nitrogen in each fertiliser. (b) Advise the farmer, stating one environmental risk of excessive fertiliser use. (N = 14, H = 1, O = 16, S = 32.)
\tcblower
\textbf{(a) Ammonium nitrate:} $M(\ce{NH4NO3}) = 14 + 4(1) + 14 + 3(16) = 80\ \mathrm{g\,mol^{-1}}$. It contains 2 N atoms, i.e. 28 g of N per mole:
\[\%N = \frac{28}{80} \times 100 = 35\%\]
\textbf{Ammonium sulphate:} $M(\ce{(NH4)2SO4}) = 2(14 + 4) + 32 + 4(16) = 132\ \mathrm{g\,mol^{-1}}$. It contains 2 N atoms (28 g):
\[\%N = \frac{28}{132} \times 100 \approx 21.2\%\]
\textbf{(b)} At equal price, \cle{ammonium nitrate} delivers more nitrogen per bag (35\% against 21.2\%), so it is the better buy for nitrogen supply. (Ammonium sulphate may still be preferred where the soil also lacks sulphur.)

\textbf{Environmental risk:} excess nitrate is washed by rain into rivers and lakes (\cle{leaching}), causing \cle{eutrophication}: algal blooms develop, and when the algae die and decay the bacteria consume the dissolved oxygen, fish die, and drinking water may exceed safe nitrate limits. Fertilisers should therefore be applied in measured doses, away from streams, and not just before heavy rain.
\end{exoresolu}

\courssub{Skill 8: Back-Titration of Ammonium in a Fertiliser}

\begin{methode}[Back-titration of an ammonium salt]
\begin{enumerate}
\item Boil a known mass of fertiliser with a known \emph{excess} of \ce{NaOH}: $\ce{NH4+ + OH- -> NH3 + H2O}$ (all the \ce{NH3} is driven off).
\item Titrate the \emph{remaining} \ce{OH-} with standard \ce{HCl}: $\ce{H+ + OH- -> H2O}$.
\item Compute: $n(\ce{NH4+}) = n(\ce{NaOH})_{\text{added}} - n(\ce{NaOH})_{\text{excess}}$, where the excess is found from the titre (do not forget the dilution factor if the solution was made up to a volumetric flask and only a portion titrated).
\item Convert to moles of salt (mind the stoichiometry: 2 \ce{NH4+} per \ce{(NH4)2SO4}), then to mass and to percentage purity.
\end{enumerate}
\textbf{Reflex:} never titrate \ce{NH4+} directly with an indicator — it is too weak an acid; the back-titration via \ce{NH3} release is the standard GCE method.
\end{methode}

\begin{exoresolu}[Purity of an ammonium sulphate fertiliser]
1.32 g of an ammonium sulphate fertiliser was boiled with $50.0\ \mathrm{cm^3}$ of $1.00\ \mathrm{mol\,dm^{-3}}$ \ce{NaOH} (excess). After cooling, the mixture was made up to $250\ \mathrm{cm^3}$ in a volumetric flask. $25.0\ \mathrm{cm^3}$ portions of this solution required $31.00\ \mathrm{cm^3}$ of $0.100\ \mathrm{mol\,dm^{-3}}$ \ce{HCl} for neutralisation. Calculate the percentage by mass of \ce{(NH4)2SO4} in the fertiliser. (N = 14, H = 1, O = 16, S = 32.)
\tcblower
\textbf{Step 1 — Moles of NaOH added:}
\[n(\ce{NaOH})_{\text{added}} = 1.00 \times \frac{50.0}{1000} = 0.0500\ \mathrm{mol}\]
\textbf{Step 2 — Moles of excess NaOH in one 25.0 cm$^3$ portion:} since $\ce{HCl + NaOH -> NaCl + H2O}$ is 1:1,
\[n(\ce{NaOH})_{\text{excess, portion}} = n(\ce{HCl}) = 0.100 \times \frac{31.00}{1000} = 3.100 \times 10^{-3}\ \mathrm{mol}\]
\textbf{Step 3 — Scale up to the full 250 cm$^3$} (factor $\dfrac{250}{25.0} = 10$):
\[n(\ce{NaOH})_{\text{excess, total}} = 3.100 \times 10^{-3} \times 10 = 0.03100\ \mathrm{mol}\]
\textbf{Step 4 — NaOH that reacted with \ce{NH4+}:}
\[n(\ce{NaOH})_{\text{reacted}} = 0.0500 - 0.03100 = 0.01900\ \mathrm{mol}\]
From $\ce{NH4+ + OH- -> NH3 + H2O}$ (1:1), $n(\ce{NH4+}) = 0.01900$ mol.
\textbf{Step 5 — Moles and mass of \ce{(NH4)2SO4}:} each formula unit contains 2 \ce{NH4+}:
\[n(\ce{(NH4)2SO4}) = \frac{0.01900}{2} = 9.500 \times 10^{-3}\ \mathrm{mol}\]
\[M(\ce{(NH4)2SO4}) = 132\ \mathrm{g\,mol^{-1}}; \qquad m = 9.500 \times 10^{-3} \times 132 = 1.254\ \mathrm{g}\]
\textbf{Step 6 — Percentage purity:}
\[\%\,\ce{(NH4)2SO4} = \frac{1.254}{1.32} \times 100 = 95.0\%\]
\textbf{Answer:} the fertiliser is $\boxed{95.0\%}$ ammonium sulphate by mass.

\textbf{Exam lesson:} a computed percentage above 100\% signals an error. The two commonest mistakes in this classic experiment are forgetting the $\times 10$ dilution factor from $25.0$ to $250\ \mathrm{cm^3}$, and forgetting the factor 2 between \ce{NH4+} and \ce{(NH4)2SO4}. Always finish with a \emph{sanity check}: the mass of pure salt found must not exceed the mass of sample taken.
\end{exoresolu}

\newpage

\coursec{Exercises}

\courssub{I check my knowledge}

\begin{exercice}[Multiple Choice Questions]
\textbf{1.} In which of the following compounds does sulphur have an oxidation state of $+4$?
\begin{enumerate}[label=\Alph*.]
\item \ce{H2S}
\item \ce{SO2}
\item \ce{SO3}
\item \ce{H2SO4}
\end{enumerate}

\textbf{2.} The catalyst used in the Contact process for the manufacture of sulphuric acid is:
\begin{enumerate}[label=\Alph*.]
\item Iron
\item Vanadium(V) oxide
\item Platinum
\item Nickel
\end{enumerate}

\textbf{3.} Which nitrogen oxide is a brown gas at room temperature?
\begin{enumerate}[label=\Alph*.]
\item \ce{N2O}
\item \ce{NO}
\item \ce{NO2}
\item \ce{N2O4}
\end{enumerate}

\textbf{4.} The Haber process for ammonia synthesis is carried out at high pressure because:
\begin{enumerate}[label=\Alph*.]
\item It increases the rate of reaction only.
\item It decreases the activation energy.
\item The forward reaction proceeds with a decrease in the number of moles of gas.
\item The reaction is endothermic.
\end{enumerate}
\end{exercice}

\begin{exercice}[True or False with Justification]
State whether each statement is True or False. Justify your answer with a chemical equation or a brief explanation.
\begin{enumerate}
\item Sulphur dioxide (\ce{SO2}) is a basic oxide.
\item Concentrated sulphuric acid acts as a dehydrating agent on sugar (\ce{C12H22O11}) leaving a black mass of carbon.
\item Ammonia gas can be dried by passing it through concentrated sulphuric acid.
\item Eutrophication is caused by the depletion of oxygen in water bodies due to excessive growth of algae fed by nitrates and phosphates.
\end{enumerate}
\end{exercice}

\begin{exercice}[Definitions and Key Terms]
Define the following terms precisely:
\begin{enumerate}
\item \textbf{Acid rain} (state the typical pH range and main acidic components).
\item \textbf{Catalyst} (mention its effect on activation energy and equilibrium position).
\item \textbf{Percentage yield} (give the formula).
\item \textbf{Le Chatelier's Principle}.
\end{enumerate}
\end{exercice}

\begin{exercice}[Direct Calculations]
\begin{enumerate}
\item Calculate the percentage by mass of nitrogen in urea, \ce{CO(NH2)2}. (Atomic masses: C=12, O=16, N=14, H=1).
\item 64 tonnes of sulphur are burned in air and the \ce{SO2} converted to \ce{H2SO4} via the Contact process with an overall yield of 95\%. Calculate the mass of sulphuric acid produced. (Molar masses: S=32, \ce{H2SO4}=98).
\item Write the balanced chemical equation, with state symbols, for the reaction between ammonia and nitric acid to form ammonium nitrate.
\item Assign the oxidation number of nitrogen in \ce{NH4+}, \ce{N2O}, \ce{NO2-}, and \ce{NO3-}.
\end{enumerate}
\end{exercice}

\courssub{I practise}

\begin{exercice}[The Contact Process -- Industrial Manufacture of \ce{H2SO4}]
\begin{enumerate}
\item Write the balanced equation for the combustion of sulphur in air.
\item Write the balanced equation for the catalytic oxidation of \ce{SO2} to \ce{SO3}. State the catalyst used, the optimum temperature, and the pressure used industrially.
\item Explain, using Le Chatelier's principle, why:
      \begin{enumerate}[label=(\alph*)]
      \item A temperature of about \SI{450}{\ensuremath{{}^{\circ}\mathrm{C}}} is used (compromise).
      \item The pressure is kept near atmospheric (not very high).
      \item Excess air (oxygen) is used.
      \end{enumerate}
\item \ce{SO3} is not dissolved directly in water. What is it dissolved in? Write the equation for the subsequent reaction with water to form \ce{H2SO4}.
\end{enumerate}
\end{exercice}

\begin{exercice}[Redox Chemistry of Sulphur and Nitrogen]
\begin{enumerate}
\item Assign the oxidation state of sulphur in: \ce{H2S}, \ce{S8}, \ce{SO2}, \ce{SO3}, \ce{H2SO4}, \ce{Na2S2O3}.
\item Assign the oxidation state of nitrogen in: \ce{NH3}, \ce{N2O}, \ce{NO}, \ce{NO2}, \ce{HNO3}, \ce{NH4NO3} (for both N atoms).
\item In the reaction: \ce{2H2S(g) + SO2(g) -> 3S(s) + 2H2O(l)}, identify the oxidising agent and the reducing agent. Write the two half-equations.
\item \ce{SO2} can act as both an oxidising and a reducing agent. Give one example of a reaction where it acts as a reducing agent and one where it acts as an oxidising agent (with balanced equations).
\end{enumerate}
\end{exercice}

\begin{exercice}[The Haber Process -- Ammonia Synthesis]
The reaction is: \ce{N2(g) + 3H2(g) <=> 2NH3(g)} \quad $\Delta H = -\SI{92}{\kilo\joule\per\mole}$.
\begin{enumerate}
\item Predict and explain the effect of \textbf{increasing the pressure} on: (a) the yield of ammonia, (b) the rate of reaction.
\item Predict and explain the effect of \textbf{increasing the temperature} on: (a) the yield of ammonia, (b) the rate of reaction.
\item State the role of the \textbf{iron catalyst}. Does it affect the yield?
\item \textbf{Calculation:} \SI{28}{\gram} of nitrogen gas reacts with excess hydrogen. If the conversion of nitrogen is 30\%, calculate the volume of ammonia gas produced at s.t.p. (Molar volume at s.t.p. = \SI{22.4}{\liter\per\mole}).
\end{enumerate}
\end{exercice}

\begin{exercice}[Fertilizers -- Composition and Economics]
\begin{enumerate}
\item Calculate the percentage by mass of nitrogen in:
      \begin{enumerate}[label=(\alph*)]
      \item Ammonium nitrate, \ce{NH4NO3}
      \item Ammonium sulphate, \ce{(NH4)2SO4}
      \item Urea, \ce{CO(NH2)2}
      \end{enumerate}
\item A farmer needs to supply \SI{46}{\kilo\gram} of nitrogen per hectare. Calculate the mass of each fertilizer required to supply this amount of nitrogen.
\item If the market prices per 50 kg bag are: \ce{NH4NO3} = 18,000 FCFA, \ce{(NH4)2SO4} = 15,000 FCFA, Urea = 20,000 FCFA, determine which fertilizer is the most economical per kilogram of nitrogen supplied. Show your working.
\item State two environmental problems associated with the excessive use of nitrogenous fertilizers.
\end{enumerate}
\end{exercice}

\courssub{I prepare for the GCE}

\begin{exercice}[Environmental Chemistry: Sulphur Oxides and Acid Rain -- 15 marks]
\begin{enumerate}
\item A roasting plant processing zinc sulphide ore releases sulphur dioxide into the atmosphere.
      \begin{enumerate}[label=(\alph*)]
      \item Write the equation for the roasting of zinc sulphide. [2 marks]
      \item Describe, with balanced equations, how \ce{SO2} in the atmosphere leads to the formation of acid rain. [4 marks]
      \item State three harmful effects of acid rain on the environment. [3 marks]
      \end{enumerate}
\item Limestone (\ce{CaCO3}) slurry is used in flue gas desulphurisation (scrubbing) to remove \ce{SO2}.
      \begin{enumerate}[label=(\alph*)]
      \item Write the balanced equation for the reaction between \ce{SO2} and \ce{CaCO3} in the presence of water/oxygen to form gypsum (\ce{CaSO4.2H2O}). [3 marks]
      \item Explain why this method is preferred over using sodium hydroxide solution. [2 marks]
      \end{enumerate}
\item State one major use of sulphuric acid in Cameroon's industry. [1 mark]
\end{enumerate}
\end{exercice}

\begin{exercice}[Laboratory Preparation and Properties of Ammonia -- 12 marks]
\begin{enumerate}
\item Describe the laboratory preparation of a dry sample of ammonia gas. Your answer must include:
      \begin{enumerate}[label=(\alph*)]
      \item The reagents used (solid). [1 mark]
      \item A labelled diagram of the apparatus setup (describe in words if drawing not possible). [3 marks]
      \item The drying agent used and why it is suitable. [2 marks]
      \item The method of collection and the reason for this choice. [2 marks]
      \end{enumerate}
\item Give one chemical test for ammonia gas and the expected observation. [1 mark]
\item Explain why concentrated sulphuric acid \textbf{cannot} be used to dry ammonia gas. Write an equation for the reaction that would occur. [2 marks]
\item Write the ionic equation for the reaction of aqueous ammonia with \ce{Cu^{2+}(aq)} to form the deep blue complex. [1 mark]
\end{enumerate}
\end{exercice}

\begin{exercice}[Data Analysis: Eutrophication and Oxygen Sag Curve -- 13 marks]
A stream flows past a fertilized farmland. Water samples were taken at five points (A, B, C, D, E) downstream from the farm entry point. The concentrations of nitrate ions (\ce{NO3-}) and dissolved oxygen (DO) were measured.

\begin{center}
\begin{tabular}{|c|c|c|}
\hline
\textbf{Sampling Point} & \textbf{\ce{[NO3-]} (\SI{}{\milli\gram\per\liter})} & \textbf{DO (\SI{}{\milli\gram\per\liter})} \\
\hline
A (Upstream) & 2 & 9.0 \\
B & 15 & 7.5 \\
C & 45 & 3.0 \\
D & 30 & 5.5 \\
E (Far downstream) & 10 & 8.0 \\
\hline
\end{tabular}
\end{center}

\begin{enumerate}
\item On graph paper (or describe the axes and plots), plot the concentration of \ce{NO3-} and DO against the sampling points (A to E) on the same axes (use two y-axes or separate graphs). [3 marks]
\item Identify the point where pollution from the farm most likely enters the stream. Justify your answer using the data. [2 marks]
\item Explain the trend in dissolved oxygen concentration from point A to point C in terms of eutrophication. Include the role of bacteria and the biochemical oxygen demand (BOD). [4 marks]
\item Explain the recovery of dissolved oxygen levels from point C to point E. [2 marks]
\item State two measures a farmer in the Cameroon highlands (e.g., around Mount Bamboutos) could take to reduce nitrate leaching into streams. [2 marks]
\end{enumerate}
\end{exercice}

\newpage

\coursec{Solutions to the exercises}

\begin{corrige}[Exercise 1 --- Multiple Choice Questions]
\textbf{1. Answer: B (\ce{SO2})}

Oxygen has oxidation state $-2$. Let $x$ be the oxidation state of S:
\begin{itemize}
\item \ce{H2S}: $x + 2(+1) = 0 \Rightarrow x = -2$
\item \ce{SO2}: $x + 2(-2) = 0 \Rightarrow x = +4$ \checkmark
\item \ce{SO3}: $x + 3(-2) = 0 \Rightarrow x = +6$
\item \ce{H2SO4}: $2(+1) + x + 4(-2) = 0 \Rightarrow x = +6$
\end{itemize}

\textbf{2. Answer: B (Vanadium(V) oxide, \ce{V2O5})}

In the Contact process, \ce{V2O5} catalyses \ce{2SO2(g) + O2(g) <=> 2SO3(g)}. Platinum works but is far more expensive and is easily poisoned by impurities such as arsenic compounds; iron and nickel are used in other processes (Haber, hydrogenation).

\textbf{3. Answer: C (\ce{NO2})}

\ce{NO2} is a pungent brown gas. \ce{N2O} and \ce{NO} are colourless; \ce{N2O4} is colourless (the brown colour of the equilibrium mixture \ce{2NO2 <=> N2O4} comes from \ce{NO2}).

\textbf{4. Answer: C}

\ce{N2(g) + 3H2(g) <=> 2NH3(g)}: 4 moles of gas $\rightarrow$ 2 moles of gas. By Le Chatelier's principle, increasing the pressure shifts the equilibrium towards the side with fewer moles of gas, i.e. towards \ce{NH3}, increasing the yield. (High pressure also increases the rate, but the thermodynamic reason for choosing it is the decrease in mole number; a catalyst, not pressure, affects activation energy; the reaction is exothermic.)
\end{corrige}

\begin{corrige}[Exercise 2 --- True or False with Justification]
\textbf{1. FALSE.} \ce{SO2} is an \emph{acidic} oxide (an oxide of a non-metal). It dissolves in water to form an acidic solution of sulphurous acid:
\[ \ce{SO2(g) + H2O(l) -> H2SO3(aq)} \]
It reacts with bases, e.g. \ce{SO2 + 2NaOH -> Na2SO3 + H2O}, which is the behaviour of an acid anhydride, not a basic oxide.

\textbf{2. TRUE.} Concentrated \ce{H2SO4} removes the elements of water (H and O in the ratio 2:1) from sucrose, leaving a black spongy mass of carbon:
\[ \ensuremath{\mathrm{C_{12}H_{22}O_{11}(s) \rightarrow [conc.\ H_{2}SO_{4}] 12C(s) + 11H_{2}O(l)}} \]
This is the classic ``black snake'' / charring of sugar demonstration.

\textbf{3. FALSE.} Ammonia is a basic gas and reacts with the acid:
\[ \ce{2NH3(g) + H2SO4(l) -> (NH4)2SO4(s)} \]
The gas would be absorbed, not dried. Ammonia is dried using a basic drying agent, quicklime (\ce{CaO}) or soda lime.

\textbf{4. TRUE.} Nitrate- and phosphate-rich run-off stimulates excessive algal growth (algal bloom). When the algae die, aerobic bacteria decompose them and consume dissolved oxygen (high BOD). The resulting oxygen depletion (hypoxia) kills fish and other aquatic life --- this sequence is eutrophication.
\end{corrige}

\begin{corrige}[Exercise 3 --- Definitions and Key Terms]
\textbf{1. Acid rain:} rain (or any precipitation) with a pH below that of normal rain ($\approx 5.6$, due to dissolved \ce{CO2}), typically in the range \textbf{pH 4--5} (sometimes lower). Its main acidic components are \textbf{sulphuric acid} (\ce{H2SO4}, from oxidation of \ce{SO2}) and \textbf{nitric acid} (\ce{HNO3}, from oxidation of \ce{NO} / \ce{NO2}).

\textbf{2. Catalyst:} a substance that increases the rate of a chemical reaction by providing an alternative reaction pathway of \textbf{lower activation energy}, without itself being consumed (it is regenerated at the end). It speeds up both forward and reverse reactions equally, so it \textbf{does not change the position of equilibrium} --- equilibrium is simply reached faster.

\textbf{3. Percentage yield:}
\[ \text{Percentage yield} = \frac{\text{actual mass (or moles) of product obtained}}{\text{theoretical mass (or moles) of product}} \times 100\,\% \]

\textbf{4. Le Chatelier's Principle:} if a constraint (change of concentration, pressure or temperature) is applied to a system at dynamic equilibrium, the system adjusts itself so as to oppose (partially counteract) the imposed change and restore a new equilibrium.
\end{corrige}

\begin{corrige}[Exercise 4 --- Direct Calculations]
\textbf{1. Percentage of nitrogen in urea \ce{CO(NH2)2}:}

Molar mass: $M = 12 + 16 + 2(14 + 2\times 1) = 60\ \mathrm{g\,mol^{-1}}$.
Mass of N per mole: $2 \times 14 = 28\ \mathrm{g}$.
\[ \%N = \frac{28}{60}\times 100 = \boxed{46.7\,\%} \]

\textbf{2. Mass of \ce{H2SO4} from 64 tonnes of S (yield 95\%):}

Overall stoichiometry: \ce{S -> SO2 -> SO3 -> H2SO4}, i.e. 1 mol S gives 1 mol \ce{H2SO4}.
\begin{align*}
n(\ce{S}) &= \frac{64\times 10^{6}\ \mathrm{g}}{32\ \mathrm{g\,mol^{-1}}} = 2.0\times 10^{6}\ \mathrm{mol}\\
m_{\text{theoretical}}(\ce{H2SO4}) &= 2.0\times 10^{6} \times 98 = 196\times 10^{6}\ \mathrm{g} = 196\ \text{tonnes}\\
m_{\text{actual}} &= 0.95 \times 196 = \boxed{186.2\ \text{tonnes}}
\end{align*}

\textbf{3. Ammonia + nitric acid:}
\[ \ce{NH3(g) + HNO3(aq) -> NH4NO3(aq)} \]
(The product may equally be written as the solid \ce{NH4NO3(s)} after evaporation; the ionic form is \ce{NH4+(aq) + NO3-(aq)}.)

\textbf{4. Oxidation numbers of nitrogen:}
\begin{itemize}
\item \ce{NH4+}: $x + 4(+1) = +1 \Rightarrow x = \mathbf{-3}$
\item \ce{N2O}: $2x + (-2) = 0 \Rightarrow x = \mathbf{+1}$
\item \ce{NO2-}: $x + 2(-2) = -1 \Rightarrow x = \mathbf{+3}$
\item \ce{NO3-}: $x + 3(-2) = -1 \Rightarrow x = \mathbf{+5}$
\end{itemize}
\end{corrige}

\begin{corrige}[Exercise 5 --- The Contact Process]
\textbf{1. Combustion of sulphur:}
\[ \ce{S(s) + O2(g) -> SO2(g)} \]

\textbf{2. Catalytic oxidation:}
\[ \ce{2SO2(g) + O2(g) <=> 2SO3(g)} \qquad \Delta H = -196\ \mathrm{kJ\,mol^{-1}} \]
Catalyst: vanadium(V) oxide, \ce{V2O5}; optimum temperature: about \SI{450}{\ensuremath{{}^{\circ}\mathrm{C}}}; pressure: about 1--2 atm (near atmospheric).

\textbf{3. Le Chatelier analysis:}
\begin{enumerate}
\item[(a)] The forward reaction is exothermic, so a \emph{low} temperature favours a high yield of \ce{SO3}. However, at low temperature the rate is too slow (and the catalyst is inactive below about \SI{400}{\ensuremath{{}^{\circ}\mathrm{C}}}). \SI{450}{\ensuremath{{}^{\circ}\mathrm{C}}} is a \textbf{compromise} between a good yield and an economically acceptable rate.
\item[(b)] The forward reaction decreases the number of gas moles ($3 \rightarrow 2$), so high pressure would favour \ce{SO3}. But the conversion at \SI{450}{\ensuremath{{}^{\circ}\mathrm{C}}} and atmospheric pressure is already about 98\%; the extra cost and safety risks of high-pressure plant are not justified by the tiny yield gain.
\item[(c)] Excess air increases the concentration of \ce{O2}; by Le Chatelier's principle the equilibrium shifts to the right, increasing the conversion of the more valuable reactant, \ce{SO2}.
\end{enumerate}

\textbf{4.} \ce{SO3} is dissolved in \textbf{concentrated sulphuric acid} to form oleum (direct dissolution in water would produce a corrosive, hard-to-condense acid mist because the reaction is too violent):
\[ \ce{SO3(g) + H2SO4(l) -> H2S2O7(l)} \]
The oleum is then diluted:
\[ \ce{H2S2O7(l) + H2O(l) -> 2H2SO4(l)} \]
\end{corrige}

\begin{corrige}[Exercise 6 --- Redox Chemistry of Sulphur and Nitrogen]
\textbf{1. Oxidation states of S:} \ce{H2S}: $-2$; \ce{S8}: $0$ (element); \ce{SO2}: $+4$; \ce{SO3}: $+6$; \ce{H2SO4}: $+6$; \ce{Na2S2O3}: $2(+1) + 2x + 3(-2) = 0 \Rightarrow x = +2$ (average).

\textbf{2. Oxidation states of N:} \ce{NH3}: $-3$; \ce{N2O}: $+1$; \ce{NO}: $+2$; \ce{NO2}: $+4$; \ce{HNO3}: $+5$; \ce{NH4NO3}: $-3$ in \ce{NH4+} and $+5$ in \ce{NO3-}.

\textbf{3.} \ce{2H2S(g) + SO2(g) -> 3S(s) + 2H2O(l)}

S in \ce{H2S} goes from $-2$ to $0$ (oxidation); S in \ce{SO2} goes from $+4$ to $0$ (reduction).
\begin{itemize}
\item \textbf{Oxidising agent:} \ce{SO2} (it is reduced).
\item \textbf{Reducing agent:} \ce{H2S} (it is oxidised).
\end{itemize}
Half-equations:
\[ \text{Oxidation: } \ce{H2S -> S + 2H+ + 2e-} \qquad \text{Reduction: } \ce{SO2 + 4H+ + 4e- -> S + 2H2O} \]

\textbf{4. Dual behaviour of \ce{SO2} (S in $+4$, an intermediate state):}
\begin{itemize}
\item \emph{As a reducing agent} (S oxidised $+4 \rightarrow +6$), e.g. with acidified potassium manganate(VII):
\[ \ce{2KMnO4 + 5SO2 + 2H2O -> K2SO4 + 2MnSO4 + 2H2SO4} \]
(purple solution decolourised), or simply \ce{2SO2 + O2 -> 2SO3}.
\item \emph{As an oxidising agent} (S reduced $+4 \rightarrow 0$), e.g. with \ce{H2S}:
\[ \ce{SO2 + 2H2S -> 3S + 2H2O} \]
\end{itemize}
\end{corrige}

\begin{corrige}[Exercise 7 --- The Haber Process]
\textbf{1. Increasing pressure:}
\begin{enumerate}
\item[(a)] \emph{Yield:} increases. The forward reaction converts 4 moles of gas into 2 moles; by Le Chatelier's principle, compression shifts the equilibrium towards the side with fewer gas moles, i.e. towards \ce{NH3}.
\item[(b)] \emph{Rate:} increases. Higher pressure means higher concentration of gas molecules, hence more frequent successful collisions per unit time.
\end{enumerate}

\textbf{2. Increasing temperature:}
\begin{enumerate}
\item[(a)] \emph{Yield:} decreases. The forward reaction is exothermic ($\Delta H < 0$); raising the temperature shifts the equilibrium in the endothermic (reverse) direction, decomposing \ce{NH3}.
\item[(b)] \emph{Rate:} increases. More molecules possess energy $\geq E_a$, so the frequency of successful collisions rises sharply (the rate constant grows exponentially with $T$). Hence the industrial compromise at about \SI{450}{\ensuremath{{}^{\circ}\mathrm{C}}}.
\end{enumerate}

\textbf{3. Iron catalyst:} it provides an alternative pathway of lower activation energy, so equilibrium is reached faster. It does \textbf{not} affect the yield (position of equilibrium), since it accelerates forward and reverse reactions equally.

\textbf{4. Volume of \ce{NH3} at s.t.p.:}
\begin{align*}
n(\ce{N2}) &= \frac{28}{28} = 1.0\ \mathrm{mol}\\
n(\ce{N2})\text{ reacted} &= 0.30 \times 1.0 = 0.30\ \mathrm{mol}\\
n(\ce{NH3}) &= 2 \times 0.30 = 0.60\ \mathrm{mol} \quad (\text{N2 : NH3 = 1:2})\\
V(\ce{NH3}) &= 0.60 \times 22.4 = \boxed{13.44\ \mathrm{dm^{3}}\ (\approx 13.4\ \mathrm{dm^{3}})}
\end{align*}
\end{corrige}

\begin{corrige}[Exercise 8 --- Fertilizers: Composition and Economics]
\textbf{1. Percentage of nitrogen:}
\begin{enumerate}
\item[(a)] \ce{NH4NO3}: $M = 14+4+14+48 = 80$; mass of N $= 28$.
\[ \%N = \frac{28}{80}\times 100 = 35.0\,\% \]
\item[(b)] \ce{(NH4)2SO4}: $M = 2(18) + 32 + 64 = 132$; mass of N $= 28$.
\[ \%N = \frac{28}{132}\times 100 = 21.2\,\% \]
\item[(c)] \ce{CO(NH2)2}: $M = 60$; mass of N $= 28$.
\[ \%N = \frac{28}{60}\times 100 = 46.7\,\% \]
\end{enumerate}

\textbf{2. Mass of fertilizer to supply \SI{46}{\kilo\gram} of N:}
\begin{align*}
m(\ce{NH4NO3}) &= \frac{46}{0.350} = 131.4\ \mathrm{kg}\\
m(\ce{(NH4)2SO4}) &= \frac{46}{0.212} = 217.0\ \mathrm{kg}\\
m(\text{urea}) &= \frac{46}{0.467} = 98.5\ \mathrm{kg}
\end{align*}

\textbf{3. Most economical fertilizer (cost per kg of N):}

Cost per kg of fertilizer: \ce{NH4NO3}: $18000/50 = 360$ FCFA/kg; \ce{(NH4)2SO4}: $15000/50 = 300$ FCFA/kg; urea: $20000/50 = 400$ FCFA/kg.

Cost per kg of N supplied:
\begin{align*}
\ce{NH4NO3}&: \frac{360}{0.350} = 1029\ \text{FCFA/kg N}\\
\ce{(NH4)2SO4}&: \frac{300}{0.212} = 1415\ \text{FCFA/kg N}\\
\text{Urea}&: \frac{400}{0.467} = 857\ \text{FCFA/kg N}
\end{align*}
\textbf{Urea is the most economical per kilogram of nitrogen} (857 FCFA/kg N), despite having the most expensive bag, because of its very high N content (46.7\%).

\textbf{4. Environmental problems of excess nitrogenous fertilizers:}
\begin{itemize}
\item \textbf{Eutrophication:} nitrate leaching into rivers and lakes causes algal blooms, oxygen depletion and death of aquatic life.
\item \textbf{Groundwater contamination:} high nitrate levels in drinking water (linked to methaemoglobinaemia in infants --- ``blue baby syndrome'').
\item (Also acceptable: soil acidification from ammonium-based fertilizers; release of the greenhouse gas \ce{N2O}.)
\end{itemize}
\end{corrige}

\begin{corrige}[Exercise 9 --- Sulphur Oxides and Acid Rain (15 marks)]
\textbf{1.(a)} Roasting of zinc blende: \hfill [2]
\[ \ce{2ZnS(s) + 3O2(g) -> 2ZnO(s) + 2SO2(g)} \]
\emph{[1 mark for correct formulae of reactants and products; 1 mark for balancing and state symbols.]}

\textbf{1.(b)} Formation of acid rain: \hfill [4]
\ce{SO2} is first oxidised in the atmosphere (by \ce{O2}, catalysed by sunlight, dust or \ce{NO2}):
\[ \ce{2SO2(g) + O2(g) -> 2SO3(g)} \]
The \ce{SO3} then dissolves in rainwater:
\[ \ce{SO3(g) + H2O(l) -> H2SO4(aq)} \]
(Alternative route also credited: \ce{SO2 + H2O -> H2SO3}, then \ce{2H2SO3 + O2 -> 2H2SO4}.) The rain therefore contains dilute sulphuric acid and falls with pH 4--5.
\emph{[1 mark per correct balanced equation (2); 1 mark for mentioning atmospheric oxidation conditions; 1 mark for linking to acidic rain / low pH.]}

\textbf{1.(c)} Three harmful effects of acid rain: \hfill [3]
\begin{itemize}
\item Acidifies lakes and streams, killing fish and aquatic organisms.
\item Leaches nutrients (and toxic \ce{Al^3+}) from soils; damages leaves and forests, reducing crop yields.
\item Corrodes limestone/marble buildings and metallic structures (roofs, bridges).
\end{itemize}
\emph{[1 mark each, any three valid effects.]}

\textbf{2.(a)} Flue gas desulphurisation: \hfill [3]
\[ \ce{2SO2(g) + 2CaCO3(s) + O2(g) + 4H2O(l) -> 2CaSO4.2H2O(s) + 2CO2(g)} \]
\emph{[1 mark reactants, 1 mark products including gypsum formula, 1 mark balancing.]}

\textbf{2.(b)} Limestone is preferred because it is \textbf{very cheap and abundant} (often quarried locally), whereas \ce{NaOH} is expensive; in addition the product, gypsum, is a \textbf{useful saleable solid} (plaster, cement, wallboard) rather than a waste solution requiring disposal. \hfill [2]

\textbf{3.} Use of \ce{H2SO4} in Cameroon: manufacture of \textbf{fertilizers} (e.g. phosphate/ammonium sulphate fertilizers); also acceptable: car battery acid, metal processing, petroleum refining. \hfill [1]
\end{corrige}

\begin{corrige}[Exercise 10 --- Laboratory Preparation of Ammonia (12 marks)]
\textbf{1.(a)} Reagents: a mixture of solid \textbf{ammonium chloride} (\ce{NH4Cl}) and solid \textbf{calcium hydroxide} (slaked lime, \ce{Ca(OH)2}). \hfill [1]
\[ \ce{2NH4Cl(s) + Ca(OH)2(s) -> CaCl2(s) + 2NH3(g) + 2H2O(g)} \]

\textbf{1.(b)} Apparatus (description): \hfill [3]
The solid mixture is placed in a round-bottomed (or test) flask held by a clamp, tilted slightly \emph{downwards} at the mouth so that condensed water does not run back onto the hot glass. A delivery tube leads the gas first through a drying tower/U-tube packed with quicklime, then into an inverted (upside-down) dry gas jar. Heat is applied gently with a Bunsen burner.
\emph{[1 mark heat source + reactants in flask; 1 mark drying unit in line; 1 mark upward delivery / inverted collection jar.]}

\textbf{1.(c)} Drying agent: \textbf{quicklime, \ce{CaO}} (or soda lime). It is suitable because it is a \textbf{basic} drying agent and therefore does not react with the basic gas \ce{NH3}. \hfill [2]

\textbf{1.(d)} Collection: by \textbf{upward delivery} (downward displacement of air) in an inverted gas jar, because ammonia is \textbf{less dense than air} and is \textbf{extremely soluble in water} (so it cannot be collected over water). \hfill [2]

\textbf{2.} Test: \ce{NH3} turns \textbf{damp red litmus paper blue} (it is the only common alkaline gas); alternatively, it gives dense white fumes of \ce{NH4Cl} with a rod dipped in concentrated \ce{HCl}. \hfill [1]

\textbf{3.} Concentrated \ce{H2SO4} is acidic and reacts with the basic ammonia instead of drying it: \hfill [2]
\[ \ce{2NH3(g) + H2SO4(l) -> (NH4)2SO4(s)} \]

\textbf{4.} Deep blue complex with \ce{Cu^2+}: \hfill [1]
\[ \ce{Cu^{2+}(aq) + 4NH3(aq) -> [Cu(NH3)4]^{2+}(aq)} \]
(Examiner's note: credit also the two-step answer showing the initial pale blue precipitate \ce{Cu(OH)2} dissolving in excess ammonia to give the deep blue tetraamminecopper(II) complex.)
\end{corrige}

\begin{corrige}[Exercise 11 --- Eutrophication and Oxygen Sag Curve (13 marks)]
\textbf{1. Graph:} \hfill [3]
$x$-axis: sampling points A $\rightarrow$ E (equally spaced). Left $y$-axis: \ce{[NO3-]} (\SI{}{\mg\per\L}), 0--50; right $y$-axis: DO (\SI{}{\mg\per\L}), 0--10. The nitrate curve rises steeply from A to a peak at C (45), then falls through D to E. The DO curve falls from 9.0 (A) to a minimum of 3.0 at C, then recovers through D to 8.0 at E --- the classic \textbf{oxygen sag curve}, roughly a mirror image of the nitrate curve.
\emph{[1 mark correct axes with labels/units; 1 mark correct nitrate plot; 1 mark correct DO plot.]}

\textbf{2. Entry point of pollution:} \hfill [2]
Between \textbf{A and B} (most evident at B). Justification: \ce{[NO3-]} jumps from \SI{2}{\mg\per\L} (A, upstream background) to \SI{15}{\mg\per\L} at B, and DO simultaneously begins to fall (9.0 $\rightarrow$ 7.5). The sharp rise in nitrate just downstream of the farm identifies the run-off entry point.

\textbf{3. Explanation of the DO fall from A to C (eutrophication):} \hfill [4]
\begin{itemize}
\item Nitrate leached from the fertilized field is a plant nutrient; it causes an \textbf{algal bloom} (rapid growth of algae and aquatic plants) near B--C.
\item The dense surface growth also shades submerged plants, reducing photosynthetic oxygen production.
\item When the algae die, they are decomposed by \textbf{aerobic bacteria}, whose respiration consumes large amounts of dissolved oxygen --- the \textbf{biochemical oxygen demand (BOD)} rises sharply.
\item Oxygen consumption by decomposition exceeds re-oxygenation from the air and photosynthesis, so DO falls to its minimum (3.0) at C; fish and other aerobic organisms suffocate.
\end{itemize}
\emph{[1 mark nutrient $\rightarrow$ algal bloom; 1 mark death of algae; 1 mark aerobic bacterial decomposition raising BOD; 1 mark consumption $>$ supply $\Rightarrow$ DO minimum.]}

\textbf{4. Recovery from C to E:} \hfill [2]
Downstream, the nitrate is diluted and taken up by plants, so algal growth and the load of dead organic matter decrease; bacterial activity and BOD fall. Meanwhile the stream is continuously \textbf{re-aerated} by dissolution of oxygen from the atmosphere (and photosynthesis resumes), so DO rises back towards its natural level (8.0 at E).

\textbf{5. Two measures for a farmer around Mount Bamboutos:} \hfill [2]
\begin{itemize}
\item Apply fertilizer in \textbf{split doses matched to crop needs} and avoid spreading just before heavy rain, so less nitrate is left to leach.
\item Maintain \textbf{buffer strips of grass/vegetation along streams} and practise contour ploughing or terracing on the steep volcanic slopes to trap run-off.
\end{itemize}
(Also credited: using organic manure/compost instead of excess mineral fertilizer; planting cover crops in the off-season.)
\end{corrige}

\newpage

\coursec{Summary sheet}

\begin{popfigure}
\begin{tikzpicture}[mindmap, mindmap]
\node[text width=0.9\linewidth, text width=0.9\linewidth, text width=0.9\linewidth, text width=0.9\linewidth, text width=0.9\linewidth, text width=0.9\linewidth, text width=0.9\linewidth, text width=0.9\linewidth, text width=0.9\linewidth, text width=0.9\linewidth, text width=0.9\linewidth, text width=0.9\linewidth, text width=0.9\linewidth, text width=0.9\linewidth, text width=0.9\linewidth, text width=0.9\linewidth, text width=0.9\linewidth, text width=0.9\linewidth, text width=0.9\linewidth, text width=0.9\linewidth, text width=0.9\linewidth, text width=0.9\linewidth, text width=0.9\linewidth, text width=0.9\linewidth, text width=0.9\linewidth, text width=0.9\linewidth, text width=0.9\linewidth, text width=0.9\linewidth, text width=0.9\linewidth, text width=0.9\linewidth, text width=0.9\linewidth, text width=0.9\linewidth, text width=0.9\linewidth, text width=0.9\linewidth, text width=0.9\linewidth, text width=0.9\linewidth, text width=0.9\linewidth, text width=0.9\linewidth, text width=0.9\linewidth, text width=0.9\linewidth, text width=0.9\linewidth, text width=0.9\linewidth, text width=0.9\linewidth, text width=0.9\linewidth, text width=0.9\linewidth, text width=0.9\linewidth, text width=0.9\linewidth, text width=0.9\linewidth, text width=0.9\linewidth, text width=0.9\linewidth, text width=0.9\linewidth, text width=0.9\linewidth, text width=0.9\linewidth, text width=0.9\linewidth, text width=0.9\linewidth, text width=0.9\linewidth, text width=0.9\linewidth, text width=0.9\linewidth, text width=0.9\linewidth, text width=0.9\linewidth, text width=0.9\linewidth, text width=0.9\linewidth, text width=0.9\linewidth, text width=0.9\linewidth, text width=0.9\linewidth, text width=0.9\linewidth, text width=0.9\linewidth, text width=0.9\linewidth, text width=0.9\linewidth, text width=0.9\linewidth, text width=0.9\linewidth, text width=0.9\linewidth, text width=0.9\linewidth, text width=0.9\linewidth, text width=0.9\linewidth, text width=0.9\linewidth, text width=0.9\linewidth, text width=0.9\linewidth, text width=0.9\linewidth, text width=0.9\linewidth, text width=0.9\linewidth, text width=0.9\linewidth, text width=0.9\linewidth, text width=0.9\linewidth, text width=0.9\linewidth, text width=0.9\linewidth, text width=0.9\linewidth, text width=0.9\linewidth, text width=0.9\linewidth, text width=0.9\linewidth, text width=0.9\linewidth, text width=0.9\linewidth, text width=0.9\linewidth, text width=0.9\linewidth, text width=0.9\linewidth, text width=0.9\linewidth, text width=0.9\linewidth, text width=0.9\linewidth, text width=0.9\linewidth, text width=0.9\linewidth, text width=0.9\linewidth, text width=0.9\linewidth, text width=0.9\linewidth, text width=0.9\linewidth, text width=0.9\linewidth, text width=0.9\linewidth, text width=0.9\linewidth, text width=0.9\linewidth, text width=0.9\linewidth, text width=0.9\linewidth, text width=0.9\linewidth, text width=0.9\linewidth, text width=0.9\linewidth, text width=0.9\linewidth, text width=0.9\linewidth, text width=0.9\linewidth, text width=0.9\linewidth, text width=0.9\linewidth, text width=0.9\linewidth, text width=0.9\linewidth, text width=0.9\linewidth, text width=0.9\linewidth, text width=0.9\linewidth, text width=0.9\linewidth, text width=0.9\linewidth, text width=0.9\linewidth, text width=0.9\linewidth, text width=0.9\linewidth, text width=0.9\linewidth, text width=0.9\linewidth, text width=0.9\linewidth, text width=0.9\linewidth, text width=0.9\linewidth, text width=0.9\linewidth, text width=0.9\linewidth, text width=0.9\linewidth, text width=0.9\linewidth, text width=0.9\linewidth, text width=0.9\linewidth, text width=0.9\linewidth, text width=0.9\linewidth, text width=0.9\linewidth, text width=0.9\linewidth, text width=0.9\linewidth, text width=0.9\linewidth, text width=0.9\linewidth, text width=0.9\linewidth, text width=0.9\linewidth, text width=0.9\linewidth, text width=0.9\linewidth, text width=0.9\linewidth, text width=0.9\linewidth, text width=0.9\linewidth, text width=0.9\linewidth, text width=0.9\linewidth, text width=0.9\linewidth, text width=0.9\linewidth, text width=0.9\linewidth, text width=0.9\linewidth, text width=0.9\linewidth, text width=0.9\linewidth, text width=0.9\linewidth, text width=0.9\linewidth, text width=0.9\linewidth, text width=0.9\linewidth, text width=0.9\linewidth, text width=0.9\linewidth, text width=0.9\linewidth, text width=0.9\linewidth, text width=0.9\linewidth, text width=0.9\linewidth, text width=0.9\linewidth, text width=0.9\linewidth, text width=0.9\linewidth, text width=0.9\linewidth, text width=0.9\linewidth, text width=0.9\linewidth, text width=0.9\linewidth, text width=0.9\linewidth, text width=0.9\linewidth, text width=0.9\linewidth, text width=0.9\linewidth, text width=0.9\linewidth, text width=0.9\linewidth, text width=0.9\linewidth, text width=0.9\linewidth, text width=0.9\linewidth, text width=0.9\linewidth, text width=0.9\linewidth, text width=0.9\linewidth, text width=0.9\linewidth, text width=0.9\linewidth, text width=0.9\linewidth, text width=0.9\linewidth, text width=0.9\linewidth, text width=0.9\linewidth, text width=0.9\linewidth, text width=0.9\linewidth, text width=0.9\linewidth, text width=0.9\linewidth, text width=0.9\linewidth, text width=0.9\linewidth, text width=0.9\linewidth, text width=0.9\linewidth, text width=0.9\linewidth, text width=0.9\linewidth, text width=0.9\linewidth, text width=0.9\linewidth, text width=0.9\linewidth, text width=0.9\linewidth, text width=0.9\linewidth, text width=0.9\linewidth, text width=0.9\linewidth, text width=0.9\linewidth, text width=0.9\linewidth, text width=0.9\linewidth, text width=0.9\linewidth, text width=0.9\linewidth, text width=0.9\linewidth, text width=0.9\linewidth, text width=0.9\linewidth, text width=0.9\linewidth, text width=0.9\linewidth, text width=0.9\linewidth, text width=0.9\linewidth, text width=0.9\linewidth, text width=0.9\linewidth, text width=0.9\linewidth, text width=0.9\linewidth, text width=0.9\linewidth, text width=0.9\linewidth, text width=0.9\linewidth, text width=0.9\linewidth, text width=0.9\linewidth, text width=0.9\linewidth, text width=0.9\linewidth, text width=0.9\linewidth, text width=0.9\linewidth, text width=0.9\linewidth, text width=0.9\linewidth, text width=0.9\linewidth, text width=0.9\linewidth, text width=0.9\linewidth, text width=0.9\linewidth, text width=0.9\linewidth, text width=0.9\linewidth, text width=0.9\linewidth, text width=0.9\linewidth, text width=0.9\linewidth, text width=0.9\linewidth, text width=0.9\linewidth, text width=0.9\linewidth, text width=0.9\linewidth, text width=0.9\linewidth, text width=0.9\linewidth, text width=0.9\linewidth, text width=0.9\linewidth, text width=0.9\linewidth, text width=0.9\linewidth, text width=0.9\linewidth, text width=0.9\linewidth, text width=0.9\linewidth, text width=0.9\linewidth, text width=0.9\linewidth, text width=0.9\linewidth, text width=0.9\linewidth, text width=0.9\linewidth, text width=0.9\linewidth, text width=0.9\linewidth, text width=0.9\linewidth, text width=0.9\linewidth, text width=0.9\linewidth, text width=0.9\linewidth, text width=0.9\linewidth, text width=0.9\linewidth, text width=0.9\linewidth, text width=0.9\linewidth, text width=0.9\linewidth, text width=0.9\linewidth, text width=0.9\linewidth, text width=0.9\linewidth, text width=0.9\linewidth, text width=0.9\linewidth, text width=0.9\linewidth, text width=0.9\linewidth, text width=0.9\linewidth, text width=0.9\linewidth, text width=0.9\linewidth, text width=0.9\linewidth, text width=0.9\linewidth, text width=0.9\linewidth, text width=0.9\linewidth, text width=0.9\linewidth, text width=0.9\linewidth, text width=0.9\linewidth, text width=0.9\linewidth, text width=0.9\linewidth, text width=0.9\linewidth, text width=0.9\linewidth, text width=0.9\linewidth, text width=0.9\linewidth, text width=0.9\linewidth, text width=0.9\linewidth, text width=0.9\linewidth, text width=0.9\linewidth, text width=0.9\linewidth, text width=0.9\linewidth, text width=0.9\linewidth, text width=0.9\linewidth, text width=0.9\linewidth, text width=0.9\linewidth, text width=0.9\linewidth, text width=0.9\linewidth, text width=0.9\linewidth, text width=0.9\linewidth, text width=0.9\linewidth, text width=0.9\linewidth, text width=0.9\linewidth, text width=0.9\linewidth, text width=0.9\linewidth, text width=0.9\linewidth, text width=0.9\linewidth, text width=0.9\linewidth, text width=0.9\linewidth, text width=0.9\linewidth, text width=0.9\linewidth, text width=0.9\linewidth, text width=0.9\linewidth, text width=0.9\linewidth, text width=0.9\linewidth, text width=0.9\linewidth, text width=0.9\linewidth, text width=0.9\linewidth, text width=0.9\linewidth, text width=0.9\linewidth, text width=0.9\linewidth, text width=0.9\linewidth, text width=0.9\linewidth, text width=0.9\linewidth, text width=0.9\linewidth, text width=0.9\linewidth, text width=0.9\linewidth, text width=0.9\linewidth, text width=0.9\linewidth, text width=0.9\linewidth, text width=0.9\linewidth, text width=0.9\linewidth, text width=0.9\linewidth, text width=0.9\linewidth, text width=0.9\linewidth, text width=0.9\linewidth, text width=0.9\linewidth, text width=0.9\linewidth, text width=0.9\linewidth, text width=0.9\linewidth, text width=0.9\linewidth, text width=0.9\linewidth, text width=0.9\linewidth, text width=0.9\linewidth, text width=0.9\linewidth, text width=0.9\linewidth, text width=0.9\linewidth, text width=0.9\linewidth, text width=0.9\linewidth, text width=0.9\linewidth, text width=0.9\linewidth, text width=0.9\linewidth, text width=0.9\linewidth, text width=0.9\linewidth, text width=0.9\linewidth, text width=0.9\linewidth, text width=0.9\linewidth, text width=0.9\linewidth, text width=0.9\linewidth, text width=0.9\linewidth, text width=0.9\linewidth, text width=0.9\linewidth, text width=0.9\linewidth, text width=0.9\linewidth, text width=0.9\linewidth, text width=0.9\linewidth, text width=0.9\linewidth, text width=0.9\linewidth, text width=0.9\linewidth, text width=0.9\linewidth, text width=0.9\linewidth, text width=0.9\linewidth, text width=0.9\linewidth, text width=0.9\linewidth, text width=0.9\linewidth, text width=0.9\linewidth, text width=0.9\linewidth, text width=0.9\linewidth, text width=0.9\linewidth, text width=0.9\linewidth, text width=0.9\linewidth, text width=0.9\linewidth, text width=0.9\linewidth, text width=0.9\linewidth, text width=0.9\linewidth, text width=0.9\linewidth, text width=0.9\linewidth, text width=0.9\linewidth, text width=0.9\linewidth, text width=0.9\linewidth, text width=0.9\linewidth, text width=0.9\linewidth, text width=0.9\linewidth, text width=0.9\linewidth, text width=0.9\linewidth, text width=0.9\linewidth, text width=0.9\linewidth, text width=0.9\linewidth, text width=0.9\linewidth, text width=0.9\linewidth, text width=0.9\linewidth, text width=0.9\linewidth, text width=0.9\linewidth, text width=0.9\linewidth, text width=0.9\linewidth, text width=0.9\linewidth, text width=0.9\linewidth, text width=0.9\linewidth, text width=0.9\linewidth, text width=0.9\linewidth, text width=0.9\linewidth, text width=0.9\linewidth, text width=0.9\linewidth, text width=0.9\linewidth, text width=0.9\linewidth, text width=0.9\linewidth, text width=0.9\linewidth, text width=0.9\linewidth, text width=0.9\linewidth, text width=0.9\linewidth, text width=0.9\linewidth, text width=0.9\linewidth, text width=0.9\linewidth, text width=0.9\linewidth, text width=0.9\linewidth, text width=0.9\linewidth, text width=0.9\linewidth, text width=0.9\linewidth, text width=0.9\linewidth, text width=0.9\linewidth, text width=0.9\linewidth, text width=0.9\linewidth, text width=0.9\linewidth, text width=0.9\linewidth, text width=0.9\linewidth, text width=0.9\linewidth, text width=0.9\linewidth, text width=0.9\linewidth, text width=0.9\linewidth, text width=0.9\linewidth, text width=0.9\linewidth, text width=0.9\linewidth, text width=0.9\linewidth, text width=0.9\linewidth, text width=0.9\linewidth, text width=0.9\linewidth, text width=0.9\linewidth, text width=0.9\linewidth, text width=0.9\linewidth, text width=0.9\linewidth, text width=0.9\linewidth, text width=0.9\linewidth, text width=0.9\linewidth, text width=0.9\linewidth, text width=0.9\linewidth, text width=0.9\linewidth, text width=0.9\linewidth, text width=0.9\linewidth, text width=0.9\linewidth, text width=0.9\linewidth, text width=0.9\linewidth, text width=0.9\linewidth, text width=0.9\linewidth, text width=0.9\linewidth, text width=0.9\linewidth, text width=0.9\linewidth, text width=0.9\linewidth, text width=0.9\linewidth, text width=0.9\linewidth, text width=0.9\linewidth, text width=0.9\linewidth, text width=0.9\linewidth, text width=0.9\linewidth, text width=0.9\linewidth, text width=0.9\linewidth, text width=0.9\linewidth, text width=0.9\linewidth, text width=0.9\linewidth, text width=0.9\linewidth, text width=0.9\linewidth, text width=0.9\linewidth, text width=0.9\linewidth, text width=0.9\linewidth, text width=0.9\linewidth, text width=0.9\linewidth, text width=0.9\linewidth, text width=0.9\linewidth, text width=0.9\linewidth, text width=0.9\linewidth, text width=0.9\linewidth, text width=0.9\linewidth, text width=0.9\linewidth, text width=0.9\linewidth, text width=0.9\linewidth, text width=0.9\linewidth, text width=0.9\linewidth, text width=0.9\linewidth, text width=0.9\linewidth, text width=0.9\linewidth, text width=0.9\linewidth, text width=0.9\linewidth, text width=0.9\linewidth, text width=0.9\linewidth, text width=0.9\linewidth, text width=0.9\linewidth, text width=0.9\linewidth, text width=0.9\linewidth, text width=0.9\linewidth, text width=0.9\linewidth, text width=0.9\linewidth, text width=0.9\linewidth, text width=0.9\linewidth, text width=0.9\linewidth, text width=0.9\linewidth, text width=0.9\linewidth, text width=0.9\linewidth, text width=0.9\linewidth, text width=0.9\linewidth, text width=0.9\linewidth, text width=0.9\linewidth, text width=0.9\linewidth, text width=0.9\linewidth, text width=0.9\linewidth, text width=0.9\linewidth, text width=0.9\linewidth, text width=0.9\linewidth, text width=0.9\linewidth, text width=0.9\linewidth, text width=0.9\linewidth, text width=0.9\linewidth, text width=0.9\linewidth, text width=0.9\linewidth, text width=0.9\linewidth, text width=0.9\linewidth, text width=0.9\linewidth, text width=0.9\linewidth, text width=0.9\linewidth, text width=0.9\linewidth, text width=0.9\linewidth, text width=0.9\linewidth, text width=0.9\linewidth, text width=0.9\linewidth, text width=0.9\linewidth, text width=0.9\linewidth, text width=0.9\linewidth, text width=0.9\linewidth, text width=0.9\linewidth, text width=0.9\linewidth, text width=0.9\linewidth, text width=0.9\linewidth, text width=0.9\linewidth, text width=0.9\linewidth, text width=0.9\linewidth, text width=0.9\linewidth, text width=0.9\linewidth, text width=0.9\linewidth, text width=0.9\linewidth, text width=0.9\linewidth, text width=0.9\linewidth, text width=0.9\linewidth, text width=0.9\linewidth, text width=0.9\linewidth, text width=0.9\linewidth, text width=0.9\linewidth, text width=0.9\linewidth, text width=0.9\linewidth, text width=0.9\linewidth, text width=0.9\linewidth, text width=0.9\linewidth, text width=0.9\linewidth, text width=0.9\linewidth, text width=0.9\linewidth, text width=0.9\linewidth, text width=0.9\linewidth, text width=0.9\linewidth, text width=0.9\linewidth, text width=0.9\linewidth, text width=0.9\linewidth, text width=0.9\linewidth, text width=0.9\linewidth, text width=0.9\linewidth, text width=0.9\linewidth, text width=0.9\linewidth, text width=0.9\linewidth, text width=0.9\linewidth, text width=0.9\linewidth, text width=0.9\linewidth, text width=0.9\linewidth, text width=0.9\linewidth, text width=0.9\linewidth, text width=0.9\linewidth, text width=0.9\linewidth, text width=0.9\linewidth, text width=0.9\linewidth, text width=0.9\linewidth, text width=0.9\linewidth, text width=0.9\linewidth, text width=0.9\linewidth, text width=0.9\linewidth, text width=0.9\linewidth, text width=0.9\linewidth, text width=0.9\linewidth, text width=0.9\linewidth, text width=0.9\linewidth, text width=0.9\linewidth, text width=0.9\linewidth, text width=0.9\linewidth, text width=0.9\linewidth, text width=0.9\linewidth, text width=0.9\linewidth, text width=0.9\linewidth, text width=0.9\linewidth, text width=0.9\linewidth, text width=0.9\linewidth, text width=0.9\linewidth, text width=0.9\linewidth, text width=0.9\linewidth, text width=0.9\linewidth, text width=0.9\linewidth, text width=0.9\linewidth, text width=0.9\linewidth, text width=0.9\linewidth, text width=0.9\linewidth, text width=0.9\linewidth, text width=0.9\linewidth, text width=0.9\linewidth, text width=0.9\linewidth, text width=0.9\linewidth, text width=0.9\linewidth, text width=0.9\linewidth, text width=0.9\linewidth, text width=0.9\linewidth, text width=0.9\linewidth, text width=0.9\linewidth, text width=0.9\linewidth, text width=0.9\linewidth, text width=0.9\linewidth, text width=0.9\linewidth, text width=0.9\linewidth, text width=0.9\linewidth, text width=0.9\linewidth, text width=0.9\linewidth, text width=0.9\linewidth, text width=0.9\linewidth, text width=0.9\linewidth, text width=0.9\linewidth, text width=0.9\linewidth, text width=0.9\linewidth, text width=0.9\linewidth, text width=0.9\linewidth, text width=0.9\linewidth, text width=0.9\linewidth, text width=0.9\linewidth, text width=0.9\linewidth, text width=0.9\linewidth, text width=0.9\linewidth, text width=0.9\linewidth, text width=0.9\linewidth, text width=0.9\linewidth, text width=0.9\linewidth, text width=0.9\linewidth, text width=0.9\linewidth, text width=0.9\linewidth, text width=0.9\linewidth, text width=0.9\linewidth, text width=0.9\linewidth, text width=0.9\linewidth, text width=0.9\linewidth, text width=0.9\linewidth, text width=0.9\linewidth, text width=0.9\linewidth, text width=0.9\linewidth, text width=0.9\linewidth, text width=0.9\linewidth, text width=0.9\linewidth, text width=0.9\linewidth, text width=0.9\linewidth, text width=0.9\linewidth, text width=0.9\linewidth, text width=0.9\linewidth, text width=0.9\linewidth, text width=0.9\linewidth, text width=0.9\linewidth, text width=0.9\linewidth, text width=0.9\linewidth, text width=0.9\linewidth, text width=0.9\linewidth, text width=0.9\linewidth, text width=0.9\linewidth, text width=0.9\linewidth, text width=0.9\linewidth, text width=0.9\linewidth, text width=0.9\linewidth, text width=0.9\linewidth, text width=0.9\linewidth, text width=0.9\linewidth, text width=0.9\linewidth, text width=0.9\linewidth, text width=0.9\linewidth, text width=0.9\linewidth, text width=0.9\linewidth, text width=0.9\linewidth, text width=0.9\linewidth, text width=0.9\linewidth, text width=0.9\linewidth, text width=0.9\linewidth, text width=0.9\linewidth, text width=0.9\linewidth, text width=0.9\linewidth, text width=0.9\linewidth, text width=0.9\linewidth, text width=0.9\linewidth, text width=0.9\linewidth, text width=0.9\linewidth, text width=0.9\linewidth, text width=0.9\linewidth, text width=0.9\linewidth, text width=0.9\linewidth, text width=0.9\linewidth, text width=0.9\linewidth, text width=0.9\linewidth, text width=0.9\linewidth, text width=0.9\linewidth, text width=0.9\linewidth, text width=0.9\linewidth, text width=0.9\linewidth, text width=0.9\linewidth, text width=0.9\linewidth, text width=0.9\linewidth, text width=0.9\linewidth, text width=0.9\linewidth, text width=0.9\linewidth, text width=0.9\linewidth, text width=0.9\linewidth, text width=0.9\linewidth, text width=0.9\linewidth, text width=0.9\linewidth, text width=0.9\linewidth, text width=0.9\linewidth, text width=0.9\linewidth, text width=0.9\linewidth, text width=0.9\linewidth, text width=0.9\linewidth, text width=0.9\linewidth, text width=0.9\linewidth, text width=0.9\linewidth, text width=0.9\linewidth, text width=0.9\linewidth, text width=0.9\linewidth, text width=0.9\linewidth, text width=0.9\linewidth, text width=0.9\linewidth, text width=0.9\linewidth, text width=0.9\linewidth, text width=0.9\linewidth, text width=0.9\linewidth, text width=0.9\linewidth, text width=0.9\linewidth, text width=0.9\linewidth, text width=0.9\linewidth, text width=0.9\linewidth, text width=0.9\linewidth, text width=0.9\linewidth, text width=0.9\linewidth, text width=0.9\linewidth, text width=0.9\linewidth, text width=0.9\linewidth, text width=0.9\linewidth, text width=0.9\linewidth, text width=0.9\linewidth, text width=0.9\linewidth, text width=0.9\linewidth, text width=0.9\linewidth, text width=0.9\linewidth, text width=0.9\linewidth, text width=0.9\linewidth, text width=0.9\linewidth, text width=0.9\linewidth, text width=0.9\linewidth, text width=0.9\linewidth, text width=0.9\linewidth, text width=0.9\linewidth, text width=0.9\linewidth, text width=0.9\linewidth, text width=0.9\linewidth, text width=0.9\linewidth, text width=0.9\linewidth, text width=0.9\linewidth, text width=0.9\linewidth, text width=0.9\linewidth, text width=0.9\linewidth, text width=0.9\linewidth, text width=0.9\linewidth, text width=0.9\linewidth, text width=0.9\linewidth, text width=0.9\linewidth, text width=0.9\linewidth, text width=0.9\linewidth, text width=0.9\linewidth, text width=0.9\linewidth, text width=0.9\linewidth, text width=0.9\linewidth, text width=0.9\linewidth, text width=0.9\linewidth, text width=0.9\linewidth, text width=0.9\linewidth, text width=0.9\linewidth, text width=0.9\linewidth, text width=0.9\linewidth, text width=0.9\linewidth, text width=0.9\linewidth, text width=0.9\linewidth, text width=0.9\linewidth, text width=0.9\linewidth, text width=0.9\linewidth, text width=0.9\linewidth, text width=0.9\linewidth, text width=0.9\linewidth, text width=0.9\linewidth, text width=0.9\linewidth, text width=0.9\linewidth, text width=0.9\linewidth, text width=0.9\linewidth, text width=0.9\linewidth, text width=0.9\linewidth, text width=0.9\linewidth, text width=0.9\linewidth, text width=0.9\linewidth, text width=0.9\linewidth, text width=0.9\linewidth, text width=0.9\linewidth, text width=0.9\linewidth, text width=0.9\linewidth, text width=0.9\linewidth, text width=0.9\linewidth, text width=0.9\linewidth, text width=0.9\linewidth, text width=0.9\linewidth, text width=0.9\linewidth, text width=0.9\linewidth, text width=0.9\linewidth, text width=0.9\linewidth, text width=0.9\linewidth, text width=0.9\linewidth, text width=0.9\linewidth, text width=0.9\linewidth, text width=0.9\linewidth, text width=0.9\linewidth, text width=0.9\linewidth, text width=0.9\linewidth, text width=0.9\linewidth, text width=0.9\linewidth, text width=0.9\linewidth, text width=0.9\linewidth, text width=0.9\linewidth, text width=0.9\linewidth, text width=0.9\linewidth, text width=0.9\linewidth, text width=0.9\linewidth, text width=0.9\linewidth, text width=0.9\linewidth, text width=0.9\linewidth, text width=0.9\linewidth, text width=0.9\linewidth, text width=0.9\linewidth, text width=0.9\linewidth, text width=0.9\linewidth, text width=0.9\linewidth, text width=0.9\linewidth, text width=0.9\linewidth, text width=0.9\linewidth, text width=0.9\linewidth, text width=0.9\linewidth, text width=0.9\linewidth, text width=0.9\linewidth, text width=0.9\linewidth, text width=0.9\linewidth, text width=0.9\linewidth, text width=0.9\linewidth, text width=0.9\linewidth, text width=0.9\linewidth, text width=0.9\linewidth, text width=0.9\linewidth, text width=0.9\linewidth, text width=0.9\linewidth, text width=0.9\linewidth, text width=0.9\linewidth, text width=0.9\linewidth, text width=0.9\linewidth, text width=0.9\linewidth, text width=0.9\linewidth, text width=0.9\linewidth, text width=0.9\linewidth, text width=0.9\linewidth, text width=0.9\linewidth, text width=0.9\linewidth, text width=0.9\linewidth, text width=0.9\linewidth, text width=0.9\linewidth, text width=0.9\linewidth, text width=0.9\linewidth, text width=0.9\linewidth, text width=0.9\linewidth, text width=0.9\linewidth, text width=0.9\linewidth, text width=0.9\linewidth, text width=0.9\linewidth, text width=0.9\linewidth, text width=0.9\linewidth, text width=0.9\linewidth, text width=0.9\linewidth, text width=0.9\linewidth, text width=0.9\linewidth, text width=0.9\linewidth, text width=0.9\linewidth, text width=0.9\linewidth, text width=0.9\linewidth, text width=0.9\linewidth, text width=0.9\linewidth, text width=0.9\linewidth, text width=0.9\linewidth, text width=0.9\linewidth, text width=0.9\linewidth, text width=0.9\linewidth, text width=0.9\linewidth, text width=0.9\linewidth, text width=0.9\linewidth, text width=0.9\linewidth, text width=0.9\linewidth, text width=0.9\linewidth, text width=0.9\linewidth, text width=0.9\linewidth, text width=0.9\linewidth, text width=0.9\linewidth, text width=0.9\linewidth, text width=0.9\linewidth, text width=0.9\linewidth, text width=0.9\linewidth, text width=0.9\linewidth, text width=0.9\linewidth, text width=0.9\linewidth, text width=0.9\linewidth, text width=0.9\linewidth, text width=0.9\linewidth, text width=0.9\linewidth, text width=0.9\linewidth, text width=0.9\linewidth, text width=0.9\linewidth, text width=0.9\linewidth, text width=0.9\linewidth, text width=0.9\linewidth, text width=0.9\linewidth, text width=0.9\linewidth, text width=0.9\linewidth, text width=0.9\linewidth, text width=0.9\linewidth, text width=0.9\linewidth, text width=0.9\linewidth, text width=0.9\linewidth, text width=0.9\linewidth, text width=0.9\linewidth, text width=0.9\linewidth, text width=0.9\linewidth, text width=0.9\linewidth, text width=0.9\linewidth, text width=0.9\linewidth, text width=0.9\linewidth, text width=0.9\linewidth, text width=0.9\linewidth, text width=0.9\linewidth, text width=0.9\linewidth, text width=0.9\linewidth, text width=0.9\linewidth, text width=0.9\linewidth, text width=0.9\linewidth, text width=0.9\linewidth, text width=0.9\linewidth, text width=0.9\linewidth, text width=0.9\linewidth, text width=0.9\linewidth, text width=0.9\linewidth, text width=0.9\linewidth, text width=0.9\linewidth, text width=0.9\linewidth, text width=0.9\linewidth, text width=0.9\linewidth, text width=0.9\linewidth, text width=0.9\linewidth, text width=0.9\linewidth, text width=0.9\linewidth, text width=0.9\linewidth, text width=0.9\linewidth, text width=0.9\linewidth, text width=0.9\linewidth, text width=0.9\linewidth, text width=0.9\linewidth, text width=0.9\linewidth, text width=0.9\linewidth, text width=0.9\linewidth, text width=0.9\linewidth, text width=0.9\linewidth, text width=0.9\linewidth, text width=0.9\linewidth, text width=0.9\linewidth, text width=0.9\linewidth, text width=0.9\linewidth, text width=0.9\linewidth, text width=0.9\linewidth, text width=0.9\linewidth, text width=0.9\linewidth, text width=0.9\linewidth, text width=0.9\linewidth, text width=0.9\linewidth, text width=0.9\linewidth, text width=0.9\linewidth, text width=0.9\linewidth, text width=0.9\linewidth, text width=0.9\linewidth, text width=0.9\linewidth, text width=0.9\linewidth, text width=0.9\linewidth, text width=0.9\linewidth, text width=0.9\linewidth, text width=0.9\linewidth, text width=0.9\linewidth, text width=0.9\linewidth, text width=0.9\linewidth, text width=0.9\linewidth, text width=0.9\linewidth, text width=0.9\linewidth, text width=0.9\linewidth, text width=0.9\linewidth, text width=0.9\linewidth, text width=0.9\linewidth, text width=0.9\linewidth, text width=0.9\linewidth, text width=0.9\linewidth, text width=0.9\linewidth, text width=0.9\linewidth, text width=0.9\linewidth, text width=0.9\linewidth, text width=0.9\linewidth, text width=0.9\linewidth, text width=0.9\linewidth, text width=0.9\linewidth, text width=0.9\linewidth, text width=0.9\linewidth, text width=0.9\linewidth, text width=0.9\linewidth, text width=0.9\linewidth, text width=0.9\linewidth, text width=0.9\linewidth, text width=0.9\linewidth, text width=0.9\linewidth, text width=0.9\linewidth, text width=0.9\linewidth, text width=0.9\linewidth, text width=0.9\linewidth, text width=0.9\linewidth, text width=0.9\linewidth, text width=0.9\linewidth, text width=0.9\linewidth, text width=0.9\linewidth, text width=0.9\linewidth, text width=0.9\linewidth, text width=0.9\linewidth, text width=0.9\linewidth, text width=0.9\linewidth, text width=0.9\linewidth, text width=0.9\linewidth, text width=0.9\linewidth, text width=0.9\linewidth, text width=0.9\linewidth, text width=0.9\linewidth, text width=0.9\linewidth, text width=0.9\linewidth, text width=0.9\linewidth, text width=0.9\linewidth, text width=0.9\linewidth, text width=0.9\linewidth, text width=0.9\linewidth, text width=0.9\linewidth, text width=0.9\linewidth, text width=0.9\linewidth, text width=0.9\linewidth, text width=0.9\linewidth, text width=0.9\linewidth, text width=0.9\linewidth, text width=0.9\linewidth, text width=0.9\linewidth, text width=0.9\linewidth, text width=0.9\linewidth, text width=0.9\linewidth, text width=0.9\linewidth, text width=0.9\linewidth, text width=0.9\linewidth, text width=0.9\linewidth, text width=0.9\linewidth, text width=0.9\linewidth, text width=0.9\linewidth, text width=0.9\linewidth, text width=0.9\linewidth, text width=0.9\linewidth, text width=0.9\linewidth, text width=0.9\linewidth, text width=0.9\linewidth, text width=0.9\linewidth, text width=0.9\linewidth, text width=0.9\linewidth, text width=0.9\linewidth, text width=0.9\linewidth, text width=0.9\linewidth, text width=0.9\linewidth, text width=0.9\linewidth, text width=0.9\linewidth, text width=0.9\linewidth, text width=0.9\linewidth, text width=0.9\linewidth, text width=0.9\linewidth, text width=0.9\linewidth, text width=0.9\linewidth, text width=0.9\linewidth, text width=0.9\linewidth, text width=0.9\linewidth, text width=0.9\linewidth, text width=0.9\linewidth, text width=0.9\linewidth, text width=0.9\linewidth, text width=0.9\linewidth, text width=0.9\linewidth, text width=0.9\linewidth, text width=0.9\linewidth, text width=0.9\linewidth, text width=0.9\linewidth, text width=0.9\linewidth, text width=0.9\linewidth, text width=0.9\linewidth, text width=0.9\linewidth, text width=0.9\linewidth, text width=0.9\linewidth, text width=0.9\linewidth, text width=0.9\linewidth, text width=0.9\linewidth, text width=0.9\linewidth, text width=0.9\linewidth, text width=0.9\linewidth, text width=0.9\linewidth, text width=0.9\linewidth, text width=0.9\linewidth, text width=0.9\linewidth, text width=0.9\linewidth, text width=0.9\linewidth, text width=0.9\linewidth, text width=0.9\linewidth, text width=0.9\linewidth, text width=0.9\linewidth, text width=0.9\linewidth, text width=0.9\linewidth, text width=0.9\linewidth, text width=0.9\linewidth, text width=0.9\linewidth, text width=0.9\linewidth, text width=0.9\linewidth, text width=0.9\linewidth, text width=0.9\linewidth, text width=0.9\linewidth, text width=0.9\linewidth, text width=0.9\linewidth, text width=0.9\linewidth, text width=0.9\linewidth, text width=0.9\linewidth, text width=0.9\linewidth, text width=0.9\linewidth, text width=0.9\linewidth, text width=0.9\linewidth, text width=0.9\linewidth, text width=0.9\linewidth, text width=0.9\linewidth, text width=0.9\linewidth, text width=0.9\linewidth, text width=0.9\linewidth, text width=0.9\linewidth, text width=0.9\linewidth, text width=0.9\linewidth, text width=0.9\linewidth, text width=0.9\linewidth, text width=0.9\linewidth, text width=0.9\linewidth, text width=0.9\linewidth, text width=0.9\linewidth, text width=0.9\linewidth, text width=0.9\linewidth, text width=0.9\linewidth, text width=0.9\linewidth, text width=0.9\linewidth, text width=0.9\linewidth, text width=0.9\linewidth, text width=0.9\linewidth, text width=0.9\linewidth, text width=0.9\linewidth, text width=0.9\linewidth, text width=0.9\linewidth, text width=0.9\linewidth, text width=0.9\linewidth, text width=0.9\linewidth, text width=0.9\linewidth, text width=0.9\linewidth, text width=0.9\linewidth, text width=0.9\linewidth, text width=0.9\linewidth, text width=0.9\linewidth, text width=0.9\linewidth, text width=0.9\linewidth, text width=0.9\linewidth, text width=0.9\linewidth, text width=0.9\linewidth, text width=0.9\linewidth, text width=0.9\linewidth, text width=0.9\linewidth, text width=0.9\linewidth, text width=0.9\linewidth, text width=0.9\linewidth, text width=0.9\linewidth, text width=0.9\linewidth, text width=0.9\linewidth, text width=0.9\linewidth, text width=0.9\linewidth, text width=0.9\linewidth, text width=0.9\linewidth, text width=0.9\linewidth, text width=0.9\linewidth, text width=0.9\linewidth, text width=0.9\linewidth, text width=0.9\linewidth, text width=0.9\linewidth, text width=0.9\linewidth, text width=0.9\linewidth, text width=0.9\linewidth, text width=0.9\linewidth, text width=0.9\linewidth, text width=0.9\linewidth, text width=0.9\linewidth, text width=0.9\linewidth, text width=0.9\linewidth, text width=0.9\linewidth, text width=0.9\linewidth, text width=0.9\linewidth, text width=0.9\linewidth, text width=0.9\linewidth, text width=0.9\linewidth, text width=0.9\linewidth, text width=0.9\linewidth, text width=0.9\linewidth, text width=0.9\linewidth, text width=0.9\linewidth, text width=0.9\linewidth, text width=0.9\linewidth, text width=0.9\linewidth, text width=0.9\linewidth, text width=0.9\linewidth, text width=0.9\linewidth, text width=0.9\linewidth, text width=0.9\linewidth, text width=0.9\linewidth, text width=0.9\linewidth, text width=0.9\linewidth, text width=0.9\linewidth, text width=0.9\linewidth, text width=0.9\linewidth, text width=0.9\linewidth, text width=0.9\linewidth, text width=0.9\linewidth, text width=0.9\linewidth, text width=0.9\linewidth, text width=0.9\linewidth, text width=0.9\linewidth, text width=0.9\linewidth, text width=0.9\linewidth, text width=0.9\linewidth, text width=0.9\linewidth, text width=0.9\linewidth, text width=0.9\linewidth, text width=0.9\linewidth, text width=0.9\linewidth, text width=0.9\linewidth, text width=0.9\linewidth, text width=0.9\linewidth, text width=0.9\linewidth, text width=0.9\linewidth, text width=0.9\linewidth, text width=0.9\linewidth, text width=0.9\linewidth, text width=0.9\linewidth, text width=0.9\linewidth, text width=0.9\linewidth, text width=0.9\linewidth, text width=0.9\linewidth, text width=0.9\linewidth, text width=0.9\linewidth, text width=0.9\linewidth, text width=0.9\linewidth, text width=0.9\linewidth, text width=0.9\linewidth, text width=0.9\linewidth, text width=0.9\linewidth, text width=0.9\linewidth, text width=0.9\linewidth, text width=0.9\linewidth, text width=0.9\linewidth, text width=0.9\linewidth, text width=0.9\linewidth, text width=0.9\linewidth, text width=0.9\linewidth, text width=0.9\linewidth, text width=0.9\linewidth, text width=0.9\linewidth, text width=0.9\linewidth, text width=0.9\linewidth, text width=0.9\linewidth, text width=0.9\linewidth, text width=0.9\linewidth, text width=0.9\linewidth, text width=0.9\linewidth, text width=0.9\linewidth, text width=0.9\linewidth, text width=0.9\linewidth, text width=0.9\linewidth, text width=0.9\linewidth, text width=0.9\linewidth, text width=0.9\linewidth, text width=0.9\linewidth, text width=0.9\linewidth, text width=0.9\linewidth, text width=0.9\linewidth, text width=0.9\linewidth, text width=0.9\linewidth, text width=0.9\linewidth, text width=0.9\linewidth, text width=0.9\linewidth, text width=0.9\linewidth, text width=0.9\linewidth, text width=0.9\linewidth, text width=0.9\linewidth, text width=0.9\linewidth, text width=0.9\linewidth, text width=0.9\linewidth, text width=0.9\linewidth, text width=0.9\linewidth, text width=0.9\linewidth, text width=0.9\linewidth, text width=0.9\linewidth, text width=0.9\linewidth, text width=0.9\linewidth, text width=0.9\linewidth, text width=0.9\linewidth, text width=0.9\linewidth, text width=0.9\linewidth, text width=0.9\linewidth, text width=0.9\linewidth, text width=0.9\linewidth, text width=0.9\linewidth, text width=0.9\linewidth, text width=0.9\linewidth, text width=0.9\linewidth, text width=0.9\linewidth, text width=0.9\linewidth, text width=0.9\linewidth, text width=0.9\linewidth, text width=0.9\linewidth, text width=0.9\linewidth, text width=0.9\linewidth, text width=0.9\linewidth, text width=0.9\linewidth, text width=0.9\linewidth, text width=0.9\linewidth, text width=0.9\linewidth, text width=0.9\linewidth, text width=0.9\linewidth, text width=0.9\linewidth, text width=0.9\linewidth, text width=0.9\linewidth, text width=0.9\linewidth, text width=0.9\linewidth, text width=0.9\linewidth, text width=0.9\linewidth, text width=0.9\linewidth, text width=0.9\linewidth, text width=0.9\linewidth, text width=0.9\linewidth, text width=0.9\linewidth, text width=0.9\linewidth, text width=0.9\linewidth, text width=0.9\linewidth, text width=0.9\linewidth, text width=0.9\linewidth, text width=0.9\linewidth, text width=0.9\linewidth, text width=0.9\linewidth, text width=0.9\linewidth, text width=0.9\linewidth, text width=0.9\linewidth, text width=0.9\linewidth, text width=0.9\linewidth, text width=0.9\linewidth, text width=0.9\linewidth, text width=0.9\linewidth, text width=0.9\linewidth, text width=0.9\linewidth, text width=0.9\linewidth, text width=0.9\linewidth, text width=0.9\linewidth, text width=0.9\linewidth, text width=0.9\linewidth, text width=0.9\linewidth, text width=0.9\linewidth, text width=0.9\linewidth, text width=0.9\linewidth, text width=0.9\linewidth, text width=0.9\linewidth, text width=0.9\linewidth, text width=0.9\linewidth, text width=0.9\linewidth, text width=0.9\linewidth, text width=0.9\linewidth, text width=0.9\linewidth, text width=0.9\linewidth, text width=0.9\linewidth, text width=0.9\linewidth, text width=0.9\linewidth, text width=0.9\linewidth, text width=0.9\linewidth, text width=0.9\linewidth, text width=0.9\linewidth, text width=0.9\linewidth, text width=0.9\linewidth, text width=0.9\linewidth, text width=0.9\linewidth, text width=0.9\linewidth, text width=0.9\linewidth, text width=0.9\linewidth, text width=0.9\linewidth, text width=0.9\linewidth, text width=0.9\linewidth, text width=0.9\linewidth, text width=0.9\linewidth, text width=0.9\linewidth, text width=0.9\linewidth, text width=0.9\linewidth, text width=0.9\linewidth, text width=0.9\linewidth, text width=0.9\linewidth, text width=0.9\linewidth, text width=0.9\linewidth, text width=0.9\linewidth, text width=0.9\linewidth, text width=0.9\linewidth, text width=0.9\linewidth, text width=0.9\linewidth, text width=0.9\linewidth, text width=0.9\linewidth, text width=0.9\linewidth, text width=0.9\linewidth, text width=0.9\linewidth, text width=0.9\linewidth, text width=0.9\linewidth, text width=0.9\linewidth, text width=0.9\linewidth, text width=0.9\linewidth, text width=0.9\linewidth, text width=0.9\linewidth, text width=0.9\linewidth, text width=0.9\linewidth, text width=0.9\linewidth, text width=0.9\linewidth, text width=0.9\linewidth, text width=0.9\linewidth, text width=0.9\linewidth, text width=0.9\linewidth, text width=0.9\linewidth, text width=0.9\linewidth, text width=0.9\linewidth, text width=0.9\linewidth, text width=0.9\linewidth, text width=0.9\linewidth, text width=0.9\linewidth, text width=0.9\linewidth, text width=0.9\linewidth, text width=0.9\linewidth, text width=0.9\linewidth, text width=0.9\linewidth, text width=0.9\linewidth, text width=0.9\linewidth, text width=0.9\linewidth, text width=0.9\linewidth, text width=0.9\linewidth, text width=0.9\linewidth, text width=0.9\linewidth, text width=0.9\linewidth, text width=0.9\linewidth, text width=0.9\linewidth, text width=0.9\linewidth, text width=0.9\linewidth, text width=0.9\linewidth, text width=0.9\linewidth, text width=0.9\linewidth, text width=0.9\linewidth, text width=0.9\linewidth, text width=0.9\linewidth, text width=0.9\linewidth, text width=0.9\linewidth, text width=0.9\linewidth, text width=0.9\linewidth, text width=0.9\linewidth, text width=0.9\linewidth, text width=0.9\linewidth, text width=0.9\linewidth, text width=0.9\linewidth, text width=0.9\linewidth, text width=0.9\linewidth, text width=0.9\linewidth, text width=0.9\linewidth, text width=0.9\linewidth, text width=0.9\linewidth, text width=0.9\linewidth, text width=0.9\linewidth, text width=0.9\linewidth, text width=0.9\linewidth, text width=0.9\linewidth, text width=0.9\linewidth, text width=0.9\linewidth, text width=0.9\linewidth, text width=0.9\linewidth, text width=0.9\linewidth, text width=0.9\linewidth, text width=0.9\linewidth, text width=0.9\linewidth, text width=0.9\linewidth, text width=0.9\linewidth, text width=0.9\linewidth, text width=0.9\linewidth, text width=0.9\linewidth, text width=0.9\linewidth, text width=0.9\linewidth, text width=0.9\linewidth, text width=0.9\linewidth, text width=0.9\linewidth, text width=0.9\linewidth, text width=0.9\linewidth, text width=0.9\linewidth, text width=0.9\linewidth, text width=0.9\linewidth, text width=0.9\linewidth, text width=0.9\linewidth, text width=0.9\linewidth, text width=0.9\linewidth, text width=0.9\linewidth, text width=0.9\linewidth, text width=0.9\linewidth, text width=0.9\linewidth, text width=0.9\linewidth, text width=0.9\linewidth, text width=0.9\linewidth, text width=0.9\linewidth, text width=0.9\linewidth, text width=0.9\linewidth, text width=0.9\linewidth, text width=0.9\linewidth, text width=0.9\linewidth, text width=0.9\linewidth, text width=0.9\linewidth, text width=0.9\linewidth, text width=0.9\linewidth, text width=0.9\linewidth, text width=0.9\linewidth, text width=0.9\linewidth, text width=0.9\linewidth, text width=0.9\linewidth, text width=0.9\linewidth, text width=0.9\linewidth, text width=0.9\linewidth, text width=0.9\linewidth, text width=0.9\linewidth, text width=0.9\linewidth, text width=0.9\linewidth, text width=0.9\linewidth, text width=0.9\linewidth, text width=0.9\linewidth, text width=0.9\linewidth, text width=0.9\linewidth, text width=0.9\linewidth, text width=0.9\linewidth, text width=0.9\linewidth, text width=0.9\linewidth, text width=0.9\linewidth, text width=0.9\linewidth, text width=0.9\linewidth, text width=0.9\linewidth, text width=0.9\linewidth, text width=0.9\linewidth, text width=0.9\linewidth, text width=0.9\linewidth, text width=0.9\linewidth, text width=0.9\linewidth, text width=0.9\linewidth, text width=0.9\linewidth, text width=0.9\linewidth, text width=0.9\linewidth, text width=0.9\linewidth, text width=0.9\linewidth, text width=0.9\linewidth, text width=0.9\linewidth, text width=0.9\linewidth, text width=0.9\linewidth, text width=0.9\linewidth, text width=0.9\linewidth, text width=0.9\linewidth, text width=0.9\linewidth, text width=0.9\linewidth, text width=0.9\linewidth, text width=0.9\linewidth, text width=0.9\linewidth, text width=0.9\linewidth, text width=0.9\linewidth, text width=0.9\linewidth, text width=0.9\linewidth, text width=0.9\linewidth, text width=0.9\linewidth, text width=0.9\linewidth, text width=0.9\linewidth, text width=0.9\linewidth, text width=0.9\linewidth, text width=0.9\linewidth, text width=0.9\linewidth, text width=0.9\linewidth, text width=0.9\linewidth, text width=0.9\linewidth, text width=0.9\linewidth, text width=0.9\linewidth, text width=0.9\linewidth, text width=0.9\linewidth, text width=0.9\linewidth, text width=0.9\linewidth, text width=0.9\linewidth, text width=0.9\linewidth, text width=0.9\linewidth, text width=0.9\linewidth, text width=0.9\linewidth, text width=0.9\linewidth, text width=0.9\linewidth, text width=0.9\linewidth, text width=0.9\linewidth, text width=0.9\linewidth, text width=0.9\linewidth, text width=0.9\linewidth, text width=0.9\linewidth, text width=0.9\linewidth, text width=0.9\linewidth, text width=0.9\linewidth, text width=0.9\linewidth, text width=0.9\linewidth, text width=0.9\linewidth, text width=0.9\linewidth, text width=0.9\linewidth, text width=0.9\linewidth, text width=0.9\linewidth, text width=0.9\linewidth, text width=0.9\linewidth, text width=0.9\linewidth, text width=0.9\linewidth, text width=0.9\linewidth, text width=0.9\linewidth, text width=0.9\linewidth, text width=0.9\linewidth, text width=0.9\linewidth, text width=0.9\linewidth, text width=0.9\linewidth, text width=0.9\linewidth, text width=0.9\linewidth, text width=0.9\linewidth, text width=0.9\linewidth, text width=0.9\linewidth, text width=0.9\linewidth, text width=0.9\linewidth, text width=0.9\linewidth, text width=0.9\linewidth, text width=0.9\linewidth, text width=0.9\linewidth, text width=0.9\linewidth, text width=0.9\linewidth, text width=0.9\linewidth, text width=0.9\linewidth, text width=0.9\linewidth, text width=0.9\linewidth, text width=0.9\linewidth, text width=0.9\linewidth, text width=0.9\linewidth, text width=0.9\linewidth, text width=0.9\linewidth, text width=0.9\linewidth, text width=0.9\linewidth, text width=0.9\linewidth, text width=0.9\linewidth, text width=0.9\linewidth, text width=0.9\linewidth, text width=0.9\linewidth, text width=0.9\linewidth, text width=0.9\linewidth, text width=0.9\linewidth, text width=0.9\linewidth, text width=0.9\linewidth, text width=0.9\linewidth, text width=0.9\linewidth, text width=0.9\linewidth, text width=0.9\linewidth, text width=0.9\linewidth, text width=0.9\linewidth, text width=0.9\linewidth, text width=0.9\linewidth, text width=0.9\linewidth, text width=0.9\linewidth, text width=0.9\linewidth, text width=0.9\linewidth, text width=0.9\linewidth, text width=0.9\linewidth, text width=0.9\linewidth, text width=0.9\linewidth, text width=0.9\linewidth, text width=0.9\linewidth, text width=0.9\linewidth, text width=0.9\linewidth, text width=0.9\linewidth, text width=0.9\linewidth, text width=0.9\linewidth, text width=0.9\linewidth, text width=0.9\linewidth, text width=0.9\linewidth, text width=0.9\linewidth, text width=0.9\linewidth, text width=0.9\linewidth, text width=0.9\linewidth, text width=0.9\linewidth, text width=0.9\linewidth, text width=0.9\linewidth, text width=0.9\linewidth, text width=0.9\linewidth, text width=0.9\linewidth, text width=0.9\linewidth, text width=0.9\linewidth, text width=0.9\linewidth, text width=0.9\linewidth, text width=0.9\linewidth, text width=0.9\linewidth, text width=0.9\linewidth, text width=0.9\linewidth, text width=0.9\linewidth, text width=0.9\linewidth, text width=0.9\linewidth, text width=0.9\linewidth, text width=0.9\linewidth, text width=0.9\linewidth, text width=0.9\linewidth, text width=0.9\linewidth, text width=0.9\linewidth, text width=0.9\linewidth, text width=0.9\linewidth, text width=0.9\linewidth, text width=0.9\linewidth, text width=0.9\linewidth, text width=0.9\linewidth, text width=0.9\linewidth, text width=0.9\linewidth, text width=0.9\linewidth, text width=0.9\linewidth, text width=0.9\linewidth, text width=0.9\linewidth, text width=0.9\linewidth, text width=0.9\linewidth, text width=0.9\linewidth, text width=0.9\linewidth, text width=0.9\linewidth, text width=0.9\linewidth, text width=0.9\linewidth, text width=0.9\linewidth, text width=0.9\linewidth, text width=0.9\linewidth, text width=0.9\linewidth, text width=0.9\linewidth, text width=0.9\linewidth, text width=0.9\linewidth, text width=0.9\linewidth, text width=0.9\linewidth, text width=0.9\linewidth, text width=0.9\linewidth, text width=0.9\linewidth, text width=0.9\linewidth, text width=0.9\linewidth, text width=0.9\linewidth, text width=0.9\linewidth, text width=0.9\linewidth, text width=0.9\linewidth, text width=0.9\linewidth, text width=0.9\linewidth, text width=0.9\linewidth, text width=0.9\linewidth, text width=0.9\linewidth, text width=0.9\linewidth, text width=0.9\linewidth, text width=0.9\linewidth, text width=0.9\linewidth, text width=0.9\linewidth, text width=0.9\linewidth, text width=0.9\linewidth, text width=0.9\linewidth, text width=0.9\linewidth, text width=0.9\linewidth, text width=0.9\linewidth, text width=0.9\linewidth, text width=0.9\linewidth, text width=0.9\linewidth, text width=0.9\linewidth, text width=0.9\linewidth, text width=0.9\linewidth, text width=0.9\linewidth, text width=0.9\linewidth, text width=0.9\linewidth, text width=0.9\linewidth, text width=0.9\linewidth, text width=0.9\linewidth, text width=0.9\linewidth, text width=0.9\linewidth, text width=0.9\linewidth, text width=0.9\linewidth, text width=0.9\linewidth, text width=0.9\linewidth, text width=0.9\linewidth, text width=0.9\linewidth, text width=0.9\linewidth, text width=0.9\linewidth, text width=0.9\linewidth, text width=0.9\linewidth, text width=0.9\linewidth, text width=0.9\linewidth, text width=0.9\linewidth, text width=0.9\linewidth, text width=0.9\linewidth, text width=0.9\linewidth, text width=0.9\linewidth, text width=0.9\linewidth, text width=0.9\linewidth, text width=0.9\linewidth, text width=0.9\linewidth, text width=0.9\linewidth, text width=0.9\linewidth, text width=0.9\linewidth, text width=0.9\linewidth, text width=0.9\linewidth, text width=0.9\linewidth, text width=0.9\linewidth, text width=0.9\linewidth, text width=0.9\linewidth, text width=0.9\linewidth, text width=0.9\linewidth, text width=0.9\linewidth, text width=0.9\linewidth, text width=0.9\linewidth, text width=0.9\linewidth, text width=0.9\linewidth, text width=0.9\linewidth, text width=0.9\linewidth, text width=0.9\linewidth, text width=0.9\linewidth, text width=0.9\linewidth, text width=0.9\linewidth, text width=0.9\linewidth, text width=0.9\linewidth, text width=0.9\linewidth, text width=0.9\linewidth, text width=0.9\linewidth, text width=0.9\linewidth, text width=0.9\linewidth, text width=0.9\linewidth, text width=0.9\linewidth, text width=0.9\linewidth, text width=0.9\linewidth, text width=0.9\linewidth, text width=0.9\linewidth, text width=0.9\linewidth, text width=0.9\linewidth, text width=0.9\linewidth, text width=0.9\linewidth, text width=0.9\linewidth, text width=0.9\linewidth, text width=0.9\linewidth, text width=0.9\linewidth, text width=0.9\linewidth, text width=0.9\linewidth, text width=0.9\linewidth, text width=0.9\linewidth, text width=0.9\linewidth, text width=0.9\linewidth, text width=0.9\linewidth, text width=0.9\linewidth, text width=0.9\linewidth, text width=0.9\linewidth, text width=0.9\linewidth, text width=0.9\linewidth, text width=0.9\linewidth, text width=0.9\linewidth, text width=0.9\linewidth, text width=0.9\linewidth, text width=0.9\linewidth, text width=0.9\linewidth, text width=0.9\linewidth, text width=0.9\linewidth, text width=0.9\linewidth, text width=0.9\linewidth, text width=0.9\linewidth, text width=0.9\linewidth, text width=0.9\linewidth, text width=0.9\linewidth, text width=0.9\linewidth, text width=0.9\linewidth, text width=0.9\linewidth, text width=0.9\linewidth, text width=0.9\linewidth, text width=0.9\linewidth, text width=0.9\linewidth, text width=0.9\linewidth, text width=0.9\linewidth, text width=0.9\linewidth, text width=0.9\linewidth, text width=0.9\linewidth, text width=0.9\linewidth, text width=0.9\linewidth, text width=0.9\linewidth, text width=0.9\linewidth, text width=0.9\linewidth, text width=0.9\linewidth, text width=0.9\linewidth, text width=0.9\linewidth, text width=0.9\linewidth, text width=0.9\linewidth, text width=0.9\linewidth, text width=0.9\linewidth, text width=0.9\linewidth, text width=0.9\linewidth, text width=0.9\linewidth, text width=0.9\linewidth, text width=0.9\linewidth, text width=0.9\linewidth, text width=0.9\linewidth, text width=0.9\linewidth, text width=0.9\linewidth, text width=0.9\linewidth, text width=0.9\linewidth, text width=0.9\linewidth, text width=0.9\linewidth, text width=0.9\linewidth, text width=0.9\linewidth, text width=0.9\linewidth, text width=0.9\linewidth, text width=0.9\linewidth, text width=0.9\linewidth, text width=0.9\linewidth, text width=0.9\linewidth, text width=0.9\linewidth, text width=0.9\linewidth, text width=0.9\linewidth, text width=0.9\linewidth, text width=0.9\linewidth, text width=0.9\linewidth, text width=0.9\linewidth, text width=0.9\linewidth, text width=0.9\linewidth, text width=0.9\linewidth, text width=0.9\linewidth, text width=0.9\linewidth, text width=0.9\linewidth, text width=0.9\linewidth, text width=0.9\linewidth, text width=0.9\linewidth, text width=0.9\linewidth, text width=0.9\linewidth, text width=0.9\linewidth, text width=0.9\linewidth, text width=0.9\linewidth, text width=0.9\linewidth, text width=0.9\linewidth, text width=0.9\linewidth, text width=0.9\linewidth, text width=0.9\linewidth, text width=0.9\linewidth, text width=0.9\linewidth, text width=0.9\linewidth, text width=0.9\linewidth, text width=0.9\linewidth, text width=0.9\linewidth, text width=0.9\linewidth, text width=0.9\linewidth, text width=0.9\linewidth, text width=0.9\linewidth, text width=0.9\linewidth, text width=0.9\linewidth, text width=0.9\linewidth, text width=0.9\linewidth, text width=0.9\linewidth, text width=0.9\linewidth, text width=0.9\linewidth, text width=0.9\linewidth, text width=0.9\linewidth, text width=0.9\linewidth, text width=0.9\linewidth, text width=0.9\linewidth, text width=0.9\linewidth, text width=0.9\linewidth, text width=0.9\linewidth, text width=0.9\linewidth, text width=0.9\linewidth, text width=0.9\linewidth, text width=0.9\linewidth, text width=0.9\linewidth, text width=0.9\linewidth, text width=0.9\linewidth, text width=0.9\linewidth, text width=0.9\linewidth, text width=0.9\linewidth, text width=0.9\linewidth, text width=0.9\linewidth, text width=0.9\linewidth, text width=0.9\linewidth, text width=0.9\linewidth, text width=0.9\linewidth, text width=0.9\linewidth, text width=0.9\linewidth, text width=0.9\linewidth, text width=0.9\linewidth, text width=0.9\linewidth, text width=0.9\linewidth, text width=0.9\linewidth, text width=0.9\linewidth, text width=0.9\linewidth, text width=0.9\linewidth, text width=0.9\linewidth, text width=0.9\linewidth, text width=0.9\linewidth, text width=0.9\linewidth, text width=0.9\linewidth, text width=0.9\linewidth, text width=0.9\linewidth, text width=0.9\linewidth, text width=0.9\linewidth, text width=0.9\linewidth, text width=0.9\linewidth, text width=0.9\linewidth, text width=0.9\linewidth, text width=0.9\linewidth, text width=0.9\linewidth, text width=0.9\linewidth, text width=0.9\linewidth, text width=0.9\linewidth, text width=0.9\linewidth, text width=0.9\linewidth, text width=0.9\linewidth, text width=0.9\linewidth, text width=0.9\linewidth, text width=0.9\linewidth, text width=0.9\linewidth, text width=0.9\linewidth, text width=0.9\linewidth, text width=0.9\linewidth, text width=0.9\linewidth, text width=0.9\linewidth, text width=0.9\linewidth, text width=0.9\linewidth, text width=0.9\linewidth, text width=0.9\linewidth, text width=0.9\linewidth, text width=0.9\linewidth, text width=0.9\linewidth, text width=0.9\linewidth, text width=0.9\linewidth, text width=0.9\linewidth, text width=0.9\linewidth, text width=0.9\linewidth, text width=0.9\linewidth, text width=0.9\linewidth, text width=0.9\linewidth, text width=0.9\linewidth, text width=0.9\linewidth, text width=0.9\linewidth, text width=0.9\linewidth, text width=0.9\linewidth, text width=0.9\linewidth, text width=0.9\linewidth, text width=0.9\linewidth, text width=0.9\linewidth, text width=0.9\linewidth, text width=0.9\linewidth, text width=0.9\linewidth, text width=0.9\linewidth, text width=0.9\linewidth, text width=0.9\linewidth, text width=0.9\linewidth, text width=0.9\linewidth, text width=0.9\linewidth, text width=0.9\linewidth, text width=0.9\linewidth, text width=0.9\linewidth, text width=0.9\linewidth, text width=0.9\linewidth, text width=0.9\linewidth, text width=0.9\linewidth, text width=0.9\linewidth, text width=0.9\linewidth, text width=0.9\linewidth, text width=0.9\linewidth, text width=0.9\linewidth, text width=0.9\linewidth, text width=0.9\linewidth, text width=0.9\linewidth, text width=0.9\linewidth, text width=0.9\linewidth, text width=0.9\linewidth, text width=0.9\linewidth, text width=0.9\linewidth, text width=0.9\linewidth, text width=0.9\linewidth, text width=0.9\linewidth, text width=0.9\linewidth, text width=0.9\linewidth, text width=0.9\linewidth, text width=0.9\linewidth, text width=0.9\linewidth, text width=0.9\linewidth, text width=0.9\linewidth, text width=0.9\linewidth, text width=0.9\linewidth, text width=0.9\linewidth, text width=0.9\linewidth, text width=0.9\linewidth, text width=0.9\linewidth, text width=0.9\linewidth, text width=0.9\linewidth, text width=0.9\linewidth, text width=0.9\linewidth, text width=0.9\linewidth, text width=0.9\linewidth, text width=0.9\linewidth, text width=0.9\linewidth, text width=0.9\linewidth, text width=0.9\linewidth, text width=0.9\linewidth, text width=0.9\linewidth, text width=0.9\linewidth, text width=0.9\linewidth, text width=0.9\linewidth, text width=0.9\linewidth, text width=0.9\linewidth, text width=0.9\linewidth, text width=0.9\linewidth, text width=0.9\linewidth, text width=0.9\linewidth, text width=0.9\linewidth, text width=0.9\linewidth, text width=0.9\linewidth, text width=0.9\linewidth, text width=0.9\linewidth, text width=0.9\linewidth, text width=0.9\linewidth, text width=0.9\linewidth, text width=0.9\linewidth, text width=0.9\linewidth, text width=0.9\linewidth, text width=0.9\linewidth, text width=0.9\linewidth, text width=0.9\linewidth, text width=0.9\linewidth, text width=0.9\linewidth, text width=0.9\linewidth, text width=0.9\linewidth, text width=0.9\linewidth, text width=0.9\linewidth, text width=0.9\linewidth, text width=0.9\linewidth, text width=0.9\linewidth, text width=0.9\linewidth, text width=0.9\linewidth, text width=0.9\linewidth, text width=0.9\linewidth, text width=0.9\linewidth, text width=0.9\linewidth, text width=0.9\linewidth, text width=0.9\linewidth, text width=0.9\linewidth, text width=0.9\linewidth, text width=0.9\linewidth, text width=0.9\linewidth, text width=0.9\linewidth, text width=0.9\linewidth, text width=0.9\linewidth, text width=0.9\linewidth, text width=0.9\linewidth, text width=0.9\linewidth, text width=0.9\linewidth]{Sulphur \& Nitrogen}
  [clockwise from=90, level 1 concept/.append style={sibling angle=60, level distance=3.4cm, text width=2.6cm, align=center, font=\small\bfseries}]
  child[concept color=popgold]{ node {S oxidation states $-2,+2,+4,+6$} }
  child[concept color=poporange]{ node {Contact process: \ce{H2SO4}} }
  child[concept color=popred]{ node {\ce{SO2}, \ce{SO3}: acid rain} }
  child[concept color=popblue]{ node {N oxidation states $-3$ to $+5$} }
  child[concept color=popgreen]{ node {Haber: \ce{NH3}; Ostwald: \ce{HNO3}} }
  child[concept color=poppurple]{ node {Fertilizers \& eutrophication} };
\end{tikzpicture}
\legende{Mind map of the chapter: sulphur and nitrogen in the environment.}
\end{popfigure}

\begin{bilan}
\textbf{Key definitions}
\begin{itemize}
\item \cle{Oxidation state}: the charge an atom would have if all its bonds were fully ionic; sulphur shows $-2$ (\ce{H2S}), $+2$ (\ce{SO}), $+4$ (\ce{SO2}, \ce{H2SO3}), $+6$ (\ce{SO3}, \ce{H2SO4}); nitrogen ranges from $-3$ (\ce{NH3}) to $+5$ (\ce{HNO3}).
\item \cle{Acid rain}: rain with pH below about 5.6, caused mainly by dissolved \ce{SO2} and nitrogen oxides forming \ce{H2SO4} and \ce{HNO3}.
\item \cle{Eutrophication}: excessive enrichment of water with nitrates and phosphates, causing algal blooms, oxygen depletion and death of aquatic life.
\item \cle{Catalyst}: a substance that increases reaction rate without being consumed (\ce{V2O5} in the Contact process, Fe in the Haber process, Pt/Rh in the Ostwald process).
\end{itemize}

\textbf{Key reactions and processes to know}
\begin{itemize}
\item Contact process: \ce{S(s) + O2(g) -> SO2(g)}; \ensuremath{\mathrm{2SO_{2}(g) + O_{2}(g) \rightleftharpoons [V_{2}O_{5}, 450\degree C] 2SO_{3}(g)}} (exothermic, moderate temperature, high pressure not needed); \ce{SO3 + H2SO4 -> H2S2O7} (oleum), then \ce{H2S2O7 + H2O -> 2H2SO4}. Never add \ce{SO3} directly to water (acid mist).
\item Properties of concentrated \ce{H2SO4}: dehydrating agent (chars sugar, dries gases), oxidising agent with hot metals: \ce{Cu + 2H2SO4 -> CuSO4 + SO2 + 2H2O}; dilute acid behaves as a typical dibasic acid.
\item Sulphur dioxide: \ce{SO2 + H2O <=> H2SO3}; bleaches by reduction; reducing agent (decolourises acidified \ce{KMnO4}).
\item Haber process: \ensuremath{\mathrm{N_{2}(g) + 3H_{2}(g) \rightleftharpoons [Fe, 450\degree C, 200 atm] 2NH_{3}(g)}}, $\Delta H < 0$; compromise conditions.
\item Test for \ce{NH3}: turns damp red litmus blue; white fumes with concentrated \ce{HCl}. Test for \ce{NH4+}: warm with \ce{NaOH}, ammonia evolved.
\item Ostwald process: \ce{4NH3 + 5O2 ->[Pt] 4NO + 6H2O}; \ce{2NO + O2 -> 2NO2}; \ce{3NO2 + H2O -> 2HNO3 + NO}.
\item Dilute \ce{HNO3} with Cu: \ce{3Cu + 8HNO3 -> 3Cu(NO3)2 + 2NO + 4H2O}; concentrated acid gives \ce{NO2} (brown fumes).
\item Fertilizers: \ce{NH4NO3}, \ce{(NH4)2SO4}, urea \ce{CO(NH2)2}, NPK; percentage of nitrogen in \ce{NH4NO3} $= \frac{28}{80}\times 100 = 35\%$.
\end{itemize}

\textbf{Methods}
\begin{itemize}
\item Assign oxidation states from rules (O usually $-2$, H usually $+1$), then identify oxidation (increase) and reduction (decrease).
\item For industrial processes, always state: equation, catalyst, temperature, pressure, and justify conditions with Le Chatelier's principle.
\item Environmental questions: name the pollutant, its source, the chemical equation of its formation, its effect, and one control measure (e.g. flue-gas desulphurisation with \ce{CaCO3}).
\end{itemize}

\textbf{Common mistakes}
\begin{itemize}
\item Writing \ce{SO3 + H2O -> H2SO4} as the industrial step: industrially \ce{SO3} is absorbed in concentrated \ce{H2SO4}, not water.
\item Confusing the bleaching action of \ce{SO2} (reduction) with that of \ce{Cl2} (oxidation).
\item Saying concentrated \ce{H2SO4} liberates \ce{H2} with metals: it is an oxidising acid and gives \ce{SO2}.
\item Forgetting state symbols and conditions in the Contact and Haber equations.
\item Attributing acid rain only to \ce{SO2}: nitrogen oxides (\ce{NO}, \ce{NO2}) also contribute.
\end{itemize}

\textbf{GCE tips}
\begin{itemize}
\item Learn the three industrial processes with exact conditions; they are asked almost every year.
\item In essay questions on pollution, structure your answer: source $\to$ equation $\to$ effect $\to$ remedy.
\item Show all steps in percentage-nitrogen calculations and give the answer to 3 significant figures.
\item Link fertilizer use in Cameroonian agriculture (maize, cocoa) to both increased yield and the risk of eutrophication of streams and lakes.
\end{itemize}
\end{bilan}

\findecours

\end{document}
